% L'adaptateur secteur — PCT 3ème — Cours signé OnBuch+
% Programme officiel MINESEC. Compiler avec : tectonic N15-adaptateur-secteur.tex
\newcommand{\DOCMATIERE}{PCT}
\newcommand{\DOCNIVEAU}{3ème}
\newcommand{\DOCMODULE}{Module 2 : Électricité}
\newcommand{\DOCLECON}{N15}
\newcommand{\DOCTITRE}{L'adaptateur secteur}
% OnBuch+ — préambule « cours signés » ENRICHI (compilé avec tectonic / XeLaTeX)
% Système visuel « Pop » d'OnBuch+ : fond crème (jamais blanc), bordures encre
% épaisses, ombres « dures » décalées, titres Archivo Black en capitales.
% Version MATHÉMATIQUES : graphes (pgfplots/tikz), boîtes pédagogiques
% (définition, propriété, méthode, attention, le savais-tu, exercice résolu,
% exercice, TP, bilan), page de garde et signature OnBuch+ ancrée en bas.
%
% Macros attendues dans main.tex AVANT \input{preamble} :
%   \DOCMATIERE  \DOCNIVEAU  \DOCMODULE  \DOCLECON (numéro)  \DOCTITRE
\documentclass[11pt]{article}

\usepackage{fontspec}
\usepackage[french]{babel}
\usepackage[a4paper,margin=2cm,top=3.4cm,bottom=2.8cm,headheight=1.4cm,footskip=1.3cm]{geometry}
\usepackage{amsmath,amssymb,amsfonts}
\usepackage{mathtools} % \xmapsto, \coloneqq... souvent utilisés par les modèles
\usepackage{cancel} % simplifications barrées (bilans de forces)
% \abs{z} (module, valeur absolue) et \norm{u} (norme), utilisés par les modèles
\providecommand{\abs}{}\let\abs\relax\DeclarePairedDelimiter{\abs}{\lvert}{\rvert}
\providecommand{\norm}{}\let\norm\relax\DeclarePairedDelimiter{\norm}{\lVert}{\rVert}
\usepackage[table]{xcolor} % [table] : fournit aussi \rowcolors (alternance de lignes)
\usepackage{pagecolor}
\usepackage{tikz}
\usetikzlibrary{arrows.meta,positioning,calc,shapes.geometric,shapes.misc,shapes.arrows,decorations.pathmorphing,decorations.pathreplacing,decorations.markings,patterns,shadows,mindmap,trees,fit,backgrounds,matrix,intersections,angles,quotes}
\usepackage{pgfplots}
\usepackage[version=4]{mhchem} % équations nucléaires \ce{^{235}_{92}U}
\usepackage[european,siunitx]{circuitikz} % schémas électriques
\pgfplotsset{compat=1.18}
\usepgfplotslibrary{fillbetween,groupplots} % aires entre courbes (calcul intégral)
\usepackage{enumitem}
\usepackage{fancyhdr}
% [most] charge toutes les bibliothèques tcolorbox (listings, minted, raster,
% documentation...) et gonfle inutilement la mémoire TeX sur un document long
% (~300 boîtes) : on ne charge que ce qui est réellement utilisé.
\usepackage[skins,breakable]{tcolorbox}
\usepackage{graphicx}
\usepackage{colortbl}
\usepackage{booktabs}
\usepackage{tabularx}
\usepackage{multirow}
\usepackage{diagbox} % case d'en-tête diagonale des tableaux à double entrée
% Colonnes extensibles centrée / gauche / droite (C, L, R), souvent utilisées par les modèles
\newcolumntype{C}{>{\centering\arraybackslash}X}
\newcolumntype{L}{>{\raggedright\arraybackslash}X}
\newcolumntype{R}{>{\raggedleft\arraybackslash}X}
\usepackage{array}
\usepackage{multicol}
\usepackage{siunitx}
% parse-numbers=false : accepte aussi \SI{2,5 \times 10^{-3}}{...}
\sisetup{locale=FR,output-decimal-marker={,},group-minimum-digits=4,parse-numbers=false}
\DeclareSIUnit\ohm{\ensuremath{\symup{\Omega}}} % U+2126 absent de la police
\DeclareSIUnit\division{div} % réglages d'oscilloscope (ms/div, V/div)
\DeclareSIUnit\tr{tr} % tours (vitesse de rotation, ex. \SI{1500}{\tr\per\minute})
\DeclareSIUnit\molar{M} % concentration molaire (ex. \SI{3}{\molar}, \SI{1}{\micro\molar})
\DeclareSIUnit\an{an} % année (le modèle confond parfois avec \year, un compteur TeX) — ex. \SI{5}{\kilo\meter\per\an}
\DeclareSIUnit\FCFA{FCFA} % franc CFA (ex. \SI{500}{\FCFA\per\kilo\gram})
\DeclareSIUnit\degreeBrix{\textdegree Brix} % teneur en sucre d'un sirop/jus (réfractométrie)
\DeclareSIUnit\month{mois} % le modèle utilise parfois \month, un compteur TeX, comme unité de durée
\DeclareSIUnit\microliter{\micro\liter} % le modèle écrit parfois \microliter en un seul token
\DeclareSIUnit\million{millions} % dénombrement (ex. \SI{15}{\million\per\mL} spermatozoïdes)
\usepackage{qrcode}
% \degree : commande fréquemment utilisée par les modèles pour un angle ou une
% température en dehors de \SI{}{} (non fournie par les paquets chargés).
\providecommand{\degree}{^{\circ}}
% \qed (fin de démonstration), utilisé par les modèles sans amsthm
\providecommand{\qed}{\hfill$\square$}
% \barre : double barre des valeurs interdites dans un tableau de variations (tkz-tab)
\providecommand{\barre}{\ensuremath{\|}}
% environnement proof (amsthm non chargé) utilisé par les modèles
\ifdefined\proof\else\newenvironment{proof}[1][Démonstration]{\par\noindent\textit{#1.}\ }{\qed\par}\fi
\usepackage{lastpage}
\usepackage{hyperref}
\hypersetup{hidelinks}

% ============================================================
% Palette « Pop » OnBuch+ — mêmes hex que la classe P (plus_ui.dart)
% ============================================================
\definecolor{poporange}{HTML}{F59321}
\definecolor{popdark}{HTML}{E07A0C}
\definecolor{popcream}{HTML}{FFF6E8}
\definecolor{popink}{HTML}{1C1714}
\definecolor{poppurple}{HTML}{7A5AE0}
\definecolor{popgreen}{HTML}{2E9E6A}
\definecolor{poppink}{HTML}{E0457B}
\definecolor{popblue}{HTML}{2F7DE1}
\definecolor{popgold}{HTML}{C98A12}
\definecolor{popmuted}{HTML}{8A7F76}
\definecolor{popline}{HTML}{EADFD2}
\definecolor{popgray}{HTML}{8A7F76}
\colorlet{popgrey}{popgray}
% Alias tolérants (le modèle nomme parfois les couleurs différemment)
\colorlet{popwhite}{white}
\colorlet{popblack}{popink}
\colorlet{popred}{poppink}
\colorlet{popyellow}{popgold}
% Teintes claires (fonds de tableaux/encadrés) — le modèle nomme parfois
% les couleurs avec un suffixe L (ex. popblueL) sans les avoir définies.
\colorlet{poporangeL}{poporange!20}
\colorlet{popdarkL}{popdark!20}
\colorlet{poppurpleL}{poppurple!20}
\colorlet{popgreenL}{popgreen!20}
\colorlet{poppinkL}{poppink!20}
\colorlet{popblueL}{popblue!20}
\colorlet{popgoldL}{popgold!20}
\colorlet{popredL}{popred!20}
\colorlet{popgrayL}{popgray!20}
% Teintes claires pour fonds de tableaux / zones
\colorlet{popblueL}{popblue!10!white}
\colorlet{poporangeL}{poporange!14!white}
\colorlet{popgreenL}{popgreen!12!white}
\colorlet{poppurpleL}{poppurple!12!white}
\colorlet{poppinkL}{poppink!12!white}
\colorlet{popgoldL}{popgold!14!white}

\pagecolor{popcream}

% ============================================================
% Typographie
% ============================================================
\defaultfontfeatures{Ligatures=TeX}
\setmainfont{PlusJakartaSans-Medium.ttf}[Path=../../../fonts/,BoldFont=PlusJakartaSans-Bold.ttf]
% Maths en sans-serif (Fira Math) avec chiffres et lettres droites de Plus
% Jakarta, pour que les formules se fondent dans le texte.
\usepackage[math-style=ISO,bold-style=ISO]{unicode-math}
\setmathfont{FiraMath-Regular.otf}[Path=../../../fonts/]
\setmathfont{PlusJakartaSans-Medium.ttf}[Path=../../../fonts/,range={up/num,up/{Latin,latin},\mathup}]
\usepackage{icomma} % virgule décimale sans espace en mode math
% Caractères mathématiques Unicode absents des polices de texte (Plus Jakarta,
% Archivo Black) : rendus par leur équivalent mathématique, en texte comme en
% formule — sinon ils s'affichent en carré vide (ex. « Arithmétique dans ℤ »).
\usepackage{newunicodechar}
\newunicodechar{ℝ}{\ensuremath{\mathbb{R}}}
\newunicodechar{ℤ}{\ensuremath{\mathbb{Z}}}
\newunicodechar{ℂ}{\ensuremath{\mathbb{C}}}
\newunicodechar{ℕ}{\ensuremath{\mathbb{N}}}
\newunicodechar{ℚ}{\ensuremath{\mathbb{Q}}}
\newunicodechar{𝔻}{\ensuremath{\mathbb{D}}}
\newunicodechar{α}{\ensuremath{\alpha}}
\newunicodechar{β}{\ensuremath{\beta}}
\newunicodechar{γ}{\ensuremath{\gamma}}
\newunicodechar{δ}{\ensuremath{\delta}}
\newunicodechar{ε}{\ensuremath{\varepsilon}}
\newunicodechar{θ}{\ensuremath{\theta}}
\newunicodechar{λ}{\ensuremath{\lambda}}
\newunicodechar{μ}{\ensuremath{\mu}}
\newunicodechar{σ}{\ensuremath{\sigma}}
\newunicodechar{τ}{\ensuremath{\tau}}
\newunicodechar{φ}{\ensuremath{\varphi}}
\newunicodechar{ψ}{\ensuremath{\psi}}
\newunicodechar{ω}{\ensuremath{\omega}}
\newunicodechar{Σ}{\ensuremath{\Sigma}}
% majuscules produites par \MakeUppercase dans les titres (α-aminés -> Α)
\newunicodechar{Α}{\ensuremath{\alpha}}
\newunicodechar{Β}{\ensuremath{\beta}}
\newunicodechar{Γ}{\ensuremath{\Gamma}}
\newunicodechar{Δ}{\ensuremath{\Delta}}
\newunicodechar{Ε}{\ensuremath{\varepsilon}}
\newunicodechar{Θ}{\ensuremath{\Theta}}
\newunicodechar{Λ}{\ensuremath{\Lambda}}
\newunicodechar{Μ}{\ensuremath{\mu}}
\newunicodechar{Τ}{\ensuremath{\tau}}
\newunicodechar{Φ}{\ensuremath{\Phi}}
\newunicodechar{Ψ}{\ensuremath{\Psi}}
\newunicodechar{Ω}{\ensuremath{\Omega}}
\newunicodechar{↦}{\ensuremath{\mapsto}}
\newunicodechar{∈}{\ensuremath{\in}}
\newunicodechar{∉}{\ensuremath{\notin}}
\newunicodechar{∧}{\ensuremath{\wedge}}
\newunicodechar{⇔}{\ensuremath{\Leftrightarrow}}
\newunicodechar{⟺}{\ensuremath{\Longleftrightarrow}}
\newunicodechar{⇒}{\ensuremath{\Rightarrow}}
\newunicodechar{⟹}{\ensuremath{\Longrightarrow}}
\newunicodechar{∘}{\ensuremath{\circ}}
\newunicodechar{∪}{\ensuremath{\cup}}
\newunicodechar{∩}{\ensuremath{\cap}}
\newunicodechar{≡}{\ensuremath{\equiv}}
\newunicodechar{⊂}{\ensuremath{\subset}}
\newunicodechar{⊥}{\ensuremath{\perp}}
\newunicodechar{∃}{\ensuremath{\exists}}
\newunicodechar{∀}{\ensuremath{\forall}}
\newunicodechar{∛}{\ensuremath{\sqrt[3]{\,}}}
\newunicodechar{∜}{\ensuremath{\sqrt[4]{\,}}}
\newunicodechar{⃗}{\ensuremath{{}^{\to}}}
\newunicodechar{ⁿ}{\ifmmode^{n}\else\textsuperscript{\ensuremath{n}}\fi}
\newunicodechar{ˣ}{\ifmmode^{x}\else\textsuperscript{\ensuremath{x}}\fi}
\newunicodechar{ᵇ}{\ifmmode^{b}\else\textsuperscript{\ensuremath{b}}\fi}
\newunicodechar{ʳ}{\ifmmode^{r}\else\textsuperscript{\ensuremath{r}}\fi}
\newunicodechar{ᵘ}{\ifmmode^{u}\else\textsuperscript{\ensuremath{u}}\fi}
\newunicodechar{ʸ}{\ifmmode^{y}\else\textsuperscript{\ensuremath{y}}\fi}
\newunicodechar{ᵃ}{\ifmmode^{a}\else\textsuperscript{\ensuremath{a}}\fi}
\newunicodechar{ᵉ}{\ifmmode^{e}\else\textsuperscript{\ensuremath{e}}\fi}
\newunicodechar{ᵏ}{\ifmmode^{k}\else\textsuperscript{\ensuremath{k}}\fi}
\newunicodechar{ᵖ}{\ifmmode^{p}\else\textsuperscript{\ensuremath{p}}\fi}
\newunicodechar{ᴺ}{\ifmmode^{N}\else\textsuperscript{\ensuremath{N}}\fi}
\newunicodechar{ᶜ}{\ifmmode^{c}\else\textsuperscript{\ensuremath{c}}\fi}
\newunicodechar{ʰ}{\ifmmode^{h}\else\textsuperscript{\ensuremath{h}}\fi}
\newunicodechar{ᵠ}{\ifmmode^{q}\else\textsuperscript{\ensuremath{q}}\fi}
\newunicodechar{ₙ}{\ifmmode_{n}\else\textsubscript{\ensuremath{n}}\fi}
\newunicodechar{ₐ}{\ifmmode_{a}\else\textsubscript{\ensuremath{a}}\fi}
\newunicodechar{ᵢ}{\ifmmode_{i}\else\textsubscript{\ensuremath{i}}\fi}
\newunicodechar{ᵣ}{\ifmmode_{r}\else\textsubscript{\ensuremath{r}}\fi}
\newunicodechar{ₖ}{\ifmmode_{k}\else\textsubscript{\ensuremath{k}}\fi}
\newunicodechar{ᵦ}{\ifmmode_{\beta}\else\textsubscript{\ensuremath{\beta}}\fi}
\newunicodechar{ₓ}{\ifmmode_{x}\else\textsubscript{\ensuremath{x}}\fi}
\newfontfamily\popdisplay{ArchivoBlack-Regular.ttf}[Path=../../../fonts/]
\newfontfamily\poplogo{SpaceGrotesk-Bold.ttf}[Path=../../../fonts/]
\newfontfamily\popsemi{PlusJakartaSans-SemiBold.ttf}[Path=../../../fonts/]
\newfontfamily\popxbold{PlusJakartaSans-ExtraBold.ttf}[Path=../../../fonts/]
\newfontfamily\popxboldls{PlusJakartaSans-ExtraBold.ttf}[Path=../../../fonts/,LetterSpace=3.0]

\setlist[enumerate]{leftmargin=1.7em,itemsep=5pt,topsep=4pt}
\setlist[itemize]{leftmargin=1.7em,itemsep=4pt,topsep=4pt,label={\color{poporange}\textbullet}}
\setlength{\parindent}{0pt}
\setlength{\parskip}{5pt}
\renewcommand{\arraystretch}{1.25}

% pgfplots : style Pop par défaut
\pgfplotsset{
  every axis/.append style={axis line style={popink,line width=1pt},tick style={popink},
    grid style={popline,line width=0.5pt},label style={font=\small},tick label style={font=\footnotesize}},
  popcurve/.style={poporange,line width=1.8pt,no markers,smooth},
}
% Même style disponible dans un \draw TikZ simple (hors pgfplots)
\tikzset{popcurve/.style={poporange,line width=1.8pt}}
\tikzset{popgrid/.style={popline,very thin}}
\colorlet{popcurve}{poporange} % le nom du style est parfois utilisé comme couleur

% ============================================================
% Pill / squiggle
% ============================================================
\newcommand{\poppill}[3]{%
  \tikz[baseline=-0.55ex]\node[draw=popink,line width=1.2pt,rounded corners=2.6pt,%
    fill=#2,inner xsep=6.5pt,inner ysep=3pt]{%
    \popxboldls\fontsize{7}{7}\selectfont\color{#3}\MakeUppercase{#1}};%
}
\newcommand{\popsquiggle}[1]{%
  \tikz[baseline=-0.4ex]{
    \draw[#1,line width=2.6pt,line cap=round]
      plot [smooth] coordinates {(0,0.09) (0.34,0.01) (0.68,0.21) (1.02,0.01) (1.36,0.10)};
  }%
}
% Mot-clé surligné (marqueur orange sous le texte)
\newcommand{\cle}[1]{{\popxbold\color{popdark}#1}}

% ============================================================
% En-tête / pied
% ============================================================
\pagestyle{fancy}
\fancyhf{}
\renewcommand{\headrulewidth}{0pt}
\renewcommand{\footrulewidth}{0pt}
\lhead{\raisebox{-0.15ex}{\poplogo\fontsize{13}{13}\selectfont\color{popink}OnBuch\color{poporange}+}}
\rhead{\poppill{\DOCMATIERE\ \textbullet\ \DOCNIVEAU\ \textbullet\ Leçon \DOCLECON}{white}{popink}}
\lfoot{\popxbold\fontsize{7.6}{7.6}\selectfont\color{popmuted}\MakeUppercase{Cours signé OnBuch+}\ \textbullet\ {\popsemi\DOCTITRE}}
\rfoot{\tikz[baseline=-0.6ex]\node[rounded corners=8.5pt,fill=popink,minimum height=17pt,minimum width=17pt,inner xsep=5pt,inner sep=0]{\popxbold\fontsize{8}{8}\selectfont\color{popcream}\thepage\,/\,\pageref*{LastPage}};}

% ============================================================
% Titres
% ============================================================
\newcommand{\chaptitre}[2]{%
  \begin{tcolorbox}[enhanced,colback=poporange,colframe=popink,boxrule=2.6pt,arc=11pt,
      drop shadow=popink,left=15pt,right=15pt,top=12pt,bottom=11pt,before skip=3pt,after skip=18pt]
    \poppill{#2}{popink}{popcream}\par\vspace{9pt}
    {\raggedright\hyphenpenalty=10000\exhyphenpenalty=10000\popdisplay\fontsize{22}{24}\selectfont\color{popcream}\MakeUppercase{#1}\par}
  \end{tcolorbox}
}
% Section numérotée : pastille orange avec le numéro + titre Archivo
\newcounter{popsec}
\newcommand{\coursec}[1]{%
  \stepcounter{popsec}\setcounter{popsub}{0}%
  \par\vspace{16pt}\noindent
  \begin{minipage}{\linewidth}
  \tikz[baseline=-0.7ex]\node[draw=popink,line width=1.6pt,fill=poporange,rounded corners=4pt,minimum size=22pt,inner sep=0]{\popdisplay\fontsize{12}{12}\selectfont\color{popcream}\thepopsec};%
  \hspace{8pt}{\popdisplay\fontsize{15}{17}\selectfont\color{popink}\MakeUppercase{#1}}\par
  \vspace{4pt}\noindent{\color{poporange}\rule{36pt}{3.6pt}}
  \end{minipage}\par\nopagebreak\vspace{8pt}
}
\newcounter{popsub}
\newcommand{\courssub}[1]{%
  \stepcounter{popsub}%
  \par\vspace{9pt}\noindent{\popxbold\fontsize{11.5}{13}\selectfont\color{popink}{\color{poporange}\thepopsec.\thepopsub}\hspace{6pt}#1}\par\nopagebreak\vspace{4pt}
}

% ============================================================
% Boîtes pédagogiques — même recette : fond blanc, bordure épaisse colorée,
% ombre dure encre, badge plein. Titre optionnel [#1].
% ============================================================
\tcbset{popbase/.style={breakable,enhanced,colback=white,boxrule=2.3pt,arc=9pt,
  shadow={3pt}{-3pt}{0pt}{popink},left=14pt,right=14pt,top=11pt,bottom=11pt,before skip=11pt,after skip=15pt,
  fontupper=\popsemi\fontsize{10.6}{14.5}\selectfont\color{popink},
  fontlower=\popsemi\fontsize{10.6}{14.5}\selectfont\color{popink}}}
% Coche dessinée (le glyphe ✓ n'existe pas dans les polices du document)
\newcommand{\popcheck}[1]{\tikz[baseline=-0.5ex]\draw[#1,line width=1.6pt,line cap=round,line join=round](-0.13,0.0)--(-0.04,-0.09)--(0.13,0.1);}
\providecommand{\coche}{\popcheck{popgreen}}
\providecommand{\resultat}[1]{\par\smallskip\noindent\fcolorbox{popgreen}{popgreenL}{\parbox{\dimexpr\linewidth-2\fboxsep-2\fboxrule\relax}{\textbf{Résultat.} #1}}\par\smallskip}
\newcommand{\popboxhead}[3]{\poppill{#1}{#2}{white}\ifx\relax#3\relax\else\hspace{6pt}{\popxbold\fontsize{10.6}{12}\selectfont\color{popink}#3}\fi\par\vspace{7pt}}

\newtcolorbox{aretenir}[1][]{popbase,colframe=popgreen,before upper={\popboxhead{À retenir}{popgreen}{#1}}}
\newtcolorbox{exemplebox}[1][]{popbase,colframe=poporange,before upper={\popboxhead{Exemple}{poporange}{#1}}}
\newtcolorbox{definition}[1][]{popbase,colframe=popblue,before upper={\popboxhead{Définition}{popblue}{#1}}}
\newtcolorbox{propriete}[1][]{popbase,colframe=poppurple,before upper={\popboxhead{Propriété}{poppurple}{#1}}}
\newtcolorbox{methode}[1][]{popbase,colframe=popgold,before upper={\popboxhead{Méthode}{popgold}{#1}}}
\newtcolorbox{attention}[1][]{popbase,colframe=poppink,before upper={\popboxhead{Attention}{poppink}{#1}}}
\newtcolorbox{experience}[1][]{popbase,colframe=popblue,colback=popblueL,before upper={\popboxhead{Expérience}{popblue}{#1}}}
% « Le savais-tu ? » : carte violette pleine (contexte, culture, Cameroun)
\newtcolorbox{savaistu}[1][]{popbase,colback=poppurple,colframe=popink,
  fontupper=\popsemi\fontsize{10.4}{14.2}\selectfont\color{white},
  before upper={\poppill{Le savais-tu\,?}{white}{poppurple}\ifx\relax#1\relax\else\hspace{6pt}{\popxbold\color{white}#1}\fi\par\vspace{7pt}}}
% Exercice résolu : énoncé en haut, \tcblower puis la solution sur fond crème
\newcounter{popexo}\newcounter{popexr}
\newtcolorbox{exoresolu}[1][]{popbase,colframe=poporange,colbacklower=popcream,
  segmentation style={popink,line width=1.2pt,dashed},
  before upper={\stepcounter{popexr}\popboxhead{Exercice résolu \thepopexr}{poporange}{#1}},
  before lower={\poppill{Solution}{popink}{popcream}\par\vspace{6pt}}}
% Exercice (non résolu en place) — la correction est dans la partie Corrigés
\newtcolorbox{exercice}[1][]{popbase,colframe=popink,
  before upper={\stepcounter{popexo}\popboxhead{Exercice \thepopexo}{popink}{#1}}}
% Corrigé
\newtcolorbox{corrige}[1][]{popbase,colframe=popgreen,colback=popgreenL,
  before upper={\popboxhead{Corrigé}{popgreen}{#1}}}
% Objectifs / prérequis / bilan
\newtcolorbox{objectifs}{popbase,colframe=popink,colback=poporangeL,
  before upper={\popboxhead{Objectifs de la leçon}{popink}{}}}
\newtcolorbox{prerequis}{popbase,colframe=popink,colback=popblueL,
  before upper={\popboxhead{Prérequis}{popblue}{}}}
\newtcolorbox{bilan}{popbase,colframe=popink,colback=white,boxrule=2.8pt,
  before upper={\popboxhead{Fiche bilan}{poporange}{L'essentiel à connaître pour l'examen}}}
% Figure encadrée avec légende
\newtcolorbox{popfigure}[1][]{enhanced,breakable=false,colback=white,colframe=popline,boxrule=1.4pt,arc=8pt,
  left=10pt,right=10pt,top=10pt,bottom=8pt,before skip=10pt,after skip=12pt,halign=center,
  lower separated=false,halign lower=center,
  fontlower=\popsemi\fontsize{9}{11.5}\selectfont\color{popmuted},
  #1}
\newcommand{\legende}[1]{\par\vspace{6pt}{\popsemi\fontsize{9}{11.5}\selectfont\color{popmuted}#1}}

% ============================================================
% Signature OnBuch+
% ============================================================
\newcommand{\poplogoinline}[1]{{\poplogo\fontsize{#1}{#1}\selectfont\color{popink}OnBuch\color{poporange}+}}
\newcommand{\signatureonbuch}{%
  \begin{tcolorbox}[enhanced,colback=popink,colframe=popink,boxrule=2.2pt,arc=10pt,
      drop shadow=poporange,left=14pt,right=14pt,top=10pt,bottom=10pt,before skip=0pt,after skip=0pt]
    \tikz[baseline=-0.7ex]\node[circle,fill=popgreen,draw=popcream,line width=1.3pt,minimum size=22pt,inner sep=0]{\popcheck{white}};%
    \hspace{9pt}\begin{minipage}[c]{0.62\linewidth}
      {\popxbold\fontsize{12}{13}\selectfont\color{popcream}Signé\ \textbullet\ {\poplogo OnBuch\color{poporange}+}}\par\vspace{2pt}
      {\raggedright\popsemi\fontsize{8}{10.5}\selectfont\color{popline}\MakeUppercase{Équipe pédagogique OnBuch+}\\\MakeUppercase{Conforme au programme officiel MINESEC}\par}
    \end{minipage}\hfill
    {\poplogo\fontsize{22}{22}\selectfont\color{popcream}OnBuch\color{poporange}+}
  \end{tcolorbox}%
}

% ============================================================
% Page de garde
% #1 titre · #2 matière · #3 niveau · #4 module · #5 n° leçon · #6 durée ·
% #7 description / objectifs (texte ou itemize)
% ============================================================
\newcommand{\pagegarde}[7]{%
  \thispagestyle{empty}
  \newgeometry{a4paper,margin=1.7cm,top=1.6cm,bottom=1.4cm}%
  \begin{tikzpicture}[remember picture,overlay]
    \fill[poporange!18] ([xshift=-3.2cm,yshift=-3.6cm]current page.north east) circle (4.6cm);
    \fill[poppurple!14] ([xshift=2.4cm,yshift=7.6cm]current page.south west) circle (3.8cm);
  \end{tikzpicture}%
  {\poplogo\fontsize{30}{30}\selectfont\color{popink}OnBuch\color{poporange}+}\hfill
  \raisebox{0.6ex}{\poppill{Cours signé \textbullet\ Premium}{popink}{popcream}}\par
  \vspace{1pt}\popsquiggle{poppurple}\par
  \vspace{8pt}
  \begin{tcolorbox}[enhanced,colback=poporange,colframe=popink,boxrule=2.8pt,arc=13pt,
      drop shadow=popink,left=18pt,right=18pt,top=12pt,bottom=13pt,after skip=10pt]
    \poppill{#2 \textbullet\ #3}{popink}{popcream}\hspace{5pt}\poppill{#4}{white}{popink}\par\vspace{8pt}
    {\popdisplay\fontsize{14}{14}\selectfont\color{popink}LEÇON #5}\par\vspace{4pt}
    {\raggedright\hyphenpenalty=10000\exhyphenpenalty=10000\popdisplay\fontsize{25}{27}\selectfont\color{popcream}\MakeUppercase{#1}\par}
    \vspace{8pt}
    {\popxbold\fontsize{9.5}{9.5}\selectfont\color{popink}\MakeUppercase{Durée conseillée : #6}}
  \end{tcolorbox}
  \begin{tcolorbox}[enhanced,colback=white,colframe=popink,boxrule=2.2pt,arc=10pt,
      drop shadow=popink,left=16pt,right=16pt,top=10pt,bottom=10pt,after skip=8pt]
    \poppill{Dans cette leçon}{poppurple}{white}\par\vspace{5pt}
    \popsemi\fontsize{9.3}{12.6}\selectfont\color{popink}#7
  \end{tcolorbox}
  \noindent\begin{minipage}[t]{0.487\linewidth}
    \begin{tcolorbox}[enhanced,colback=popgreen,colframe=popink,boxrule=2.2pt,arc=10pt,
        drop shadow=popink,left=11pt,right=11pt,top=10pt,bottom=10pt,height=3.1cm,valign=center]
      \tcbox[colback=white,colframe=popink,boxrule=1.3pt,arc=5pt,boxsep=0pt,nobeforeafter,
        left=3pt,right=3pt,top=3pt,bottom=3pt,valign=center]{\qrcode[height=1.9cm]{https://play.google.com/store/apps/details?id=com.luvvix.onbuch&hl=fr}}%
      \hspace{8pt}\begin{minipage}[c]{0.5\linewidth}
        {\popdisplay\fontsize{11}{12.5}\selectfont\color{white}\MakeUppercase{Télécharge OnBuch}}\par\vspace{4pt}
        {\raggedright\popsemi\fontsize{8.4}{11}\selectfont\color{white}Gratuit \textbullet\ Google Play\\Tuteur IA, annales, concours, résultats\par}
      \end{minipage}
    \end{tcolorbox}
  \end{minipage}\hfill
  \begin{minipage}[t]{0.487\linewidth}
    \begin{tcolorbox}[enhanced,colback=poppurple,colframe=popink,boxrule=2.2pt,arc=10pt,
        drop shadow=popink,left=11pt,right=11pt,top=10pt,bottom=10pt,height=3.1cm,valign=center]
      \tikz[baseline=-0.7ex]\node[circle,fill=white,draw=popink,line width=1.3pt,minimum size=1.6cm,inner sep=0]{\popdisplay\fontsize{24}{24}\selectfont\color{poppurple}+};%
      \hspace{8pt}\begin{minipage}[c]{0.6\linewidth}
        {\popdisplay\fontsize{11}{12.5}\selectfont\color{white}\MakeUppercase{Passe à OnBuch+}}\par\vspace{4pt}
        {\raggedright\popsemi\fontsize{8.4}{11}\selectfont\color{white}Tous les cours signés, Tuteur IA illimité, fiches PDF — onglet OnBuch+ de l'app\par}
      \end{minipage}
    \end{tcolorbox}
  \end{minipage}
  \vspace{8pt}
  \popcontenu
  \vfill
  \signatureonbuch
  \restoregeometry
  \newpage
}

% Rangée « Ce cours contient » (page de garde)
\newcommand{\poptile}[3]{% #1 couleur #2 gros texte #3 légende
  \begin{tcolorbox}[enhanced,colback=#1,colframe=popink,boxrule=2pt,arc=9pt,shadow={2.5pt}{-2.5pt}{0pt}{popink},
      left=6pt,right=6pt,top=9pt,bottom=9pt,halign=center,valign=center,height=1.75cm,width=\linewidth,nobeforeafter]
    {\popdisplay\fontsize{12.5}{13}\selectfont\color{white}\MakeUppercase{#2}}\par\vspace{3pt}
    {\centering\popsemi\fontsize{7.8}{9.5}\selectfont\color{white}#3\par}
  \end{tcolorbox}}
\newcommand{\popcontenu}{%
  \par\noindent{\popxboldls\fontsize{7.5}{7.5}\selectfont\color{popmuted}CE COURS SIGNÉ CONTIENT}\par\vspace{7pt}
  \noindent\begin{minipage}[t]{0.235\linewidth}\poptile{popblue}{Cours}{complet, illustré, pas à pas}\end{minipage}\hfill
  \begin{minipage}[t]{0.235\linewidth}\poptile{poporange}{Méthodes}{et exercices résolus}\end{minipage}\hfill
  \begin{minipage}[t]{0.235\linewidth}\poptile{poppink}{Exercices}{type examen + corrigés}\end{minipage}\hfill
  \begin{minipage}[t]{0.235\linewidth}\poptile{popgreen}{Fiche bilan}{l'essentiel à retenir}\end{minipage}\par}

% Sommaire
\newtcolorbox{sommaire}{enhanced,colback=white,colframe=popink,boxrule=2.2pt,arc=10pt,
  drop shadow=popink,left=18pt,right=18pt,top=13pt,bottom=13pt,before skip=6pt,after skip=20pt,
  before upper={\poppill{Sommaire}{poppurple}{white}\par\vspace{9pt}\popsemi\fontsize{10.6}{14.5}\selectfont\color{popink}}}

% Signature de fin de cours — poussée tout en bas de la dernière page
\newcommand{\findecours}{%
  \par\vfill\nopagebreak
  \begin{center}\popsquiggle{poporange}\end{center}\vspace{4pt}
  \signatureonbuch
}
\AtBeginDocument{\def\llbracket{\mathopen{[\mkern-3.5mu[}}\def\rrbracket{\mathclose{]\mkern-3.5mu]}}} % intervalles d'entiers (glyphe ⟦ absent de la police)
% Glyphes absents de Fira Math : redéfinis à partir de glyphes disponibles
\AtBeginDocument{%
  \def\setminus{\mathbin{\backslash}}\let\smallsetminus\setminus
  \def\vdots{\mathord{\vbox{\baselineskip3.2pt\lineskiplimit0pt\kern4pt\hbox{.}\hbox{.}\hbox{.}}}}%
  \def\ddots{\mathinner{\mkern1mu\raise6.5pt\hbox{.}\mkern2mu\raise3.6pt\hbox{.}\mkern2mu\raise0.7pt\hbox{.}\mkern1mu}}%
  \def\triangle{\mathord{\scalebox{0.9}{$\symup{\Delta}$}}}%
  \def\nabla{\mathord{\raisebox{\depth}{\scalebox{1}[-1]{$\symup{\Delta}$}}}}%
  \def\checkmark{\popcheck{popdark}}%
  \def\star{\mathbin{\ast}}\def\bigstar{\mathord{\ast}}%
  \def\diamond{\mathbin{\scalebox{0.6}{$\Diamond$}}}%
  \def\bigcup{\mathop{\mathchoice{\raisebox{-0.3em}{\scalebox{1.7}{$\cup$}}}{\scalebox{1.25}{$\cup$}}{\cup}{\cup}}\displaylimits}%
  \def\bigcap{\mathop{\mathchoice{\raisebox{-0.3em}{\scalebox{1.7}{$\cap$}}}{\scalebox{1.25}{$\cap$}}{\cap}{\cap}}\displaylimits}%
  \def\lesseqqgtr{\mathrel{\lessgtr}}\def\bigoplus{\mathop{\oplus}\displaylimits}%
}
\providecommand{\ssi}{si et seulement si} % abréviation fréquente
\AtBeginDocument{\providecommand{\wideparen}{\overparen}} % arc de cercle

%% FIGFIX-BEGIN v4 — correctifs globaux des figures (docs/onbuch-generation/tools/figfix)
\usepackage{adjustbox}\usepackage{etoolbox}
% toute figure TikZ/pgfplots plus large que la ligne est réduite à la largeur disponible
\BeforeBeginEnvironment{tikzpicture}{\begin{adjustbox}{max width=\linewidth}}
\AfterEndEnvironment{tikzpicture}{\end{adjustbox}}
% légendes pgfplots lisibles par-dessus les courbes
\pgfplotsset{every axis/.append style={legend style={fill=white,fill opacity=0.92,text opacity=1,draw=black!15,inner sep=3pt}}}
% étiquettes d'arêtes (syntaxe edge["..."]) avec fond blanc
\tikzset{every edge quotes/.append style={fill=white,inner sep=1.2pt}}
%% FIGFIX-END
\begin{document}

\pagegarde{L'adaptateur secteur}{PCT}{3ème}{Module 2 : Électricité}{N15}{1 séance (1 h chacune)}{Cette leçon étudie l'adaptateur secteur, appareil omniprésent dans la vie courante, qui abaisse et redresse la tension du réseau pour alimenter les appareils électroniques. Elle met en évidence ses deux fonctions essentielles — transformation et redressement — à partir d'observations expérimentales simples, et insiste sur les règles de sécurité liées à son utilisation.
\vspace{4pt}
\begin{multicols}{2}\small
\begin{itemize}
\item À la fin de cette leçon, je sais définir un adaptateur secteur et citer des exemples d'utilisation dans la vie courante
\item À la fin de cette leçon, je sais expliquer pourquoi la tension du secteur ne peut pas alimenter directement un appareil électronique
\item À la fin de cette leçon, je sais décrire le rôle du transformateur (abaissement de la tension alternative)
\item À la fin de cette leçon, je sais expliquer la fonction de redressement (conversion alternatif → continu) à partir d'observations à l'oscilloscope ou au voltmètre
\item À la fin de cette leçon, je sais identifier sur la plaque signalétique d'un adaptateur les grandeurs d'entrée et de sortie (tension, nature du courant)
\item À la fin de cette leçon, je sais schématiser la chaîne de conversion : secteur → transformateur → redresseur → appareil
\end{itemize}\end{multicols}}

\begin{sommaire}
\begin{multicols}{2}
\begin{enumerate}
\item Pourquoi un adaptateur ? Fonction générale
\item Première fonction : la transformation (abaissement de tension)
\item Deuxième fonction : le redressement (alternatif → continu)
\item Schéma de principe et chaîne de conversion de l'adaptateur
\item Sécurité et bon usage des adaptateurs
\item Activité d'intégration
\item Méthodes et exercices résolus
\item Exercices
\item Corrigés des exercices
\item Fiche bilan
\end{enumerate}
\end{multicols}
\end{sommaire}

\begin{objectifs}
\begin{itemize}
\item À la fin de cette leçon, je sais définir un adaptateur secteur et citer des exemples d'utilisation dans la vie courante
\item À la fin de cette leçon, je sais expliquer pourquoi la tension du secteur ne peut pas alimenter directement un appareil électronique
\item À la fin de cette leçon, je sais décrire le rôle du transformateur (abaissement de la tension alternative)
\item À la fin de cette leçon, je sais expliquer la fonction de redressement (conversion alternatif → continu) à partir d'observations à l'oscilloscope ou au voltmètre
\item À la fin de cette leçon, je sais identifier sur la plaque signalétique d'un adaptateur les grandeurs d'entrée et de sortie (tension, nature du courant)
\item À la fin de cette leçon, je sais schématiser la chaîne de conversion : secteur → transformateur → redresseur → appareil
\item À la fin de cette leçon, je sais appliquer les consignes de sécurité liées à l'utilisation des adaptateurs secteur
\item À la fin de cette leçon, je sais choisir un adaptateur compatible avec un appareil donné à partir des indications portées sur les deux
\end{itemize}
\end{objectifs}

\begin{prerequis}
\begin{itemize}
\item Tension continue et tension alternative : distinction et symboles (leçon sur les tensions)
\item Utilisation du voltmètre (positions AC et DC)
\item Le secteur délivre une tension alternative de valeur efficace 220 V au Cameroun
\item Notions de circuit électrique simple et de sécurité électrique (dangers du courant)
\item Lecture d'un oscillogramme simple (leçon sur les tensions périodiques)
\end{itemize}
\end{prerequis}

\newpage

\coursec{Pourquoi un adaptateur ? Fonction générale}

À Yaoundé comme à Douala, chaque famille possède plusieurs petits boîtiers branchés sur les prises : chargeurs de téléphone, adaptateurs de radio, de lampe rechargeable, bloc d'alimentation d'ordinateur portable. Pourtant, le secteur ENEO délivre une tension \emph{alternative} de \SI{220}{V}, dangereuse, alors qu'un téléphone portable n'a besoin que de \SI{5}{V} \emph{continus}. Comment ce petit boîtier noir transforme-t-il une tension élevée et alternative en une tension faible et continue ? Que risque-t-on si l'on branche un adaptateur inadapté ou défectueux ? C'est ce que nous allons découvrir.

\courssub{Le constat : deux mondes électriques incompatibles}

\textbf{Observation de la vie courante}

Regardons autour de nous et relevons les indications portées sur quelques appareils électriques de la maison :

\begin{itemize}
\item un téléphone portable se recharge sous \SI{5}{V} continus ;
\item une radio à piles fonctionne sous \SI{6}{V} ou \SI{9}{V} continus ;
\item une lampe torche à DEL fonctionne sous \SI{3}{V} continus ;
\item un ordinateur portable est alimenté sous \SI{12}{V} ou \SI{19}{V} continus ;
\item une petite lampe de bicyclette est prévue pour \SI{6}{V}.
\end{itemize}

Or la prise murale, alimentée par le réseau ENEO, délivre une tension de \SI{220}{V}, et cette tension est \cle{alternative} : elle change de sens un très grand nombre de fois par seconde (50 fois par seconde au Cameroun).

\begin{aretenir}[Le problème posé]
Les petits appareils électroniques fonctionnent avec une \cle{tension continue} de \cle{faible valeur} (quelques volts), alors que le secteur fournit une \cle{tension alternative} de \cle{forte valeur} (\SI{220}{V}). Il y a donc un double écart :
\begin{itemize}
\item un écart de \textbf{valeur} : \SI{220}{V} contre quelques volts ;
\item un écart de \textbf{nature} : courant alternatif contre courant continu.
\end{itemize}
\end{aretenir}

\textbf{Expérience de mise en évidence : peut-on brancher directement ?}

\begin{experience}[Brancher une lampe de 6 V sur le secteur : une idée dangereuse]
\textbf{Question.} Une élève possède une petite lampe de bicyclette prévue pour \SI{6}{V}. Elle se demande si elle peut la brancher directement sur la prise de courant de la maison.

\textbf{Hypothèse.} « La lampe va simplement briller plus fort. »

\textbf{Expérience.} Cette expérience est \textbf{interdite en vraie grandeur} car elle est mortelle pour l'expérimentateur. On la simule en toute sécurité : on branche une lampe de \SI{6}{V} sur un générateur de laboratoire dont on augmente progressivement la tension.

\textbf{Observations.}
\begin{itemize}
\item À \SI{6}{V}, la lampe brille normalement.
\item À \SI{12}{V}, elle brille très fort, le filament devient blanc.
\item Au-delà, le filament \textbf{fond} : la lampe est détruite instantanément, parfois avec projection de verre.
\end{itemize}

\textbf{Interprétation.} Une lampe prévue pour \SI{6}{V} ne supporte pas une tension très supérieure : le courant qui la traverse devient trop intense et détruit le filament par échauffement. Branchée sur \SI{220}{V}, elle serait détruite en une fraction de seconde, avec un risque d'incendie et d'électrocution.

\textbf{Conclusion.} Il est \textbf{impossible et dangereux} de brancher directement un petit appareil basse tension sur le secteur. Il faut un \cle{appareil intermédiaire} qui adapte la tension : c'est l'\cle{adaptateur secteur}.
\end{experience}

\begin{attention}[Danger absolu]
Ne jamais, sous aucun prétexte, brancher un appareil prévu pour quelques volts (lampe de \SI{6}{V}, téléphone, radio à piles) directement sur une prise de \SI{220}{V}. Les conséquences sont : destruction immédiate de l'appareil, projection de débris, incendie, et surtout \textbf{électrocution} de la personne qui manipule. La tension de \SI{220}{V} peut tuer.
\end{attention}

\courssub{Définition et fonction générale de l'adaptateur secteur}

\begin{definition}[Adaptateur secteur]
L'\cle{adaptateur secteur} est un appareil qui \textbf{convertit la tension alternative du secteur} (\SI{220}{V}) \textbf{en une tension continue de faible valeur} (\SI{3}{V}, \SI{5}{V}, \SI{9}{V}, \SI{12}{V}\ldots) adaptée au fonctionnement d'un appareil électrique ou électronique.
\end{definition}

L'adaptateur remplit donc une \cle{fonction générale} : \textbf{adapter} l'énergie électrique du réseau aux besoins de l'appareil. Cette fonction générale se décompose en deux fonctions techniques principales que nous allons étudier en détail :

\begin{enumerate}
\item la \cle{transformation} : abaisser la tension de \SI{220}{V} à quelques volts (grâce à un \emph{transformateur}) ;
\item le \cle{redressement} : convertir la tension alternative en tension continue (grâce à des composants appelés \emph{diodes}).
\end{enumerate}

\begin{exemplebox}[Exemples familiers d'adaptateurs]
\begin{itemize}
\item le \textbf{chargeur de téléphone} : il reçoit \SI{220}{V} alternatifs et fournit \SI{5}{V} continus ;
\item le \textbf{bloc d'alimentation d'ordinateur portable} : il fournit souvent \SI{19}{V} continus ;
\item l'\textbf{adaptateur de radio} ou de lampe rechargeable : il fournit \SI{6}{V}, \SI{9}{V} ou \SI{12}{V} continus ;
\item le petit boîtier noir de la télévision par satellite ou du modem internet.
\end{itemize}
\end{exemplebox}

\begin{popfigure}
\begin{center}
\begin{tikzpicture}[font=\small, >=Stealth]
% Prise secteur
\draw[thick, popdark, fill=popblueL] (0,0) rectangle (2.2,2);
\fill[popdark] (0.8,1) circle (0.12);
\fill[popdark] (1.4,1) circle (0.12);
\node[align=center] at (1.1,2.4) {\textbf{Prise secteur}};
\node[align=center, popblue] at (1.1,-0.5) {\SI{220}{V} $\sim$\\ 50 Hz};
% Flèche 1
\draw[->, very thick, poporange] (2.4,1) -- (3.8,1);
% Adaptateur
\draw[thick, popdark, fill=poporangeL, rounded corners=4pt] (4,0) rectangle (7,2);
\node[align=center] at (5.5,1) {\textbf{Adaptateur}\\ \textbf{secteur}};
\node[align=center, poporange] at (5.5,-0.5) {transformation\\ + redressement};
% Flèche 2
\draw[->, very thick, popgreen] (7.2,1) -- (8.6,1);
% Téléphone
\draw[thick, popdark, fill=popgreenL, rounded corners=6pt] (8.8,-0.2) rectangle (10.6,2.2);
\draw[popdark] (9.1,0.2) rectangle (10.3,1.7);
\node[align=center] at (9.7,2.7) {\textbf{Téléphone}};
\node[align=center, popgreen!60!black] at (9.7,-0.8) {\SI{5}{V} \text{continu}};
\end{tikzpicture}
\end{center}
\legende{Chaîne de conversion : l'adaptateur secteur reçoit la tension alternative de \SI{220}{V} de la prise ENEO et fournit au téléphone une tension continue de \SI{5}{V}.}
\end{popfigure}

\courssub{Lire la plaque signalétique d'un adaptateur}

Tout adaptateur porte, gravée ou imprimée sur son boîtier, une \cle{plaque signalétique} qui résume ses caractéristiques électriques. Savoir la lire est indispensable pour utiliser l'appareil correctement et en sécurité.

\begin{popfigure}
\begin{center}
\begin{tikzpicture}[font=\small]
\draw[thick, popdark, fill=popcream, rounded corners=4pt] (0,0) rectangle (9,4.6);
\node[align=center, popdark] at (4.5,4.1) {\textbf{ADAPTATEUR SECTEUR -- Mod\`ele ONB-05}};
\draw[popline] (0.4,3.7) -- (8.6,3.7);
\node[align=left, anchor=west] at (0.6,3.0) {\textbf{INPUT (entr\'ee) :} \SI{220}{V} $\sim$ \quad \SI{50}{Hz}};
\node[align=left, anchor=west] at (0.6,2.1) {\textbf{OUTPUT (sortie) :} \SI{5}{V} \text{continu} \quad \SI{1}{A}};
\node[align=left, anchor=west, popmuted] at (0.6,1.2) {Ne pas ouvrir -- Usage int\'erieur uniquement};
\node[align=left, anchor=west, popmuted] at (0.6,0.5) {Fabriqu\'e selon les normes en vigueur};
\end{tikzpicture}
\end{center}
\legende{Plaque signalétique schématique d'un adaptateur : la ligne INPUT décrit ce que l'adaptateur \emph{reçoit} du secteur ; la ligne OUTPUT décrit ce qu'il \emph{fournit} à l'appareil.}
\end{popfigure}

\begin{methode}[Lire la plaque signalétique d'un adaptateur]
\begin{enumerate}
\item Repérer la ligne \textbf{INPUT} (en français : \emph{entrée}) : elle indique la tension que l'adaptateur \textbf{reçoit} du secteur, par exemple \SI{220}{V} $\sim$ \SI{50}{Hz}.
\item Repérer la ligne \textbf{OUTPUT} (en français : \emph{sortie}) : elle indique la tension que l'adaptateur \textbf{fournit} à l'appareil, par exemple \SI{5}{V} continu, ainsi que l'intensité maximale qu'il peut débiter (par exemple \SI{1}{A}).
\item Identifier la \textbf{nature du courant} grâce aux symboles : le symbole $\sim$ signifie \emph{alternatif} ; le symbole d'un trait plein surmontant un trait pointillé signifie \emph{continu}.
\item Vérifier la \textbf{compatibilité} : la tension de sortie (OUTPUT) doit correspondre exactement à la tension prévue pour l'appareil à alimenter.
\end{enumerate}
\end{methode}

\begin{aretenir}[Les deux symboles à connaître]
\begin{itemize}
\item Le symbole $\sim$ (petite vague) désigne une tension ou un courant \cle{alternatif} : c'est le cas du secteur ENEO.
\item Le symbole d'un trait continu au-dessus d'un trait pointillé (--- $\cdot\cdot\cdot$) désigne une tension ou un courant \cle{continu} : c'est le cas de la sortie de l'adaptateur, des piles et des batteries.
\end{itemize}
\end{aretenir}

\begin{exoresolu}[Lecture d'une plaque signalétique]
Énoncé. Sur le chargeur d'un téléphone, on lit : INPUT : \SI{220}{V} $\sim$ \SI{50}{Hz} ; OUTPUT : \SI{5}{V} continu, \SI{2}{A}.
\begin{enumerate}
\item Quelle tension ce chargeur reçoit-il du secteur ? De quelle nature est-elle ?
\item Quelle tension fournit-il au téléphone ? De quelle nature est-elle ?
\item Que signifie l'indication \SI{2}{A} ?
\end{enumerate}
\tcblower
Solution détaillée.
\begin{enumerate}
\item Le chargeur reçoit du secteur une tension de \SI{220}{V} de nature \textbf{alternative} (symbole $\sim$), de fréquence \SI{50}{Hz}. C'est la tension normale du réseau ENEO.
\item Il fournit au téléphone une tension de \SI{5}{V} de nature \textbf{continue} (trait plein sur trait pointillé). L'adaptateur a donc abaissé la tension (de \SI{220}{V} à \SI{5}{V}) et changé sa nature (d'alternative en continue).
\item L'indication \SI{2}{A} est l'\textbf{intensité maximale} du courant que le chargeur peut débiter : il peut alimenter tout appareil réclamant au plus \SI{2}{A} sous \SI{5}{V}.
\end{enumerate}
\end{exoresolu}

\begin{attention}[Erreur fréquente : confondre entrée et sortie]
Beaucoup d'élèves confondent la \textbf{tension d'entrée} (\SI{220}{V}) et la \textbf{tension de sortie} (par exemple \SI{5}{V}). Retiens bien :
\begin{itemize}
\item ce n'est \textbf{pas} l'appareil (téléphone, radio) qui reçoit \SI{220}{V} : il ne reçoit que la faible tension de sortie ;
\item ce n'est \textbf{pas} non plus l'adaptateur qui « consomme » \SI{220}{V} pour son propre fonctionnement : il \emph{reçoit} \SI{220}{V} du secteur et les \emph{convertit} ;
\item brancher un appareil prévu pour \SI{5}{V} sur un adaptateur qui sort \SI{12}{V} le détruit, même si l'adaptateur est de bonne qualité.
\end{itemize}
Au BEPC, quand on te donne une plaque signalétique, cite toujours les deux tensions en précisant leur rôle : entrée (reçue) et sortie (fournie).
\end{attention}

\begin{savaistu}[Le secteur ENEO au Cameroun]
Au Cameroun, la société ENEO distribue l'énergie électrique sous une tension alternative de \SI{220}{V} et de fréquence \SI{50}{Hz}, comme dans la plupart des pays d'Afrique et d'Europe. C'est pourquoi les adaptateurs vendus à Yaoundé, Douala ou Garoua indiquent presque toujours INPUT : \SI{220}{V} $\sim$ \SI{50}{Hz}. Attention en voyage : dans certains pays (comme les États-Unis), le secteur délivre \SI{110}{V} à \SI{60}{Hz} ; il faut alors vérifier que l'adaptateur accepte cette tension d'entrée, sinon il peut être endommagé. De nombreux chargeurs modernes sont dits « universels » : leur plaque indique INPUT : 100--\SI{240}{V} $\sim$, ce qui leur permet de fonctionner partout dans le monde.
\end{savaistu}

\courssub{Première fonction : la transformation (abaissement de la tension)}

\textbf{Le composant clé : le transformateur}

À l'intérieur de l'adaptateur (souvent le gros bloc lourd), se trouve un \cle{transformateur}. C'est un composant statique (sans pièces mobiles) qui ne fonctionne \textbf{qu'avec du courant alternatif}.

\begin{experience}[Le transformateur : abaisser la tension alternative]
\textbf{Matériel.} Un transformateur de laboratoire (primaire 220 V, secondaire 6 V ou 12 V), un voltmètre alternatif, des fils de connexion.

\textbf{Protocole.}
\begin{enumerate}
\item Brancher le primaire du transformateur sur le secteur \SI{220}{V} $\sim$ (prudence !).
\item Mesurer la tension aux bornes du secondaire avec le voltmètre en mode alternatif.
\item Relever la valeur.
\end{enumerate}

\textbf{Observation.} La tension mesurée aux bornes du secondaire est faible (ex. \SI{6}{V} $\sim$ ou \SI{12}{V} $\sim$). Elle est toujours \textbf{alternative} (l'aiguille du voltmètre ou l'affichage oscille si on mesure une tension alternative pure, ou le symbole $\sim$ s'affiche).

\textbf{Interprétation.} Le transformateur a abaissé la tension. Il est composé de deux enroulements de fil de cuivre (le \emph{primaire} et le \emph{secondaire}) enroulés sur un même noyau en fer doux (feuilles de tôle isolées pour limiter les courants de Foucault). Le courant alternatif du primaire crée un champ magnétique alternatif dans le noyau, qui induit une tension alternative dans le secondaire.

\textbf{Propriété (Rapport de transformation).} Pour un transformateur idéal (sans pertes) :
\[
\frac{U_{\text{secondaire}}}{U_{\text{primaire}}} = \frac{N_{\text{secondaire}}}{N_{\text{primaire}}}
\]
où $U$ est la tension efficace (en V) et $N$ le nombre de spires.
Comme $N_{\text{secondaire}} < N_{\text{primaire}}$, on a $U_{\text{secondaire}} < U_{\text{primaire}}$ : c'est un \textbf{transformateur abaisseur}.

\textbf{Conclusion.} Le transformateur réalise la \textbf{fonction de transformation} : il abaisse la tension alternative de \SI{220}{V} à une tension alternative de faible valeur (ex. \SI{9}{V} $\sim$). Il assure aussi l'\textbf{isolement galvanique} entre le secteur (dangereux) et la sortie (sécurité).
\end{experience}

\begin{propriete}[Transformateur abaisseur]
Un transformateur abaisseur a moins de spires au secondaire qu'au primaire ($N_2 < N_1$). Il fournit une tension alternative secondaire $U_2$ plus faible que la tension primaire $U_1$ : $U_2 < U_1$. L'intensité secondaire $I_2$ est plus grande que l'intensité primaire $I_1$ (conservation de l'énergie : $U_1 I_1 \approx U_2 I_2$).
\end{propriete}

\begin{attention}[Le transformateur ne change pas la nature du courant !]
À la sortie du transformateur, la tension est \textbf{toujours alternative} (symbole $\sim$). Elle a juste une valeur plus faible. Il ne faut pas confondre « tension faible » et « tension continue ». Une tension alternative de \SI{9}{V} peut encore faire clignoter une DEL ou détruire un circuit électronique prévu pour du continu. Il faut une deuxième étape : le redressement.
\end{attention}

\courssub{Deuxième fonction : le redressement (alternatif $\rightarrow$ continu)}

\textbf{Le composant clé : la diode}

Une \cle{diode} est un composant électronique qui ne laisse passer le courant que dans \textbf{un seul sens}. Elle agit comme un clapet anti-retour.

\begin{definition}[Diode — symbole et convention]
Le symbole de la diode est un triangle pointant vers une barre verticale. Le courant ne peut circuler que dans le sens de la flèche (de l'anode vers la cathode).
\begin{center}
\begin{tikzpicture}[font=\small, >=Stealth]
\draw[thick] (0,0) -- (1,1) -- (0,2) -- cycle; % triangle
\draw[thick] (1,1) -- (1.5,1); % sortie
\draw[thick] (0,1) -- (-0.5,1); % entrée
\draw[very thick] (1,0) -- (1,2); % barre
\node at (-0.8,1) {A (anode)};
\node at (1.8,1) {K (cathode)};
\draw[->, thick, poporange] (0.5,1.3) -- (0.5,0.7);
\node[poporange] at (0.5,1.5) {Sens passant};
\end{tikzpicture}
\end{center}
\end{definition}

\textbf{Redressement simple (alternance positive seulement)}

Si on place une diode en série avec la sortie du transformateur :
\begin{itemize}
\item Pendant l'alternance positive de la tension alternative, la diode est \emph{passante} : le courant circule, la tension aux bornes de la charge suit la tension du générateur.
\item Pendant l'alternance négative, la diode est \emph{bloquée} : aucun courant ne circule, la tension aux bornes de la charge est nulle.
\end{itemize}
On obtient une tension \textbf{unilatérale} (toujours positive) mais \textbf{pulsée} (elle tombe à zéro 50 fois par seconde). Ce n'est pas encore une tension continue stable.

\textbf{Redressement double alternance (pont de Graetz) — utilisé dans les adaptateurs}

Pour utiliser les deux alternances, on utilise \textbf{quatre diodes} montées en \textbf{pont} (pont de Graetz).

\begin{popfigure}
\begin{center}
\begin{tikzpicture}[font=\small, >=Stealth]
% Transformateur secondaire
\draw[thick] (0,2) -- (0,0);
\draw[thick] (3,2) -- (3,0);
\node at (1.5,2.4) {Secondaire transfo ($\sim$)};
% 4 diodes
\draw[thick] (0,1) -- (0.5,1);
\draw[thick] (3,1) -- (2.5,1);
% Diode haut gauche (D1)
\draw[thick] (0.5,1.5) -- (1,1) -- (0.5,0.5) -- cycle;
\draw[very thick] (1,0.5) -- (1,1.5);
\draw[thick] (1,1) -- (1.5,1);
% Diode haut droite (D2)
\draw[thick] (2.5,1.5) -- (2,1) -- (2.5,0.5) -- cycle;
\draw[very thick] (2,0.5) -- (2,1.5);
\draw[thick] (1.5,1) -- (2,1);
% Diode bas gauche (D3)
\draw[thick] (0.5,1.5) -- (1,2) -- (0.5,2.5) -- cycle;
\draw[very thick] (1,2) -- (1,2.5);
\draw[thick] (1,2) -- (1.5,2);
% Diode bas droite (D4)
\draw[thick] (2.5,1.5) -- (2,2) -- (2.5,2.5) -- cycle;
\draw[very thick] (2,2) -- (2,2.5);
\draw[thick] (1.5,2) -- (2,2);
% Charge
\draw[thick] (1.5,1) -- (1.5,0.5) node[ground] {};
\draw[thick] (1.5,2) -- (1.5,2.5);
\draw[thick, poporange] (1.5,0.5) -- (2.5,0.5) node[right] {+};
\draw[thick, popblue] (1.5,2.5) -- (2.5,2.5) node[right] {-};
\node[poporange] at (2.5,0.2) {Sortie +};
\node[popblue] at (2.5,2.8) {Sortie -};
\end{tikzpicture}
\end{center}
\legende{Schéma de principe d'un pont de Graetz (redresseur double alternance). Les flèches indiquent le sens de circulation du courant dans la charge pour les deux alternances.}
\end{popfigure}

\begin{methode}[Fonctionnement du pont de Graetz]
\begin{enumerate}
\item \textbf{Alternance positive} (borne haute du transfo +, borne basse -) : Les diodes D1 et D4 sont passantes ; D2 et D3 bloquées. Le courant traverse la charge de gauche à droite (du + vers le -).
\item \textbf{Alternance négative} (borne haute du transfo -, borne basse +) : Les diodes D2 et D3 sont passantes ; D1 et D4 bloquées. Le courant traverse \textbf{toujours} la charge de gauche à droite.
\end{enumerate}
\textbf{Résultat :} La tension aux bornes de la charge est toujours de la même polarité. On a redressé les deux alternances.
\end{methode}

\textbf{Lissage (filtrage) — pour obtenir une tension continue stable}

La tension redressée est encore pulsée (elle varie entre 0 et le maximum). Pour la rendre \textbf{continue} (quasi plate), on ajoute un \cle{condensateur de lissage} (ou de filtrage) en parallèle sur la charge.

\begin{popfigure}
\begin{center}
\begin{tikzpicture}[font=\small, >=Stealth]
\begin{axis}[
width=\linewidth, height=5cm,
axis lines=middle,
xlabel={Temps $t$}, ylabel={Tension $u$},
xtick=\empty, ytick=\empty,
domain=0:4*pi, samples=400,
xmin=0, xmax=13, ymin=-1.5, ymax=3.5,
clip=false
]
% Alternative (pointillés)
\addplot[popmuted, dashed, thick] {sin(deg(x))};
\node[popmuted, anchor=south west] at (axis cs:0.5,1) {$u_{\text{alt}}$};
% Redressée simple (plein trait orange)
\addplot[poporange, thick] {max(sin(deg(x)),0)};
\node[poporange, anchor=south west] at (axis cs:0.5,1.5) {$u_{\text{redressée simple}}$};
% Redressée double + lissage (vert)
\addplot[popgreen, very thick, domain=0:12.56, samples=500] {0.64 + 0.08*cos(2*deg(x))}; % approximation visuelle lissage
\node[popgreen, anchor=south west] at (axis cs:0.5,2.5) {$u_{\text{lissée}}$};
\end{axis}
\end{tikzpicture}
\end{center}
\legende{Évolution de la tension : alternative $\rightarrow$ redressée simple $\rightarrow$ redressée double alternance lissée (quasi continue).}
\end{popfigure}

\begin{aretenir}[Les trois étapes de la conversion à l'intérieur de l'adaptateur]
\begin{enumerate}
\item \textbf{Transformation} : Transformateur abaisseur $\rightarrow$ Tension alternative faible (ex. \SI{9}{V} $\sim$).
\item \textbf{Redressement} : Pont de diodes (Graetz) $\rightarrow$ Tension unilatérale pulsée (toujours positive, mais qui tombe à 0).
\item \textbf{Lissage (Filtrage)} : Condensateur $\rightarrow$ Tension \textbf{continue} quasi constante (ex. \SI{9}{V} ---).
\end{enumerate}
\textit{Remarque :} Les adaptateurs modernes (chargeurs de téléphone) utilisent une technologie « à découpage » (plus légère, plus compacte) qui réalise ces mêmes fonctions à haute fréquence, mais le principe logique reste : Transformation $\rightarrow$ Redressement $\rightarrow$ Lissage/Régulation.
\end{aretenir}

\courssub{Schéma de principe et chaîne de conversion de l'adaptateur}

On résume l'adaptateur par un \textbf{schéma-blocs} (chaîne de conversion d'énergie) :

\begin{popfigure}
\begin{center}
\begin{tikzpicture}[font=\small, >=Stealth, node distance=1.2cm and 0.5cm]
% Styles
\tikzset{bloc/.style={draw, thick, rounded corners, fill=popblueL, minimum width=2.5cm, minimum height=1.2cm, align=center},
fleche/.style={->, very thick, popdark},
texte/.style={font=\small, align=center}}
% Nœuds
\node[bloc, align=center] (secteur){SECTEUR ENEO\\ \SI{220}{V} $\sim$ \SI{50}{Hz}};
\node[bloc, right=of secteur, align=center] (transfo){TRANSFORMATEUR\\ (Abaissement)\\ \SI{220}{V} $\sim$ $\rightarrow$ \SI{9}{V} $\sim$};
\node[bloc, right=of transfo, align=center] (redresseur){REDRESSEUR\\ (Pont de diodes)\\ $\sim$ $\rightarrow$ Unilatérale pulsée};
\node[bloc, right=of redresseur, align=center] (filtre){FILTRE / LISSAGE\\ (Condensateur)\\ Pulsée $\rightarrow$ Continue};
\node[bloc, right=of filtre, align=center] (sortie){SORTIE\\ \SI{9}{V} ---};
% Flèches
\draw[fleche] (secteur.east) -- (transfo.west) node[midway, above, texte] {Entrée};
\draw[fleche] (transfo.east) -- (redresseur.west) node[midway, above, texte] {Alt. faible};
\draw[fleche] (redresseur.east) -- (filtre.west) node[midway, above, texte] {Pulsée};
\draw[fleche] (filtre.east) -- (sortie.west) node[midway, above, texte] {Continue};
\end{tikzpicture}
\end{center}
\legende{Chaîne de conversion complète d'un adaptateur secteur linéaire classique.}
\end{popfigure}

\begin{propriete}[Chaîne de conversion de l'adaptateur secteur]
L'adaptateur secteur est un \textbf{convertisseur} d'énergie électrique. Sa chaîne de conversion comporte quatre maillons :
\begin{enumerate}
\item \textbf{Entrée} : Tension alternative secteur (\SI{220}{V} $\sim$).
\item \textbf{Transformation} : Transformateur abaisseur $\rightarrow$ Tension alternative de sécurité (TBTS, < 50 V $\sim$).
\item \textbf{Redressement} : Pont de diodes (Graetz) $\rightarrow$ Tension unilatérale pulsée.
\item \textbf{Lissage / Régulation} : Condensateur (+ circuit régulateur) $\rightarrow$ \textbf{Sortie} : Tension continue stabilisée (ex. \SI{5}{V} ---).
\end{enumerate}
\end{propriete}

\courssub{Sécurité et bon usage des adaptateurs}

\begin{aretenir}[Règles de sécurité indispensables]
\begin{enumerate}
\item \textbf{Vérifier la tension de sortie (OUTPUT)} : Elle doit correspondre \textbf{exactement} à la tension de l'appareil (ex. \SI{5}{V} pour un téléphone). Une tension supérieure détruit l'appareil ; une tension inférieure ne le fait pas fonctionner ou l'endommage.
\item \textbf{Vérifier l'intensité maximale (OUTPUT)} : L'adaptateur doit pouvoir fournir \textbf{au moins} l'intensité demandée par l'appareil (ex. téléphone \SI{1,5}{A} $\rightarrow$ adaptateur \SI{2}{A} minimum). Un adaptateur trop faible surchauffe et peut prendre feu.
\item \textbf{Vérifier la polarité de la fiche} : Le schéma sur l'adaptateur (souvent un cercle avec un + au centre et - à l'extérieur, ou l'inverse) doit correspondre à celui de l'appareil. Inverser la polarité détruit l'électronique.
\item \textbf{Ne pas couvrir l'adaptateur} pendant la charge (sur un lit, sous un oreiller) : il dégage de la chaleur (pertes par effet Joule dans le transformateur et les composants). Risque d'incendie.
\item \textbf{Ne pas utiliser un adaptateur déformé, qui bourdonne anormalement, qui sent le brûlé ou dont le fil est dénudé.} Le jeter et le remplacer par un modèle certifié (marque CE, normes).
\item \textbf{Débrancher l'adaptateur de la prise} quand il ne charge rien : il consomme quand même un peu d'énergie (veillée) et peut surchauffer en cas de surtension (orages fréquents au Cameroun).
\end{enumerate}
\end{aretenir}

\begin{attention}[Adaptateurs « génériques » ou contrefaits]
Sur les marchés (Marché Central, Mokolo, etc.), on trouve des adaptateurs très bon marché, sans marque, légers (pas de transformateur en fer, électronique de découpage bas de gamme). Ils ne respectent souvent pas l'isolement secteur/sortie : en cas de défaut, le \SI{220}{V} peut se retrouver sur la sortie \SI{5}{V} $\rightarrow$ \textbf{électrocution mortelle} en touchant le téléphone en charge. Privilégier les adaptateurs d'origine ou de marques reconnues, vendus chez des revendeurs agréés.
\end{attention}

\begin{savaistu}[Le chargeur « universel » à découpage (SMPS)]
Les anciens adaptateurs (lourds) contiennent un gros transformateur 50 Hz. Les chargeurs modernes de téléphone (légers, compacts) sont des \textbf{alimentations à découpage} (Switch Mode Power Supply).
\begin{enumerate}
\item Le \SI{220}{V} $\sim$ est redressé et lissé \textbf{directement} en \SI{300}{V} --- (dangereux !).
\item Un transistor commute cette haute tension à haute fréquence (\SI{50}{kHz} à \SI{100}{kHz}) dans un \textbf{petit transformateur} à ferrite (très léger).
\item La sortie est redressée, lissée et régulée.
\end{enumerate}
Avantages : légers, universels (INPUT 100--\SI{240}{V} $\sim$), rendement élevé. Inconvénients : plus complexes, parasites haute fréquence, danger interne \SI{300}{V} --- (jamais ouvrir !).
\end{savaistu}

\begin{exercice}[À toi de jouer]
\begin{enumerate}
\item Pourquoi ne peut-on pas brancher directement une radio prévue pour \SI{9}{V} continus sur une prise de courant ?
\item Sur l'adaptateur d'une lampe rechargeable, on lit : INPUT : \SI{220}{V} $\sim$ ; OUTPUT : \SI{12}{V} continu, \SI{0,5}{A}. Recopie et complète : « L'adaptateur reçoit une tension \ldots de valeur \ldots et fournit une tension \ldots de valeur \ldots. »
\item Un élève affirme : « Mon téléphone fonctionne sous \SI{220}{V} puisqu'il est branché sur la prise. » Corrige cette erreur en deux phrases.
\item \textbf{Schéma.} Complète la chaîne de conversion ci-dessous en nommant les 4 blocs et les tensions entre eux :
\begin{center}
\begin{tikzpicture}[font=\small, >=Stealth]
\node[draw, thick, rounded corners, fill=popblueL, minimum width=2cm] (a) {SECTEUR};
\node[draw, thick, rounded corners, fill=poporangeL, minimum width=2.5cm, right=1.5cm of a] (b) {?};
\node[draw, thick, rounded corners, fill=poporangeL, minimum width=2.5cm, right=1.5cm of b] (c) {?};
\node[draw, thick, rounded corners, fill=poporangeL, minimum width=2.5cm, right=1.5cm of c] (d) {?};
\node[draw, thick, rounded corners, fill=popgreenL, minimum width=2cm, right=1.5cm of d] (e) {SORTIE};
\draw[->, thick] (a) -- (b) node[midway, above] {\SI{220}{V} $\sim$};
\draw[->, thick] (b) -- (c) node[midway, above] {?};
\draw[->, thick] (c) -- (d) node[midway, above] {?};
\draw[->, thick] (d) -- (e) node[midway, above] {?};
\end{tikzpicture}
\end{center}
\item \textbf{Choix d'adaptateur.} Ton ordinateur portable nécessite \SI{19}{V} --- et \SI{3,42}{A}. Tu as deux adaptateurs de récupération :
\begin{itemize}
\item Adaptateur A : OUTPUT \SI{19}{V} --- \SI{2,1}{A}
\item Adaptateur B : OUTPUT \SI{19}{V} --- \SI{4,7}{A}
\end{itemize}
Lequel peux-tu utiliser sans risque ? Justifie.
\end{enumerate}
\end{exercice}

\begin{corrige}[Exercice : À toi de jouer]
\begin{enumerate}
\item Parce que la prise délivre \SI{220}{V} alternatifs, une tension beaucoup trop élevée et de mauvaise nature : la radio serait détruite instantanément et il y aurait un danger d'incendie et d'électrocution. Il faut un adaptateur qui fournisse \SI{9}{V} continus.
\item « L'adaptateur reçoit une tension \textbf{alternative} de valeur \textbf{\SI{220}{V}} et fournit une tension \textbf{continue} de valeur \textbf{\SI{12}{V}}. »
\item C'est faux : le téléphone ne reçoit jamais \SI{220}{V}. Le chargeur (adaptateur) convertit les \SI{220}{V} alternatifs du secteur en \SI{5}{V} continus, et c'est cette faible tension continue qui alimente le téléphone.
\item Chaîne complète :
\begin{itemize}
\item Bloc 1 : \textbf{Transformateur} (ou Transformation) — Tension entre 1 et 2 : \textbf{\SI{9}{V} $\sim$ (ou valeur basse alternative)}.
\item Bloc 2 : \textbf{Redresseur / Pont de diodes} — Tension entre 2 et 3 : \textbf{Unilatérale pulsée}.
\item Bloc 3 : \textbf{Filtre / Lissage (Condensateur)} — Tension entre 3 et 4 (Sortie) : \textbf{\SI{9}{V} --- (Continue)}.
\end{itemize}
\item \textbf{Adaptateur B.} L'adaptateur A ne peut fournir que \SI{2,1}{A} max, or l'ordinateur demande \SI{3,42}{A}. L'adaptateur A surchaufferait, sa tension s'effondrerait, et il risquerait l'incendie. L'adaptateur B fournit jusqu'à \SI{4,7}{A} > \SI{3,42}{A} : il fonctionne en toute sécurité (il ne donnera que le courant demandé par l'ordi).
\end{enumerate}
\end{corrige}

\begin{aretenir}[L'essentiel de la leçon — L'adaptateur secteur]
\begin{itemize}
\item Le secteur ENEO délivre une tension \textbf{alternative} de \SI{220}{V} (\SI{50}{Hz}) ; les petits appareils fonctionnent sous une tension \textbf{continue} de quelques volts.
\item Brancher directement un appareil basse tension sur le secteur est \textbf{impossible et dangereux} (destruction, incendie, électrocution).
\item L'\cle{adaptateur secteur} convertit la tension alternative du secteur en une tension continue de faible valeur grâce à trois fonctions successives :
\begin{enumerate}
\item \textbf{Transformation} (transformateur abaisseur) : \SI{220}{V} $\sim$ $\rightarrow$ quelques volts $\sim$ (isolement de sécurité).
\item \textbf{Redressement} (pont de diodes / Graetz) : alternative $\rightarrow$ unilatérale pulsée.
\item \textbf{Lissage / Régulation} (condensateur + régulateur) : pulsée $\rightarrow$ \textbf{continue stabilisée}.
\end{enumerate}
\item La plaque signalétique se lit en deux lignes : \textbf{INPUT} (entrée : ce que l'adaptateur reçoit du secteur) et \textbf{OUTPUT} (sortie : ce qu'il fournit à l'appareil), avec les symboles $\sim$ (alternatif) et --- $\cdot\cdot\cdot$ (continu).
\item \textbf{Sécurité} : Vérifier tension, intensité, polarité ; ne pas couvrir ; éviter les contrefaçons ; débrancher après usage.
\end{itemize}
\end{aretenir}

\coursec{Première fonction : la transformation (abaissement de tension)}

Dans la section précédente, nous avons vu que nos petits appareils (téléphone, radio, lampe rechargeable) ont besoin d'une tension \emph{continue} et \emph{faible}, alors que le secteur ENEO fournit une tension \emph{alternative} de \SI{220}{V}. L'adaptateur secteur doit donc accomplir deux tâches : d'abord \cle{abaisser} la tension, ensuite la \cle{redresser}. Dans cette section, nous étudions la première tâche : comment passer de \SI{220}{V} alternatifs à une tension alternative beaucoup plus faible, par exemple \SI{12}{V} ? L'appareil qui réalise cette prouesse s'appelle le \cle{transformateur}.

\courssub{Le problème posé}

\begin{exemplebox}[Une situation de la vie courante]
Le chargeur de ton téléphone porte une petite étiquette : « Entrée : 220 V $\sim$ ; Sortie : 5 V --- ». Le secteur délivre \SI{220}{V} : si l'on branchait directement la batterie du téléphone sur la prise, elle serait détruite instantanément, et il y aurait un grave danger d'incendie. Il faut donc \textbf{diviser la tension par plus de 40}, sans gaspiller d'énergie, et sans contact direct entre le circuit dangereux (le secteur) et le circuit que tu touches (ton téléphone). Comment faire ?
\end{exemplebox}

On pourrait penser à utiliser un simple résistor pour « absorber » la tension excédentaire. Mais un résistor traversé par un courant chauffe : il gaspille l'énergie en chaleur, et il ne protège pas l'utilisateur car le circuit reste électriquement relié au secteur. Il faut une solution plus astucieuse, qui utilise deux phénomènes remarquables : un courant électrique peut créer un champ magnétique, et un champ magnétique \emph{variable} peut, à son tour, créer une tension dans une bobine voisine, \textbf{sans aucun contact électrique}. C'est le principe du transformateur.

\courssub{Le transformateur : description et constitution}

\begin{definition}[Le transformateur]
Un \cle{transformateur} est un appareil qui \textbf{modifie la valeur d'une tension alternative} grâce à \textbf{deux bobines} (enroulements de fil de cuivre isolé) \textbf{couplées par un noyau de fer doux}, sans contact électrique entre elles.
\begin{itemize}
\item La bobine reliée à la source (le secteur) s'appelle le \cle{primaire} ; elle possède $N_1$ spires et reçoit la tension $U_1$.
\item La bobine reliée à l'appareil à alimenter s'appelle le \cle{secondaire} ; elle possède $N_2$ spires et délivre la tension $U_2$.
\item Le \cle{noyau de fer doux} est un cadre de fer qui canalise le champ magnétique d'une bobine vers l'autre.
\end{itemize}
\end{definition}

\begin{popfigure}
\begin{center}
\begin{tikzpicture}[scale=1.0, every node/.style={font=\small}]
% Noyau de fer doux (cadre)
\draw[line width=5pt, popmuted] (0.4,0) rectangle (5.6,3.4);
\draw[line width=5pt, white] (0.9,0.5) rectangle (5.1,2.9);
\draw[line width=1pt, popdark] (0.4,0) rectangle (5.6,3.4);
\draw[line width=1pt, popdark] (0.9,0.5) rectangle (5.1,2.9);
% Bobine primaire (colonne gauche)
\foreach \y in {0.75,1.0,1.25,1.5,1.75,2.0,2.25,2.5,2.65}{
  \draw[popblue, line width=1.4pt] (0.15,\y) arc[start angle=-90, end angle=270, x radius=0.12, y radius=0.125];
}
\draw[popblue, line width=1.4pt] (0.15,0.75) -- (-0.9,0.75);
\draw[popblue, line width=1.4pt] (0.15,2.65) -- (-0.9,2.65);
% Source 220 V
\draw[popdark] (-0.9,0.75) -- (-1.6,0.75);
\draw[popdark] (-0.9,2.65) -- (-1.6,2.65);
\draw[popdark, line width=1pt] (-1.6,1.7) circle (0.45);
\node[popdark] at (-1.6,1.7) {$\sim$};
\node[left, popblue] at (-2.15,1.7) {$U_1 = \SI{220}{V}$};
\node[popblue, align=center] at (1.6,-0.55) {Primaire : $N_1$ spires};
% Bobine secondaire (colonne droite) : moins de spires
\foreach \y in {1.15,1.4,1.65,1.9,2.15}{
  \draw[poporange, line width=1.4pt] (5.85,\y) arc[start angle=-90, end angle=270, x radius=0.12, y radius=0.125];
}
\draw[poporange, line width=1.4pt] (5.85,1.15) -- (6.8,1.15);
\draw[poporange, line width=1.4pt] (5.85,2.15) -- (6.8,2.15);
\draw[popdark] (6.8,1.15) -- (7.5,1.15);
\draw[popdark] (6.8,2.15) -- (7.5,2.15);
\draw[popdark, line width=1pt] (7.5,1.65) circle (0.4);
\node[popdark] at (7.5,1.65) {V};
\node[right, poporange] at (8.0,1.65) {$U_2 = \SI{12}{V}$};
\node[poporange, align=center] at (4.4,-0.55) {Secondaire : $N_2$ spires};
% Noyau label
\node[popmuted, fill=white, inner sep=1pt] at (3,3.4) {\footnotesize Noyau de fer doux};
% Flèche champ magnétique
\draw[-{Stealth}, popgreen, line width=1.2pt] (2.2,1.7) -- (3.8,1.7);
\node[popgreen, below] at (3,1.55) {\footnotesize champ magnétique variable};
\end{tikzpicture}
\end{center}
\legende{Schéma de principe du transformateur abaisseur. Le primaire ($N_1$ spires, bleu) est branché sur le secteur \SI{220}{V} $\sim$ ; le secondaire ($N_2$ spires, orange) délivre une tension plus faible, ici \SI{12}{V} $\sim$. Les deux bobines sont enroulées sur le même noyau de fer doux mais \textbf{ne se touchent pas} : l'énergie passe par le champ magnétique.}
\end{popfigure}

\begin{attention}[Aucun contact électrique entre les bobines]
Le fil de cuivre des bobines est recouvert d'un vernis isolant, et les deux bobines sont séparées. Il n'y a donc \textbf{aucun passage direct de courant} entre le primaire et le secondaire. C'est une sécurité essentielle : le circuit de ton téléphone n'est jamais relié électriquement au secteur \SI{220}{V}. Ne dis jamais que « le courant passe du primaire au secondaire » : c'est l'\emph{énergie} qui passe, transportée par le champ magnétique.
\end{attention}

\courssub{Comment ça marche ? Explication qualitative}

Le fonctionnement du transformateur repose sur deux phénomènes que tu as déjà rencontrés :

\begin{enumerate}
\item \textbf{Un courant crée un champ magnétique.} Une bobine parcourue par un courant se comporte comme un aimant. Si le courant est \emph{alternatif}, il change de sens 50 fois par seconde (fréquence du réseau ENEO : $f = \SI{50}{Hz}$) : le champ magnétique créé par le primaire est donc \cle{variable} — il grandit, s'annule, s'inverse, en permanence.
\item \textbf{Un champ magnétique variable crée une tension.} Lorsqu'une bobine est plongée dans un champ magnétique \emph{variable}, il apparaît à ses bornes une tension alternative : c'est le phénomène d'\cle{induction électromagnétique}. Le noyau de fer doux guide ce champ variable jusqu'au secondaire, où une tension $U_2$ apparaît.
\end{enumerate}

\begin{aretenir}[Le mécanisme en une phrase]
Le primaire, parcouru par un courant \textbf{alternatif}, crée un \textbf{champ magnétique variable} dans le noyau de fer ; ce champ variable \textbf{induit} une tension \textbf{alternative} aux bornes du secondaire. La valeur de cette tension dépend du nombre de spires de chaque bobine.
\end{aretenir}

\courssub{Expérience de classe : la tension de sortie dépend du nombre de spires}

\begin{experience}[Influence du nombre de spires sur la tension secondaire]
\textbf{Matériel :} un transformateur démontable (noyau en U avec barreau amovible), deux bobines de nombres de spires différents, un générateur basse tension réglé sur \SI{12}{V} $\sim$, deux voltmètres en mode alternatif (AC), des fils de connexion.

\textbf{Protocole :}
\begin{enumerate}
\item On ferme le noyau de fer avec son barreau, puis on branche le primaire ($N_1 = 600$ spires) sur le générateur réglé à $U_1 = \SI{12}{V}$ $\sim$.
\item On mesure la tension $U_2$ aux bornes du secondaire pour différentes valeurs de $N_2$.
\item On consigne les résultats dans un tableau et on calcule les rapports $\dfrac{U_2}{U_1}$ et $\dfrac{N_2}{N_1}$.
\end{enumerate}

\textbf{Observations (résultats typiques) :}

\begin{center}
\begin{tabularx}{\linewidth}{|c|X|X|X|X|}
\hline
\rowcolor{popblueL} Expérience & $N_1$ (spires) & $N_2$ (spires) & $U_2$ mesurée (V) & $\dfrac{N_2}{N_1}$ \\
\hline
1 & 600 & 1200 & 24 & 2 \\
\hline
2 & 600 & 600 & 12 & 1 \\
\hline
3 & 600 & 300 & 6 & 0,5 \\
\hline
4 & 600 & 150 & 3 & 0,25 \\
\hline
\end{tabularx}
\end{center}

\textbf{Interprétation :} dans chaque ligne, le rapport des tensions est égal au rapport des nombres de spires. Par exemple, à l'expérience 3 : $\dfrac{U_2}{U_1} = \dfrac{6}{12} = 0,5$ et $\dfrac{N_2}{N_1} = \dfrac{300}{600} = 0,5$. On vérifie de même à l'expérience 1 : $\dfrac{U_2}{U_1} = \dfrac{24}{12} = 2 = \dfrac{N_2}{N_1}$.

\textbf{Conclusion :} la tension de sortie d'un transformateur dépend du rapport des nombres de spires. Si $N_2 < N_1$, la tension diminue : le transformateur est \cle{abaisseur}. Si $N_2 > N_1$, elle augmente : il est \cle{élévateur}.
\end{experience}

\begin{propriete}[Relation fondamentale du transformateur — admise, établie expérimentalement]
Pour un transformateur parfait, le rapport des tensions est égal au rapport des nombres de spires :
\[ \frac{U_2}{U_1} = \frac{N_2}{N_1} \]
avec $U_1$ et $U_2$ les tensions (en volts) au primaire et au secondaire, $N_1$ et $N_2$ les nombres de spires correspondants. Cette relation est \textbf{admise} en classe de 3ème : elle se vérifie expérimentalement, sa démonstration n'est pas exigible au BEPC. En revanche, savoir l'\emph{utiliser} dans un calcul simple est exigible.
\end{propriete}

\begin{exemplebox}[Transformateur abaisseur et élévateur]
\begin{itemize}
\item \textbf{Abaisseur :} le secondaire a \emph{moins} de spires que le primaire ($N_2 < N_1$), donc $U_2 < U_1$. C'est le cas dans l'adaptateur secteur : on passe de \SI{220}{V} à \SI{12}{V}.
\item \textbf{Élévateur :} le secondaire a \emph{plus} de spires ($N_2 > N_1$), donc $U_2 > U_1$. On l'utilise par exemple pour élever la tension avant le transport de l'électricité sur les lignes à haute tension.
\end{itemize}
\end{exemplebox}

\begin{exoresolu}[Calcul de la tension de sortie d'un adaptateur]
Énoncé. Le transformateur d'un adaptateur secteur possède $N_1 = 1100$ spires au primaire et $N_2 = 60$ spires au secondaire. Il est branché sur le secteur $U_1 = \SI{220}{V}$ $\sim$. Calculer la tension $U_2$ disponible au secondaire. Ce transformateur est-il abaisseur ou élévateur ?
\tcblower
Solution détaillée.
\begin{enumerate}
\item On écrit la relation du transformateur : $\dfrac{U_2}{U_1} = \dfrac{N_2}{N_1}$.
\item On isole l'inconnue $U_2$ : $U_2 = U_1 \times \dfrac{N_2}{N_1}$.
\item Application numérique :
\begin{align*}
U_2 &= 220 \times \frac{60}{1100} \\
&= 220 \times 0,0545 \\
&\approx \SI{12}{V}
\end{align*}
\item Comparaison : $N_2 = 60 < N_1 = 1100$, donc $U_2 < U_1$ : le transformateur est bien un \textbf{abaisseur}.
\end{enumerate}
Conclusion : le secondaire délivre une tension alternative d'environ \SI{12}{V}, environ 18 fois plus faible que le secteur. C'est exactement le rôle du premier étage de l'adaptateur secteur.
\end{exoresolu}

\courssub{Le transformateur ne fonctionne qu'en courant alternatif}

\begin{experience}[Et si l'on alimente le primaire en courant continu ?]
\textbf{Protocole :} on reprend le montage précédent, mais on remplace le générateur de tension alternative par un générateur de tension \emph{continue} de \SI{12}{V}. On mesure la tension au secondaire.

\textbf{Observation :} le voltmètre branché au secondaire indique \textbf{zéro volt} (à part un bref « saut » de l'aiguille à la fermeture et à l'ouverture du circuit). Au primaire, le courant continu est important et la bobine chauffe dangereusement.

\textbf{Interprétation :} avec un courant continu, le champ magnétique créé par le primaire est \emph{constant}. Or seul un champ magnétique \emph{variable} peut induire une tension dans le secondaire. Champ constant $\Rightarrow$ aucune tension induite.

\textbf{Conclusion :} le transformateur \textbf{ne fonctionne qu'en courant alternatif}.
\end{experience}

\begin{attention}[Erreur fréquente au BEPC]
Beaucoup d'élèves écrivent : « le transformateur transforme le courant continu en courant alternatif » ou « le transformateur fonctionne avec n'importe quel courant ». C'est \textbf{faux} ! Retiens bien :
\begin{itemize}
\item Le transformateur ne fonctionne \textbf{qu'en courant alternatif} : alimenté en continu, il ne délivre \textbf{aucune tension} au secondaire et risque de chauffer et de griller.
\item Le transformateur ne change \textbf{pas la nature} du courant : il reçoit de l'alternatif et fournit de l'alternatif (de valeur différente). C'est le \emph{redresseur}, étudié dans la section suivante, qui convertira l'alternatif en continu.
\end{itemize}
\end{attention}

\courssub{Bilan : le rôle du transformateur dans l'adaptateur secteur}

Dans l'adaptateur secteur, le transformateur est le \textbf{premier étage} de la chaîne de conversion. Branché sur la prise, il abaisse la tension du secteur de \SI{220}{V} $\sim$ à une tension alternative faible, par exemple \SI{12}{V} $\sim$, tout en isolant électriquement ton appareil du réseau dangereux. Mais la tension obtenue est encore \emph{alternative} : elle change de sens 50 fois par seconde, ce qui ne convient pas à une batterie ou à un circuit électronique. Il restera donc à la \emph{redresser} : ce sera l'objet de la section suivante.

\begin{savaistu}[Les transformateurs sur les poteaux ENEO]
Tu as sûrement remarqué, sur certains poteaux électriques au Cameroun, de gros boîtiers métalliques gris avec des ailettes de refroidissement : ce sont des \textbf{transformateurs de distribution}. L'électricité est transportée dans les quartiers sous une tension de \SI{15000}{V} (pour limiter les pertes en chaleur dans les câbles), puis ces transformateurs l'abaissent à \SI{220}{V} pour alimenter les habitations. À l'inverse, à la sortie des centrales hydroélectriques (Songloulou, Edéa, Memve'ele), d'énormes transformateurs \emph{élévateurs} portent la tension à \SI{90000}{V} ou \SI{225000}{V} pour le transport sur longue distance. Le transformateur est donc partout autour de toi, du minuscule chargeur de téléphone au poste de transformation de ton quartier !
\end{savaistu}

\begin{exercice}[À toi de jouer]
Un transformateur possède $N_1 = 800$ spires au primaire et $N_2 = 40$ spires au secondaire. On le branche sur le secteur $U_1 = \SI{220}{V}$ $\sim$.
\begin{enumerate}
\item Ce transformateur est-il abaisseur ou élévateur ? Justifier.
\item Calculer la tension $U_2$ au secondaire.
\item On branche par erreur le primaire sur une batterie de \SI{12}{V} continu. Que vaut la tension au secondaire ? Pourquoi ?
\end{enumerate}
\end{exercice}

\begin{corrige}[Exercice]
\begin{enumerate}
\item $N_2 = 40 < N_1 = 800$ : le secondaire a moins de spires que le primaire, donc $U_2 < U_1$ : c'est un transformateur \textbf{abaisseur}.
\item $U_2 = U_1 \times \dfrac{N_2}{N_1} = 220 \times \dfrac{40}{800} = 220 \times 0,05 = \SI{11}{V}$.
\item La tension au secondaire est \textbf{nulle} : la batterie fournit un courant \emph{continu}, le champ magnétique créé est constant, et seul un champ magnétique \emph{variable} peut induire une tension au secondaire. Le transformateur ne fonctionne qu'en courant alternatif.
\end{enumerate}
\end{corrige}

\coursec{Deuxième fonction : le redressement (alternatif → continu)}

Dans la section précédente, nous avons vu que le \cle{transformateur} abaisse la tension du secteur : de \SI{220}{V} alternatif, on obtient par exemple \SI{12}{V} alternatif. La tension est devenue sans danger. Mais un problème subsiste : ton téléphone portable, ta radio ou ta calculatrice ne fonctionnent pas avec une tension \emph{alternative}. Ils exigent une tension \emph{continue}, comme celle d'une pile. C'est ici qu'intervient la deuxième fonction de l'adaptateur : le \cle{redressement}.

\courssub{Le problème posé : une tension faible, mais encore alternative}

\begin{exemplebox}[Situation de la vie courante]
Tu branches ton téléphone sur son chargeur. À l'intérieur, le transformateur vient d'abaisser la tension à quelques volts. Mais la batterie du téléphone se comporte comme une pile : elle possède une borne $+$ et une borne $-$, et elle doit être rechargée par un courant qui circule \textbf{toujours dans le même sens}. Or le courant alternatif change de sens 50 fois par seconde (fréquence $f = \SI{50}{Hz}$ au Cameroun). Brancher directement la sortie du transformateur sur la batterie reviendrait à la charger puis la décharger 50 fois par seconde : elle ne se chargerait jamais !
\end{exemplebox}

\begin{experience}[Observation de la tension à la sortie du transformateur]
\textbf{Matériel :} un transformateur 220 V / 12 V, un oscilloscope (ou un voltmètre en position AC), des fils de connexion.

\textbf{Protocole :}
\begin{enumerate}
\item Brancher le primaire du transformateur sur le secteur \SI{220}{V}.
\item Relier le secondaire aux bornes de l'oscilloscope.
\item Observer l'oscillogramme, puis mesurer la tension avec un voltmètre en position AC, puis en position DC.
\end{enumerate}

\textbf{Observations :}
\begin{itemize}
\item L'oscillogramme est une \cle{sinusoïde} : la courbe monte, redescend, passe par zéro, devient négative, puis remonte, de façon régulière.
\item Le voltmètre en position AC indique environ \SI{12}{V}.
\item Le voltmètre en position DC indique \SI{0}{V} : la \textbf{valeur moyenne} de la tension est nulle, car les alternances positives et négatives se compensent exactement.
\end{itemize}

\textbf{Interprétation :} la tension alternative change de signe périodiquement. Sa moyenne est nulle : elle ne peut donc pas alimenter un appareil qui exige une tension de signe constant.
\end{experience}

\begin{popfigure}
\begin{center}
\begin{tikzpicture}
\begin{axis}[
width=0.85\linewidth, height=5.2cm,
xlabel={$t$ (ms)}, ylabel={$u$ (V)},
xmin=0, xmax=40, ymin=-15, ymax=15,
xtick={0,10,20,30,40}, ytick={-12,0,12},
grid=both, grid style={popline!50},
axis lines=middle,
]
\addplot[popcurve, domain=0:40, samples=300] {12*sin(360*x/20)};
\end{axis}
\end{tikzpicture}
\end{center}
\legende{Tension sinusoïdale à la sortie du transformateur : elle est alternativement positive et négative, sa valeur moyenne est nulle (période $T = \SI{20}{ms}$, soit $f = \SI{50}{Hz}$).}
\end{popfigure}

\courssub{La diode : un composant à sens unique}

Pour transformer une tension alternative en tension de signe constant, il faut un composant capable de « trier » le sens du courant. Ce composant existe : c'est la \cle{diode}.

\begin{definition}[La diode]
Une \cle{diode} est un composant électrique à deux bornes (l'\textbf{anode} A et la \textbf{cathode} K) qui ne laisse passer le courant que dans \textbf{un seul sens}, appelé \textbf{sens passant} : de l'anode vers la cathode. Dans l'autre sens (sens bloqué), elle se comporte comme un interrupteur ouvert.
\end{definition}

\begin{popfigure}
\begin{center}
\begin{tikzpicture}[line cap=round, >=Stealth]
% Symbole de la diode
\draw[thick, popdark] (-3,0) -- (-1,0);
\draw[very thick, popblue, fill=popblueL] (-1,-0.5) -- (-1,0.5) -- (0.2,0) -- cycle;
\draw[very thick, popblue] (0.2,-0.5) -- (0.2,0.5);
\draw[thick, popdark] (0.2,0) -- (2.2,0);
\node[below, popdark] at (-2.4,-0.1) {\textbf{A}};
\node[below, popdark] at (1.6,-0.1) {\textbf{K}};
\node[above, popdark] at (-0.4,0.6) {anode};
\node[above, popdark] at (0.9,0.6) {cathode};
% Sens du courant autorisé
\draw[->, very thick, popgreen] (-2.6,1.4) -- (1.8,1.4) node[midway, above, popgreen!60!black] {sens passant : courant autorisé};
% Sens interdit
\draw[->, very thick, poppink] (1.8,-1.4) -- (-2.6,-1.4) node[midway, below, poppink!70!black] {sens bloqué : courant interdit};
\node[popmuted] at (5.2,0) {\textbf{= valve à sens unique}};
\end{tikzpicture}
\end{center}
\legende{Symbole de la diode : le triangle indique le sens passant (de A vers K), la barre représente la cathode. La diode agit comme un clapet anti-retour.}
\end{popfigure}

\begin{savaistu}[Une image pour comprendre]
La diode fonctionne comme la valve d'une chambre à air de moto-taxi : l'air peut entrer pour gonfler le pneu, mais il ne peut pas ressortir par la valve. De même, le courant peut traverser la diode dans le sens passant, mais il est bloqué dans l'autre sens.
\end{savaistu}

\begin{experience}[Mise en évidence du sens passant]
\textbf{Matériel :} un générateur de tension alternative (GBF), une diode, une DEL (diode électroluminescente), un résistor de protection, des fils.

\textbf{Protocole :}
\begin{enumerate}
\item Réaliser un circuit série : GBF, diode, résistor, DEL.
\item Observer la DEL, puis inverser le sens de la diode et observer à nouveau.
\end{enumerate}

\textbf{Observations :}
\begin{itemize}
\item La DEL ne brille pas en permanence : elle \textbf{clignote}, elle ne s'allume que par intermittence.
\item En inversant la diode, la DEL clignote toujours, mais pendant les alternances opposées.
\end{itemize}

\textbf{Interprétation :} pendant une alternance sur deux, le courant est dans le sens passant de la diode : il circule et la DEL brille. Pendant l'autre alternance, la diode bloque le courant : la DEL s'éteint. La diode a donc \textbf{supprimé une alternance sur deux}.
\end{experience}

\courssub{Le redressement : rendre la tension unidirectionnelle}

\begin{definition}[Le redressement]
Le \cle{redressement} est la conversion d'une tension \textbf{alternative} (qui change de signe) en une tension \textbf{unidirectionnelle} (qui garde toujours le même signe), à l'aide de \textbf{diodes}. On dit que la tension est \emph{redressée}.
\end{definition}

Avec une seule diode, on supprime les alternances négatives : c'est le \textbf{redressement simple alternance}. Mais on perd la moitié de l'énergie. Pour être plus efficace, on utilise un \textbf{pont de diodes} : il \emph{inverse} les alternances négatives au lieu de les supprimer. Toutes les alternances deviennent positives : c'est le \textbf{redressement double alternance}.

\begin{popfigure}
\begin{center}
\begin{tikzpicture}
\begin{axis}[
width=0.85\linewidth, height=5.2cm,
xlabel={$t$ (ms)}, ylabel={$u$ (V)},
xmin=0, xmax=40, ymin=-2, ymax=15,
xtick={0,10,20,30,40}, ytick={0,12},
grid=both, grid style={popline!50},
axis lines=middle,
]
\addplot[popcurve, domain=0:40, samples=300] {abs(12*sin(360*x/20))};
\end{axis}
\end{tikzpicture}
\end{center}
\legende{Tension redressée double alternance (pont de diodes) : les alternances négatives ont été inversées. La tension reste toujours positive, mais elle varie encore : on dit qu'elle est \emph{ondulée}.}
\end{popfigure}

\begin{methode}[Interpréter un oscillogramme de tension redressée]
\begin{enumerate}
\item Repérer l'axe horizontal des temps (niveau zéro de la tension).
\item Observer la position de la courbe par rapport à cet axe.
\item Si la courbe reste \textbf{toujours du même côté} de l'axe (toujours au-dessus ou toujours en dessous), la tension est \textbf{unidirectionnelle} : elle a été redressée.
\item Si la courbe \textbf{traverse} l'axe et change de signe, la tension est encore \textbf{alternative} : le redressement n'a pas eu lieu (ou la diode est défectueuse).
\end{enumerate}
\end{methode}

\begin{attention}[Erreur fréquente au BEPC]
Ne dis jamais que la tension redressée est « parfaitement continue ». Après redressement, la tension est bien \textbf{unidirectionnelle} (elle ne change plus de signe), mais elle est encore \textbf{ondulée} : sa valeur monte et descend entre 0 et sa valeur maximale. Une tension continue parfaite, comme celle d'une pile, garde une valeur \emph{constante}. C'est le rôle du filtrage de s'en approcher.
\end{attention}

\courssub{Complément : le filtrage (hors programme strict)}

\begin{savaistu}[Le condensateur de filtrage]
Pour rendre la tension redressée presque continue, les adaptateurs utilisent un \cle{condensateur} branché après le pont de diodes. Ce composant se charge quand la tension monte et se décharge quand elle descend, comme un petit réservoir : il « \textbf{lisse} » la tension. On obtient alors une tension presque constante, très proche de celle d'une pile. Cette étude dépasse le programme de 3ème : retiens simplement que le filtrage \emph{améliore} la tension redressée.
\end{savaistu}

\begin{savaistu}[Le pont de Graetz]
Le montage de 4 diodes utilisé pour redresser les deux alternances s'appelle le \textbf{pont de Graetz}, du nom du physicien allemand Leo Graetz qui le décrivit en 1897. Grâce à lui, les deux alternances sont utilisées : la conversion est deux fois plus efficace qu'avec une seule diode. Ce pont est présent dans quasiment tous les chargeurs de téléphone et adaptateurs que tu utilises chaque jour au Cameroun.
\end{savaistu}

\begin{exoresolu}[Identifier une tension redressée]
Un élève branche un oscilloscope à la sortie d'un adaptateur ouvert, juste après le pont de diodes. Il observe une courbe qui reste toujours au-dessus de l'axe des temps, mais qui monte et descend régulièrement entre \SI{0}{V} et \SI{12}{V}. Un camarade affirme : « La tension est continue, l'adaptateur est en panne si elle varie. » Qui a raison ? Justifier.
\tcblower
\textbf{Solution :} L'élève observateur a raison, son camarade se trompe.
\begin{itemize}
\item La courbe reste \textbf{toujours du même côté} de l'axe des temps : la tension ne change plus de signe. Elle est donc bien \textbf{redressée} (unidirectionnelle) : le pont de diodes fonctionne correctement.
\item La tension varie encore entre 0 et \SI{12}{V} : elle est \textbf{ondulée}, ce qui est normal \emph{avant} le filtrage par le condensateur.
\item Une tension vraiment continue (comme celle d'une pile) garderait une valeur constante. Ici, l'adaptateur n'est pas en panne : l'observation est faite avant l'étage de filtrage.
\end{itemize}
\textbf{Conclusion :} une tension redressée est unidirectionnelle mais pas nécessairement constante.
\end{exoresolu}

\begin{aretenir}[L'essentiel sur le redressement]
\begin{itemize}
\item La \cle{diode} ne laisse passer le courant que dans un seul sens (sens passant, de l'anode vers la cathode).
\item Le \cle{redressement} convertit une tension alternative en tension \textbf{unidirectionnelle} (de signe constant), grâce à des diodes.
\item Le \textbf{pont de diodes} (pont de Graetz, 4 diodes) redresse les deux alternances : la tension obtenue est toujours positive mais \textbf{ondulée}.
\item Le \textbf{condensateur de filtrage} lisse ensuite cette tension pour la rendre presque continue (complément hors programme).
\end{itemize}
\end{aretenir}

\begin{exercice}[À toi de jouer]
On visualise à l'oscilloscope deux tensions : la tension A est une sinusoïde qui passe alternativement au-dessus et en dessous de l'axe des temps ; la tension B reste toujours au-dessus de l'axe mais oscille entre 0 et \SI{9}{V}.
\begin{enumerate}
\item Laquelle de ces deux tensions est alternative ? Justifier.
\item Laquelle est redressée ? Justifier.
\item La tension B est-elle continue ? Expliquer.
\end{enumerate}
\end{exercice}

\begin{corrige}[Exercice]
\begin{enumerate}
\item La tension A est alternative : elle change de signe périodiquement (elle traverse l'axe des temps).
\item La tension B est redressée : elle garde toujours le même signe (toujours positive), elle est unidirectionnelle.
\item Non, la tension B n'est pas continue : sa valeur n'est pas constante, elle oscille entre 0 et \SI{9}{V}. Elle est unidirectionnelle mais ondulée. Il faudrait un filtrage pour la rendre presque continue.
\end{enumerate}
\end{corrige}

\coursec{Schéma de principe et chaîne de conversion de l'adaptateur}

Nous venons de voir, dans les sections précédentes, les deux grandes fonctions de l'adaptateur secteur : le \cle{transformateur} abaisse la tension du secteur de 220 V à une basse tension alternative, puis le \cle{redresseur} à diodes transforme cette tension alternative en tension continue. Il est temps de rassembler tout cela dans un seul schéma : comment ces fonctions s'enchaînent-elles à l'intérieur du petit boîtier noir branché sur la prise de ta radio ou de ton téléphone ? C'est ce que nous allons découvrir, étape par étape, en suivant le trajet de l'électricité depuis la prise murale jusqu'à l'appareil.

\courssub{Synthèse : deux fonctions successives}

\begin{propriete}[Structure générale d'un adaptateur secteur]
Un \cle{adaptateur secteur} réalise deux fonctions successives :
\begin{enumerate}
\item la \cle{transformation} : un transformateur abaisseur ramène la tension du secteur (220 V alternatif) à une basse tension alternative (par exemple 9 V ou 12 V) ;
\item le \cle{redressement} : un pont de diodes convertit cette basse tension alternative en tension continue.
\end{enumerate}
Un étage de \cle{filtrage} (condensateur) peut s'ajouter pour rendre la tension de sortie plus stable, presque parfaitement continue.
\end{propriete}

L'ordre de ces fonctions n'est pas un détail : il est imposé par la physique des composants. Le transformateur ne fonctionne \emph{qu'en courant alternatif} (il utilise les variations du courant pour créer une tension induite au secondaire). Il faut donc d'abord abaisser la tension pendant qu'elle est encore alternative, et seulement ensuite la redresser.

\begin{attention}[Ne pas inverser l'ordre des fonctions]
Erreur fréquente au BEPC : écrire que l'adaptateur « redresse puis transforme ». C'est impossible !
\begin{itemize}
\item Le transformateur a besoin d'une tension \textbf{alternative} pour fonctionner : si l'on redressait d'abord le 220 V, le transformateur ne pourrait plus rien faire (et serait dangereusement traversé par un courant continu).
\item On \textbf{abaisse d'abord} la tension (transformateur), \textbf{puis on redresse} (pont de diodes).
\end{itemize}
Retiens l'ordre : \textbf{T}ransformation puis \textbf{R}edressement — « \textbf{T} avant \textbf{R} », comme dans l'alphabet.
\end{attention}

\courssub{La chaîne de conversion de l'adaptateur}

Suivons le voyage de l'électricité, de la prise ENEO jusqu'à ta radio. À chaque étape, la tension change de forme : on peut la visualiser avec un oscilloscope.

\begin{popfigure}
\begin{center}
\begin{tikzpicture}[font=\small, >=Stealth]
% Blocs
\draw[fill=popblueL, draw=popblue, thick, rounded corners] (0,0) rectangle (2.2,1.4);
\node[align=center] at (1.1,0.7) {\textbf{Secteur}\\ 220 V $\sim$};
\draw[fill=poporangeL, draw=poporange, thick, rounded corners] (3.6,0) rectangle (6.2,1.4);
\node[align=center] at (4.9,0.7) {\textbf{Transformateur}\\ abaisseur};
\draw[fill=popgreenL, draw=popgreen, thick, rounded corners] (7.6,0) rectangle (10.0,1.4);
\node[align=center] at (8.8,0.7) {\textbf{Redresseur}\\ pont de diodes};
\draw[fill=poppurpleL, draw=poppurple, thick, rounded corners] (11.4,0) rectangle (13.8,1.4);
\node[align=center] at (12.6,0.7) {\textbf{Appareil}\\ (radio...)};
% Flèches
\draw[->, very thick, popdark] (2.2,0.7) -- (3.6,0.7);
\draw[->, very thick, popdark] (6.2,0.7) -- (7.6,0.7);
\draw[->, very thick, popdark] (10.0,0.7) -- (11.4,0.7);
% Oscillogramme 1 : sinus 220 V
\begin{scope}[shift={(0,-2.6)}]
\draw[->] (0,0) -- (2.4,0); \draw[->] (0,-0.9) -- (0,1.0);
\draw[popblue, thick, domain=0:2.2, samples=80] plot (\x, {0.75*sin(360*\x/1.1)});
\node[below, align=center] at (1.1,-0.95) {220 V $\sim$\\ sinusoïdal};
\end{scope}
% Oscillogramme 2 : sinus 12 V
\begin{scope}[shift={(3.7,-2.6)}]
\draw[->] (0,0) -- (2.4,0); \draw[->] (0,-0.9) -- (0,1.0);
\draw[poporange, thick, domain=0:2.2, samples=80] plot (\x, {0.35*sin(360*\x/1.1)});
\draw[dashed, popmuted] (0,0.75) -- (2.2,0.75);
\node[below, align=center] at (1.1,-0.95) {12 V $\sim$\\ sinusoïdal (plus petit)};
\end{scope}
% Oscillogramme 3 : redressé
\begin{scope}[shift={(7.7,-2.6)}]
\draw[->] (0,0) -- (2.4,0); \draw[->] (0,-0.9) -- (0,1.0);
\draw[popgreen, thick, domain=0:2.2, samples=120] plot (\x, {0.35*abs(sin(360*\x/1.1))});
\node[below, align=center] at (1.1,-0.95) {12 V redressé\\ (une seule polarité)};
\end{scope}
% Oscillogramme 4 : continu
\begin{scope}[shift={(11.5,-2.6)}]
\draw[->] (0,0) -- (2.4,0); \draw[->] (0,-0.9) -- (0,1.0);
\draw[poppurple, very thick] (0,0.35) -- (2.2,0.35);
\node[below, align=center] at (1.1,-0.95) {$\approx$ 12 V continu\\ (filtré)};
\end{scope}
% Légendes étapes
\node[above, popblue] at (1.1,1.5) {\footnotesize étape 0};
\node[above, poporange] at (4.9,1.5) {\footnotesize étape 1 : transformation};
\node[above, popgreen] at (8.8,1.5) {\footnotesize étape 2 : redressement};
\node[above, poppurple] at (12.6,1.5) {\footnotesize étape 3 : utilisation};
\end{tikzpicture}
\end{center}
\legende{Schéma-bloc de la chaîne de conversion d'un adaptateur secteur. Sous chaque étape, l'allure de la tension observée à l'oscilloscope : sinusoïdale 220 V, sinusoïdale abaissée, redressée (les alternances négatives sont « relevées »), puis presque continue après filtrage.}
\end{popfigure}

Lisons cette chaîne dans le détail :

\begin{enumerate}
\item \textbf{Entrée : le secteur 220 V $\sim$.} La tension délivrée par ENEO est alternative sinusoïdale : elle effectue 50 périodes par seconde (fréquence 50 Hz), c'est-à-dire qu'elle change de sens 100 fois par seconde, et sa valeur efficace est 220 V. C'est une tension dangereuse, inutilisable directement par un appareil électronique.
\item \textbf{Étape 1 : le transformateur abaisseur.} Le secondaire du transformateur fournit une tension toujours alternative et sinusoïdale, de même fréquence (50 Hz), mais d'amplitude beaucoup plus faible : par exemple 12 V. Sur l'oscillogramme, la courbe a la même forme, elle est juste « écrasée » verticalement.
\item \textbf{Étape 2 : le pont de diodes redresse.} Les diodes ne laissent passer le courant que dans un sens. Les alternances négatives sont « relevées » : la tension ne change plus de polarité. Elle ondule encore (elle retombe à zéro entre deux bosses), mais elle reste toujours du même côté : c'est une tension \textbf{redressée}.
\item \textbf{Étape 3 (souvent présente) : le filtrage.} Un condensateur placé en sortie se charge pendant les « bosses » et se décharge entre elles : il lisse la tension. On obtient une tension presque constante, par exemple $\approx$ 12 V continu.
\item \textbf{Sortie : l'appareil.} La radio, le téléphone ou le petit ventilateur reçoit une basse tension continue, adaptée à son fonctionnement et sans danger pour l'utilisateur.
\end{enumerate}

\begin{aretenir}[L'essentiel de la chaîne]
\[\text{Secteur } 220~V \sim \;\longrightarrow\; \textbf{Transformateur} \;\longrightarrow\; 12~V \sim \;\longrightarrow\; \textbf{Pont de diodes} \;\longrightarrow\; 12~V \text{ redressé} \;\longrightarrow\; (\text{filtrage}) \;\longrightarrow\; \approx 12~V \text{ continu}\]
Le transformateur change la \textbf{valeur} de la tension (220 V $\rightarrow$ 12 V) ; le redresseur change sa \textbf{nature} (alternatif $\rightarrow$ continu).
\end{aretenir}

\courssub{Le circuit interne simplifié d'un adaptateur}

Ouvrons (en imagination !) le boîtier d'un adaptateur. On y trouve deux sous-ensembles faciles à reconnaître : le transformateur (un bloc avec deux bobinages sur une carcasse de fer doux) et quatre diodes montées en pont.

\begin{popfigure}
\begin{center}
\begin{tikzpicture}[font=\small, >=Stealth]
% Prise secteur
\draw[fill=popblueL, draw=popblue, thick] (0,-0.6) rectangle (1.4,0.6);
\node[align=center] at (0.7,0) {220 V\\ $\sim$};
% Fils vers primaire
\draw[thick] (1.4,0.4) -- (2.2,0.4);
\draw[thick] (1.4,-0.4) -- (2.2,-0.4);
% Transformateur : primaire (4 spires)
\draw[poporange, very thick] (2.2,0.4) -- (2.6,0.4);
\draw[poporange, very thick] (2.2,-0.4) -- (2.6,-0.4);
\foreach \y in {-0.4,-0.2,0,0.2} {\draw[poporange, thick] (2.6,\y) arc[start angle=-90, end angle=90, radius=0.1];}
% Noyau
\draw[very thick, popdark] (3.0,-0.7) -- (3.0,0.7);
\draw[very thick, popdark] (3.15,-0.7) -- (3.15,0.7);
% Secondaire (2 spires)
\draw[popgold, thick] (3.5,-0.2) arc[start angle=-90, end angle=90, radius=0.1];
\draw[popgold, thick] (3.5,0) arc[start angle=-90, end angle=90, radius=0.1];
\draw[popgold, very thick] (3.5,-0.2) -- (3.9,-0.2);
\draw[popgold, very thick] (3.5,0.2) -- (3.9,0.2);
\node[below, align=center, poporange] at (2.6,-0.85) {primaire\\ (beaucoup de spires)};
\node[below, align=center, popgold!60!popdark] at (3.7,-0.85) {secondaire\\ (peu de spires)};
\node[above, popdark] at (3.05,0.75) {noyau de fer};
% Fils du secondaire vers le pont (entrées gauche et droite du losange)
\draw[thick] (3.9,0.2) -- (4.3,0.2) -- (4.3,0.0) -- (4.8,0.0);
\draw[thick] (3.9,-0.2) -- (4.5,-0.2) -- (4.5,-1.6) -- (7.6,-1.6) -- (7.6,0.0);
% Pont de diodes en losange : sommets G(4.8,0) D(7.6,0) H(6.2,1.0) B(6.2,-1.0)
% Diode 1 : G -> H
\draw[thick] (4.8,0) -- (5.35,0.4);
\draw[fill=popgreen, draw=popgreen] (5.35,0.32) -- (5.35,0.48) -- (5.65,0.62) -- cycle;
\draw[popgreen, very thick] (5.57,0.5) -- (5.73,0.74);
\draw[thick] (5.65,0.62) -- (6.2,1.0);
% Diode 2 : D -> H
\draw[thick] (7.6,0) -- (7.05,0.4);
\draw[fill=popgreen, draw=popgreen] (7.05,0.32) -- (7.05,0.48) -- (6.75,0.62) -- cycle;
\draw[popgreen, very thick] (6.83,0.5) -- (6.67,0.74);
\draw[thick] (6.75,0.62) -- (6.2,1.0);
% Diode 3 : B -> G
\draw[thick] (6.2,-1.0) -- (5.65,-0.62);
\draw[fill=popgreen, draw=popgreen] (5.73,-0.74) -- (5.57,-0.5) -- (5.35,-0.4) -- cycle;
\draw[popgreen, very thick] (5.35,-0.32) -- (5.35,-0.48);
\draw[thick] (5.35,-0.4) -- (4.8,0);
% Diode 4 : B -> D
\draw[thick] (6.2,-1.0) -- (6.75,-0.62);
\draw[fill=popgreen, draw=popgreen] (6.67,-0.74) -- (6.83,-0.5) -- (7.05,-0.4) -- cycle;
\draw[popgreen, very thick] (7.05,-0.32) -- (7.05,-0.48);
\draw[thick] (7.05,-0.4) -- (7.6,0);
% Sortie + (sommet haut) et - (sommet bas)
\draw[thick] (6.2,1.0) -- (6.2,1.6) -- (8.8,1.6);
\draw[thick] (6.2,-1.0) -- (6.2,-2.2) -- (8.8,-2.2);
% Appareil
\draw[fill=poppurpleL, draw=poppurple, thick] (8.8,-2.5) rectangle (10.8,1.9);
\node[align=center] at (9.8,-0.3) {Appareil\\ (charge)};
\node[poppurple] at (8.55,1.6) {\textbf{+}};
\node[poppurple] at (8.55,-2.2) {\textbf{--}};
% Étiquettes
\node[align=center, popgreen!60!popdark] at (6.2,2.1) {pont de 4 diodes (redresseur)};
\draw[->, popgreen!60!popdark] (6.2,1.95) -- (6.2,1.15);
\node[align=center, popblue] at (0.7,-1.2) {entrée secteur};
\node[align=center, poppurple] at (9.8,-3.0) {sortie basse tension continue};
\node[align=center, popmuted] at (5.0,0.55) {12 V $\sim$};
\end{tikzpicture}
\end{center}
\legende{Circuit interne simplifié d'un adaptateur secteur : la prise 220 V alimente le primaire du transformateur ; le secondaire (moins de spires) délivre une basse tension alternative au pont de quatre diodes montées en losange, qui impose au courant de toujours traverser l'appareil dans le même sens. La sortie possède un pôle + et un pôle -- fixes.}
\end{popfigure}

\begin{methode}[Activité de lecture : identifier les éléments sur le schéma d'un adaptateur]
Face au schéma d'un adaptateur (comme celui de la figure ci-dessus), procède dans l'ordre :
\begin{enumerate}
\item \textbf{Repère l'entrée} : c'est le côté relié à la prise, marqué 220 V $\sim$.
\item \textbf{Repère le transformateur} : cherche les deux bobinages face à face séparés par des traits verticaux (le noyau de fer). Le bobinage relié au secteur est le primaire ; l'autre est le secondaire.
\item \textbf{Repère les diodes} : cherche les symboles en triangle avec une barre, groupés par quatre en « pont » (en losange).
\item \textbf{Repère la sortie} : c'est le côté relié à l'appareil, marqué d'un pôle + et d'un pôle -- (signe d'une tension continue) ou de la mention « continu » (symbole DC).
\item \textbf{Vérifie l'ordre} : en partant de la prise, tu dois rencontrer le transformateur \emph{avant} les diodes.
\end{enumerate}
\end{methode}

\begin{exemplebox}[Exemple chiffré : l'adaptateur d'une radio 220 V $\sim$ / 9 V continu]
Sur la plaque d'un adaptateur de radio, on lit : « Entrée : 220 V $\sim$ ; Sortie : 9 V continu (DC) ». Voici ce qui se passe à l'intérieur :
\begin{itemize}
\item le \textbf{transformateur} abaisse la tension de 220 V $\sim$ à environ 9 V $\sim$ (toujours alternatif, même fréquence 50 Hz) ;
\item le \textbf{pont de diodes} redresse cette tension : elle ne change plus de polarité ;
\item après filtrage, la \textbf{sortie} est une tension continue d'environ 9 V, qui alimente la radio en toute sécurité.
\end{itemize}
Les deux indications de la plaque racontent donc les deux fonctions : 220 V $\sim$ (avant transformation) et 9 V continu (après transformation et redressement).
\end{exemplebox}

\courssub{Vérification expérimentale globale : la sortie est-elle vraiment continue ?}

Comment prouver, au laboratoire, que la tension de sortie d'un adaptateur est bien continue ? Un simple multimètre suffit, à condition de l'utiliser dans les deux positions : DC (courant continu) et AC (courant alternatif, symbole $\sim$).

\begin{experience}[Mesure de la tension de sortie d'un adaptateur]
\textbf{Matériel :} un adaptateur secteur (par exemple 12 V continu), un multimètre, deux fils de connexion.

\textbf{Protocole :}
\begin{enumerate}
\item Branche l'adaptateur sur la prise (manipulation réalisée par le professeur ou sous sa surveillance).
\item Place le multimètre en voltmètre, position \textbf{DC} (calibre 20 V), et mesure la tension entre les bornes de sortie. Note la valeur.
\item Sans rien changer d'autre, place le multimètre en position \textbf{AC} (calibre 20 V) et mesure à nouveau. Note la valeur.
\end{enumerate}

\textbf{Observations :}
\begin{itemize}
\item En position DC, le voltmètre indique une valeur proche de la tension annoncée (par exemple \SI{12}{V}).
\item En position AC, le voltmètre indique une valeur quasi nulle (proche de \SI{0}{V}).
\end{itemize}

\textbf{Interprétation :} un voltmètre en position DC ne mesure correctement que les tensions continues ; en position AC, il ne mesure correctement que les tensions alternatives. Ici, la tension est « vue » en DC et « invisible » en AC : la sortie de l'adaptateur est donc bien une tension \textbf{continue}.
\end{experience}

\begin{methode}[Vérifier la nature de la tension de sortie d'un adaptateur]
\begin{enumerate}
\item Identifier sur la plaque de l'adaptateur la tension de sortie annoncée et sa nature ($\sim$ pour alternatif, DC ou « continu » pour continu).
\item Choisir un calibre de voltmètre supérieur à la tension annoncée (par exemple 20 V pour un adaptateur 9 V ou 12 V).
\item Mesurer en position \textbf{DC} : une valeur nette et stable confirme la présence d'une tension continue.
\item Mesurer en position \textbf{AC} : une valeur quasi nulle confirme qu'il n'y a (presque) plus de composante alternative.
\item Conclure : si DC donne une valeur et AC presque rien, la sortie est continue ; l'adaptateur contient donc un transformateur \emph{et} un redresseur.
\end{enumerate}
\end{methode}

\begin{attention}[Pièges de la mesure]
\begin{itemize}
\item \textbf{Ne jamais} mesurer la tension de sortie en position « ampèremètre » : tu créerais un court-circuit qui grillerait le fusible du multimètre ou l'adaptateur.
\item En position DC, si l'affichage est négatif, ce n'est pas une panne : tu as simplement inversé les fils. Cela prouve au contraire que la tension a une polarité fixe, donc qu'elle est continue.
\item Une toute petite valeur en position AC (quelques dixièmes de volt) est normale : c'est l'ondulation résiduelle, c'est-à-dire le filtrage qui n'est pas parfait.
\end{itemize}
\end{attention}

\begin{exoresolu}[Lecture de plaque et mesures]
Sur la plaque d'un adaptateur, on lit : « Input : 220 V $\sim$ 50 Hz — Output : 12 V DC ». Un élève mesure la tension de sortie : il obtient \SI{11,8}{V} en position DC et \SI{0,2}{V} en position AC.
\begin{enumerate}
\item Quelles sont les deux fonctions réalisées par cet adaptateur ?
\item Les mesures sont-elles cohérentes avec la plaque ? Justifier.
\end{enumerate}
\tcblower
\textbf{Solution détaillée.}
\begin{enumerate}
\item La plaque indique une entrée alternative 220 V et une sortie continue 12 V. L'adaptateur réalise donc :
\begin{itemize}
\item la \textbf{transformation} (abaissement de 220 V à 12 V) par le transformateur ;
\item le \textbf{redressement} (alternatif $\rightarrow$ continu) par le pont de diodes, suivi d'un filtrage.
\end{itemize}
\item Oui, les mesures sont cohérentes :
\begin{itemize}
\item en position DC, on lit \SI{11,8}{V}, très proche des \SI{12}{V} annoncés (l'écart de \SI{0,2}{V} vient des tolérances de fabrication) : la tension continue est bien présente ;
\item en position AC, on lit seulement \SI{0,2}{V}, valeur quasi nulle : il ne reste presque plus de composante alternative, le redressement et le filtrage ont bien fonctionné.
\end{itemize}
Conclusion : la sortie est bien une tension continue d'environ 12 V, conforme à la mention DC de la plaque.
\end{enumerate}
\end{exoresolu}

\courssub{Et les chargeurs de téléphone modernes ?}

Les chargeurs de téléphone actuels sont beaucoup plus petits et plus légers que les vieux adaptateurs à transformateur. Ont-ils changé de principe ? Pas vraiment.

\begin{savaistu}[Adaptateurs miniatures et électronique de commutation]
Les chargeurs modernes utilisent l'\cle{électronique de commutation} : au lieu d'un gros transformateur fonctionnant à 50 Hz, ils découpent la tension à très haute fréquence grâce à des composants électroniques rapides, ce qui permet d'utiliser un transformateur minuscule. Mais les fonctions restent exactement les mêmes : \textbf{abaisser} la tension du secteur puis \textbf{redresser} pour fournir une basse tension continue (souvent 5 V continu sur un port USB). La chaîne de conversion que tu as apprise reste donc valable pour tous les adaptateurs, anciens ou modernes.

Autre piège de la vie courante : certains adaptateurs fournissent une sortie \textbf{alternative} (plaque marquée « Output : 12 V $\sim$ » ou « 12 V AC »). Ils contiennent un transformateur mais \emph{pas} de redresseur. Si tu branches un tel adaptateur sur un appareil prévu pour du continu, il ne fonctionnera pas (ou sera abîmé). Moralité : \textbf{lis toujours la plaque} avant de brancher un adaptateur — vérifie la valeur (en volts) \emph{et} la nature ($\sim$ ou continu) de la tension de sortie.
\end{savaistu}

\begin{exercice}[À toi de jouer]
\begin{enumerate}
\item Recopie et complète la chaîne de conversion : Secteur 220 V $\sim$ $\rightarrow$ \ldots\ldots\ldots $\rightarrow$ 6 V $\sim$ $\rightarrow$ \ldots\ldots\ldots $\rightarrow$ 6 V continu.
\item Un élève affirme : « L'adaptateur redresse d'abord le 220 V, puis abaisse la tension. » Explique pourquoi c'est faux.
\item On mesure la sortie d'un adaptateur : \SI{0,1}{V} en DC et \SI{9}{V} en AC. Quelle est la nature de la tension de sortie ? Cet adaptateur contient-il un redresseur ?
\end{enumerate}
\end{exercice}

\begin{corrige}[Exercice]
\begin{enumerate}
\item Secteur 220 V $\sim$ $\rightarrow$ \textbf{transformateur abaisseur} $\rightarrow$ 6 V $\sim$ $\rightarrow$ \textbf{pont de diodes (redresseur)} $\rightarrow$ 6 V continu.
\item C'est faux car le transformateur ne fonctionne qu'avec une tension alternative : s'il recevait une tension déjà redressée, il ne pourrait pas l'abaisser. On transforme (abaisse) d'abord, puis on redresse.
\item La tension est « vue » en AC (\SI{9}{V}) et quasi nulle en DC : c'est une tension \textbf{alternative}. Cet adaptateur ne contient donc \textbf{pas} de redresseur : c'est un simple transformateur abaisseur (sortie marquée $\sim$).
\end{enumerate}
\end{corrige}

\begin{aretenir}[Bilan de la section]
\begin{itemize}
\item Un adaptateur secteur enchaîne deux fonctions, dans cet ordre : \textbf{transformation} (abaissement de la tension alternative) puis \textbf{redressement} (conversion en tension continue), souvent complétées par un filtrage.
\item La chaîne de conversion s'écrit : 220 V $\sim$ $\rightarrow$ transformateur $\rightarrow$ basse tension $\sim$ $\rightarrow$ pont de diodes $\rightarrow$ tension redressée $\rightarrow$ (filtrage) $\rightarrow$ basse tension continue.
\item On vérifie la nature continue de la sortie avec un voltmètre : valeur nette en position DC, valeur quasi nulle en position AC.
\item Les chargeurs modernes remplissent les mêmes fonctions grâce à l'électronique de commutation ; certains adaptateurs n'ont pas de redresseur (sortie $\sim$) : toujours lire la plaque.
\end{itemize}
\end{aretenir}

\coursec{Sécurité et bon usage des adaptateurs}

Dans la section précédente, nous avons démonté (sur le papier !) la chaîne de conversion d'un adaptateur : la prise \SI{220}{V} alimente le transformateur qui abaisse la tension, puis le pont de diodes redresse le courant, le condensateur le filtre et l'appareil reçoit une tension continue de quelques volts. Nous savons donc maintenant \emph{comment} cet appareil fonctionne. Reste une question vitale : comment l'utiliser \emph{sans danger} ? Car si la sortie de l'adaptateur délivre une tension faible et inoffensive, son entrée, elle, reste branchée sur le secteur \SI{220}{V} — une tension mortelle. Cette dernière section te donne les règles de sécurité et de bon usage que tout élève de 3ème doit connaître et appliquer à la maison.

\courssub{Pourquoi faut-il être prudent avec un adaptateur ?}

\courssub{Le danger vient du secteur, pas de la sortie}

Reprenons la situation de départ. Au quartier, tu branches le chargeur du téléphone de ta mère sur la prise murale alimentée par ENEO. À l'intérieur du petit boîtier plastique, le primaire du transformateur est \emph{directement relié} aux deux broches de la prise, donc au secteur \SI{220}{V}. Or le corps humain conduit le courant électrique : une tension de \SI{220}{V} appliquée entre deux points du corps provoque un passage de courant à travers les tissus, le cœur et les muscles. C'est le \cle{choc électrique}, qui peut être mortel.

\begin{attention}[Danger : un adaptateur reste relié au secteur \SI{220}{V}]
Tant que l'adaptateur est branché sur la prise, l'intérieur de son boîtier est au potentiel du secteur \SI{220}{V}, même si aucun appareil n'est connecté à sa sortie et même si l'appareil est éteint. Il est donc \textbf{formellement interdit d'ouvrir un adaptateur} et \textbf{interdit d'utiliser un adaptateur endommagé}. La tension de sortie (5 V, 9 V, 12 V) est sans danger ; la tension d'entrée (\SI{220}{V}) est mortelle.
\end{attention}

\begin{aretenir}[L'idée essentielle]
Un adaptateur est une frontière entre deux mondes : le côté « secteur » (\SI{220}{V}, dangereux, à ne jamais toucher) et le côté « sortie » (basse tension continue, sans danger). Toutes les règles de sécurité consistent à protéger cette frontière : garder le boîtier fermé et intact, garder les contacts au sec, et choisir la bonne tension de sortie.
\end{aretenir}

\courssub{Les six règles d'or}

\textbf{Règle 1 : ne jamais ouvrir un adaptateur, même débranché.}\par
On pourrait croire qu'une fois débranché, l'adaptateur est inoffensif. C'est faux ! Souviens-toi de la chaîne de conversion : après le pont de diodes se trouve un \cle{condensateur} de filtrage. Or un condensateur est un réservoir de charges électriques : il peut \emph{rester chargé} plusieurs minutes après le débranchement et délivrer une décharge douloureuse à qui touche ses bornes. De plus, ouvrir le boîtier détruit l'isolation prévue par le fabricant. Si un adaptateur tombe en panne, on ne le répare pas soi-même : on le remplace ou on le confie à un technicien qualifié.

\textbf{Règle 2 : ne pas utiliser un adaptateur endommagé.}\par
Un câble dénudé (fils de cuivre visibles), un boîtier fissuré, une trace de brûlure ou de plastique fondu sont des signaux d'alarme. L'isolation n'est plus garantie : le secteur \SI{220}{V} peut devenir accessible au toucher, ou deux fils peuvent se toucher et provoquer un \cle{court-circuit} (échauffement brutal, étincelles, début d'incendie). Un adaptateur endommagé se jette (dans une filière de déchets électroniques si possible) et se remplace. Aucun bout de ruban adhésif ne rendra sa sécurité à un câble dénudé.

\textbf{Règle 3 : choisir une tension de sortie exactement égale à celle de l'appareil.}\par
Chaque appareil est conçu pour fonctionner sous une tension précise, indiquée sur son étiquette (par exemple « \emph{input : 9 V ⎓} »). L'adaptateur de remplacement doit délivrer \emph{exactement} cette tension : 9 V pour un appareil 9 V, jamais 12 V, jamais 6 V. Nous étudierons en détail ce qui se passe en cas d'erreur dans l'étude de cas ci-dessous.

\textbf{Règle 4 : vérifier la polarité et l'intensité maximale.}\par
La tension de sortie est continue : elle possède donc un pôle $+$ et un pôle $-$. Sur la fiche cylindrique de sortie, le $+$ peut être le contact central ou le contact extérieur selon les modèles. Un petit symbole dessiné sur l'adaptateur et sur l'appareil l'indique : il faut qu'ils \emph{correspondent}. Brancher une polarité inversée peut détruire les composants électroniques de l'appareil. Il faut aussi vérifier que l'adaptateur peut fournir une intensité au moins égale à celle demandée par l'appareil (voir la méthode plus bas).

\textbf{Règle 5 : jamais les mains mouillées ; débrancher en cas d'orage ou d'odeur de brûlé.}\par
L'eau (surtout l'eau du robinet, qui contient des sels dissous) conduit le courant. Manipuler une prise ou un adaptateur avec les mains mouillées, c'est offrir au courant un chemin à travers ton corps. De même, en cas d'orage, la foudre peut provoquer des \cle{surtensions} brutales sur le réseau : on débranche les appareils sensibles. Enfin, une odeur de brûlé ou un boîtier anormalement chaud signale un défaut interne : on débranche immédiatement, sans attendre.

\textbf{Règle 6 : ne pas surcharger les multiprises.}\par
Dans beaucoup de maisons, une seule prise murale alimente une multiprise sur laquelle on branche la télévision, le décodeur, deux chargeurs de téléphone et le ventilateur. Chaque appareil tire un courant, et tous ces courants s'additionnent dans le même câble de la multiprise. Si le courant total dépasse ce que le câble peut supporter, celui-ci \emph{chauffe} (effet Joule) : la gaine fond, et l'incendie peut démarrer. On répartit donc les gros consommateurs sur des prises différentes et on évite les multiprises « en cascade ».

\begin{popfigure}
\begin{center}
\begin{tikzpicture}[scale=1, every node/.style={font=\small}]
% --- Pictogramme 1 : mains mouillées barrées ---
\begin{scope}[shift={(0,0)}]
  \draw[fill=popblueL, rounded corners=2pt] (-0.9,-0.6) rectangle (0.9,0.6);
  \node[font=\footnotesize\bfseries, popdark] at (0,0.9) {Mains mouillées};
  % main stylisée
  \draw[fill=popgoldL, rounded corners=1pt] (-0.45,-0.35) rectangle (0.15,0.15);
  \foreach \x in {-0.42,-0.28,-0.14,0.0} {\draw[fill=popgoldL, rounded corners=1pt] (\x,0.15) rectangle (\x+0.1,0.42);}
  % gouttes
  \draw[fill=popblue] (0.45,0.1) circle (0.07);
  \draw[fill=popblue] (0.6,-0.15) circle (0.07);
  \draw[fill=popblue] (0.45,-0.4) circle (0.07);
  % interdiction
  \draw[popink, line width=2pt] (0,0) circle (0.85);
  \draw[popink, line width=2pt] (-0.6,0.6) -- (0.6,-0.6);
\end{scope}
% --- Pictogramme 2 : adaptateur ouvert barré ---
\begin{scope}[shift={(4.2,0)}]
  \node[font=\footnotesize\bfseries, popdark] at (0,0.9) {Boîtier ouvert};
  % adaptateur ouvert
  \draw[fill=popmuted!30, rounded corners=2pt] (-0.7,-0.45) rectangle (0.3,0.45);
  \draw[fill=popmuted!50, rounded corners=2pt, rotate=25] (0.35,-0.5) rectangle (1.0,0.1);
  % composants visibles
  \draw[fill=poporangeL] (-0.5,-0.25) rectangle (-0.15,0.25);
  \draw[fill=popgreenL] (-0.05,-0.15) circle (0.12);
  \draw[popdark] (0.3,0.2) -- (0.75,0.35);
  % éclair danger
  \draw[fill=popgold] (-0.35,0.28) -- (-0.2,0.05) -- (-0.3,0.05) -- (-0.15,-0.2) -- (-0.28,-0.02) -- (-0.18,-0.02) -- cycle;
  % interdiction
  \draw[popink, line width=2pt] (0.1,0) circle (0.85);
  \draw[popink, line width=2pt] (-0.5,0.6) -- (0.7,-0.6);
\end{scope}
% --- Pictogramme 3 : symbole de polarité ---
\begin{scope}[shift={(8.4,0)}]
  \node[font=\footnotesize\bfseries, popdark] at (0,0.9) {Polarité de la fiche};
  % symbole polarité standard
  \draw[popdark, line width=1.2pt] (-0.75,0.1) circle (0.16);
  \draw[popdark, line width=1.2pt] (-0.75,-0.45) -- (-0.75,-0.15);
  \draw[popdark, line width=1.2pt] (0.75,0.1) -- (0.75,-0.2);
  \draw[popdark, line width=1.2pt] (0.75,-0.2) arc (0:-180:0.18);
  \draw[popdark, line width=1.2pt] (-0.59,0.1) -- (-0.15,0.1);
  \draw[popdark, line width=1.2pt] (0.15,0.1) -- (0.75,0.1);
  \node[font=\bfseries, popdark] at (-0.95,0.1) {$+$};
  \node[font=\bfseries, popdark] at (0.98,0.1) {$-$};
  \node[font=\footnotesize, popdark, align=center] at (0,-0.75) {le $+$ au centre\\ le $-$ à l'extérieur};
\end{scope}
\end{tikzpicture}
\end{center}
\legende{Pictogrammes de sécurité. À gauche : il est interdit de manipuler un adaptateur ou une prise avec les mains mouillées. Au centre : il est interdit d'ouvrir le boîtier d'un adaptateur, même débranché (le secteur \SI{220}{V} et le condensateur chargé restent dangereux). À droite : le symbole de polarité indique que le contact central de la fiche est le pôle $+$ et le contact extérieur le pôle $-$ ; il doit être identique sur l'adaptateur et sur l'appareil.}
\end{popfigure}

\courssub{Comment choisir un adaptateur de remplacement ?}

Ton adaptateur est perdu ou en panne : comment le remplacer sans risque ? Trois informations figurent sur l'étiquette de l'appareil (ou de l'ancien adaptateur) : la \cle{tension de sortie} (en volts), la \cle{polarité} (symbole $+/-$) et l'\cle{intensité maximale} (en ampères ou milliampères).

\begin{methode}[Choisir un adaptateur de remplacement en trois vérifications]
\begin{enumerate}
\item \textbf{Même tension de sortie} : la tension de l'adaptateur neuf doit être \emph{exactement égale} à celle indiquée sur l'appareil (9 V pour 9 V). Ni plus, ni moins.
\item \textbf{Même polarité} : le symbole de polarité de la fiche doit être identique sur l'adaptateur et sur l'appareil.
\item \textbf{Intensité égale ou supérieure} : l'intensité maximale que peut fournir l'adaptateur doit être \emph{au moins égale} à celle demandée par l'appareil. Une intensité disponible plus grande ne pose aucun problème.
\end{enumerate}
\end{methode}

\begin{attention}[Erreur fréquente : « un adaptateur plus puissant va griller mon appareil »]
Beaucoup d'élèves (et d'adultes !) croient qu'un adaptateur affichant \SI{1}{A} « force » un courant de \SI{1}{A} dans l'appareil. C'est faux. L'intensité indiquée sur l'adaptateur est l'intensité \emph{maximale disponible}, c'est-à-dire ce qu'il \emph{peut} fournir. C'est l'appareil qui « décide » du courant qu'il prend, selon ses besoins : un appareil prévu pour \SI{500}{mA} branché sur un adaptateur \SI{1}{A} ne prendra que \SI{500}{mA}. En revanche, un adaptateur dont l'intensité maximale est \emph{trop faible} (par exemple \SI{300}{mA} pour un appareil qui demande \SI{500}{mA}) va surchauffer et s'user prématurément. Règle pratique : intensité de l'adaptateur $\geq$ intensité demandée par l'appareil.
\end{attention}

\begin{exemplebox}[Remplacer un adaptateur 9 V ⎓ 500 mA]
L'adaptateur d'un poste de radio porte la mention « sortie : 9 V ⎓ 500 mA ». Il est en panne. Trois possibilités au marché :
\begin{itemize}
\item Adaptateur A : 9 V ⎓ 1 A, même polarité $\Rightarrow$ \textbf{acceptable} : même tension, même polarité, intensité disponible supérieure ($\SI{1}{A} > \SI{500}{mA}$). La radio ne prendra que ses 500 mA.
\item Adaptateur B : 9 V ⎓ 300 mA $\Rightarrow$ \textbf{déconseillé} : l'intensité disponible est insuffisante, l'adaptateur va chauffer.
\item Adaptateur C : 12 V ⎓ 1 A $\Rightarrow$ \textbf{interdit} : la tension est trop élevée, la radio peut être détruite.
\end{itemize}
\end{exemplebox}

\courssub{Étude de cas : que se passe-t-il si l'on branche un adaptateur 12 V sur un appareil prévu pour 5 V ?}

\begin{experience}[Observation et raisonnement]
\textbf{Document observé.} Un téléphone portable dont l'étiquette indique « entrée : 5 V ⎓ » est branché par erreur sur un adaptateur délivrant 12 V. Au bout de quelques secondes, le boîtier du téléphone chauffe anormalement, une odeur de composant brûlé se dégage et l'écran ne s'allume plus.

\textbf{Interprétation pas à pas.}
\begin{enumerate}
\item Les composants électroniques du téléphone (circuits intégrés, condensateurs, batterie) sont dimensionnés pour supporter une tension maximale voisine de 5 V.
\item En appliquant 12 V, soit plus du double, on impose aux composants une contrainte électrique qu'ils ne peuvent supporter : l'isolation interne de certains composants claque.
\item Le courant qui traverse alors les circuits devient excessif ; par effet Joule, la puissance dissipée en chaleur augmente fortement : d'où l'échauffement et l'odeur de brûlé.
\item Les composants détruits ne se réparent pas : l'appareil est hors d'usage.
\end{enumerate}

\textbf{Conclusion.} Une tension de sortie \emph{supérieure} à la tension prévue détériore ou détruit l'appareil, parfois avec dégagement de chaleur et risque d'incendie. Une tension \emph{inférieure} (par exemple 3 V au lieu de 5 V) ne détruit généralement pas l'appareil, mais celui-ci fonctionne mal ou pas du tout. D'où la règle absolue : la tension de l'adaptateur doit être \emph{exactement} celle de l'appareil.
\end{experience}

\begin{exoresolu}[Le bon choix au quartier]
\textbf{Énoncé.} Marlyse a perdu l'adaptateur de sa lampe rechargeable. Sur la lampe, on lit : « entrée : 6 V ⎓ 800 mA », avec le symbole de polarité « $+$ au centre ». Au magasin, le vendeur propose trois adaptateurs :
\begin{itemize}
\item Adaptateur 1 : sortie 12 V ⎓ 1 A, polarité « $+$ au centre » ;
\item Adaptateur 2 : sortie 6 V ⎓ 500 mA, polarité « $+$ au centre » ;
\item Adaptateur 3 : sortie 6 V ⎓ 1 A, polarité « $+$ au centre ».
\end{itemize}
Lequel doit-elle choisir ? Justifier en écartant les deux autres.
\tcblower
\textbf{Solution détaillée.}
Appliquons la méthode en trois vérifications.

\emph{Adaptateur 1 (12 V ⎓ 1 A)} : la tension de sortie (12 V) est le double de la tension prévue (6 V). D'après l'étude de cas, une tension trop élevée détruit les composants de la lampe. \textbf{À écarter.}

\emph{Adaptateur 2 (6 V ⎓ 500 mA)} : la tension est correcte (6 V) et la polarité est bonne. Mais l'intensité maximale disponible (500 mA) est \emph{inférieure} à celle demandée par la lampe (800 mA) : l'adaptateur va surchauffer et s'user rapidement. \textbf{À écarter.}

\emph{Adaptateur 3 (6 V ⎓ 1 A)} : tension exacte (6 V), polarité identique, et intensité disponible ($\SI{1}{A} = \SI{1000}{mA}$) \emph{supérieure} aux 800 mA demandés. La lampe ne prendra que les 800 mA dont elle a besoin : aucun danger. \textbf{C'est le bon choix.}

\textbf{Conclusion} : Marlyse doit acheter l'adaptateur 3 (6 V ⎓ 1 A, $+$ au centre).
\end{exoresolu}

\courssub{Tableau récapitulatif de correspondance appareil / adaptateur}

\begin{center}
\begin{tabularx}{\linewidth}{|l|X|X|X|}\hline
\rowcolor{popblueL} \textbf{Critère} & \textbf{Où le lire} & \textbf{Règle} & \textbf{Exemple} \\ \hline
Tension de sortie & Étiquette de l'appareil (V ⎓) & Exactement égale & Appareil 9 V $\Rightarrow$ adaptateur 9 V \\ \hline
Polarité & Symbole $+/-$ sur la fiche & Identique & $+$ au centre des deux côtés \\ \hline
Intensité maximale & Étiquette (A ou mA) & Égale ou supérieure & Appareil 500 mA $\Rightarrow$ adaptateur 500 mA ou 1 A \\ \hline
État du matériel & Inspection visuelle & Aucun dommage & Pas de câble dénudé ni de boîtier fissuré \\ \hline
\end{tabularx}
\end{center}

\begin{savaistu}[Un adaptateur branché à vide consomme quand même de l'énergie]
Touche le boîtier de ton chargeur resté branché alors que le téléphone est déjà déconnecté : il est tiède ! C'est la preuve qu'il consomme un peu d'énergie électrique, même « à vide » : le primaire du transformateur reste traversé par un petit courant, et une partie de cette énergie est perdue en chaleur (effet Joule). Multipliée par des millions de chargeurs branchés en permanence, cette consommation invisible devient énorme à l'échelle d'un pays. Le geste économique et écologique est simple : \textbf{débrancher l'adaptateur après usage}, ou utiliser une multiprise à interrupteur que l'on éteint le soir. Ton porte-monnaie et la planète te diront merci !
\end{savaistu}

\begin{exercice}[À toi de jouer]
\begin{enumerate}
\item Citer deux raisons pour lesquelles il est interdit d'ouvrir un adaptateur, même débranché.
\item L'adaptateur d'un ventilateur rechargeable indique « sortie : 5 V ⎓ 2 A ». Peut-on le remplacer par un adaptateur 5 V ⎓ 1 A ? Par un adaptateur 9 V ⎓ 2 A ? Par un adaptateur 5 V ⎓ 3 A ? Justifier chaque réponse.
\item Pourquoi est-il dangereux de brancher de nombreux adaptateurs sur une même multiprise ?
\end{enumerate}
\end{exercice}

\begin{corrige}[Exercice — À toi de jouer]
\begin{enumerate}
\item D'une part, l'intérieur du boîtier est relié au secteur \SI{220}{V} lorsque l'adaptateur est branché (risque de choc électrique mortel) ; d'autre part, le condensateur de filtrage peut rester chargé après le débranchement et provoquer une décharge. Ouvrir détruit aussi l'isolation prévue par le fabricant.
\item Adaptateur 5 V ⎓ 1 A : non, l'intensité disponible (\SI{1}{A}) est inférieure aux \SI{2}{A} demandés, il va surchauffer. Adaptateur 9 V ⎓ 2 A : non, la tension est trop élevée, le ventilateur peut être détruit. Adaptateur 5 V ⎓ 3 A : oui, même tension, intensité disponible supérieure ; le ventilateur ne prendra que ses \SI{2}{A}.
\item Les courants de tous les appareils s'additionnent dans le câble de la multiprise. Si le courant total est trop grand, le câble chauffe par effet Joule : la gaine peut fondre et provoquer un incendie.
\end{enumerate}
\end{corrige}

\begin{aretenir}[L'essentiel de la section]
\begin{itemize}
\item L'adaptateur est branché sur le secteur \SI{220}{V} : on ne l'ouvre jamais, même débranché (condensateur chargé), et on n'utilise jamais un adaptateur endommagé.
\item Pas de manipulation les mains mouillées ; on débranche en cas d'orage ou d'odeur de brûlé ; on ne surcharge pas les multiprises.
\item Pour remplacer un adaptateur : \textbf{même tension}, \textbf{même polarité}, \textbf{intensité égale ou supérieure}. Une intensité disponible plus grande est sans danger : c'est l'appareil qui prélève le courant dont il a besoin.
\item Une tension trop élevée détruit l'appareil ; une tension trop faible l'empêche de fonctionner correctement.
\item Un adaptateur branché à vide consomme un peu d'énergie : on le débranche après usage.
\end{itemize}
\end{aretenir}

\newpage

\coursec{Activité d'intégration}

\courssub{Objectifs de l'activité}

À la fin de cette activité, tu seras capable de :

\begin{itemize}
\item lire et interpréter la \cle{plaque signalétique} d'un adaptateur secteur (tension d'entrée, tension de sortie, nature du courant, intensité maximale) ;
\item reconnaître sur un \cle{oscillogramme} une tension alternative sinusoïdale et une tension continue (redressée puis filtrée) ;
\item identifier les fonctions contenues dans un adaptateur : \cle{transformation} seule ou transformation + \cle{redressement} ;
\item choisir l'adaptateur compatible avec un appareil donné et justifier le rejet des autres ;
\item schématiser la \cle{chaîne de conversion} d'un adaptateur (transformateur $\rightarrow$ redresseur $\rightarrow$ filtre) ;
\item appliquer et formuler les règles de sécurité liées à l'utilisation des adaptateurs.
\end{itemize}

\courssub{Situation}

\textbf{Document 1 : le kiosque de recharge.}

Dans le quartier Nkol-Eton à Yaoundé, Monsieur Etoundi tient un kiosque de recharge de téléphones, très fréquenté les jours de délestage d'ENEO. Les clients déposent leurs appareils et repartent avec une batterie pleine pour quelques centaines de francs CFA. Ce matin, un jeune client, Arnaud, élève en classe de 3ème, remarque que trois adaptateurs traînent sur la table, branchés sur une multiprise. Il lit attentivement leurs plaques signalétiques :

\begin{center}
\begin{tabularx}{\linewidth}{|l|X|X|}
\hline
\rowcolor{popblueL} \textbf{Adaptateur} & \textbf{Entrée (INPUT)} & \textbf{Sortie (OUTPUT)} \\
\hline
A & \SI{220}{V} $\sim$ (alternatif), \SI{50}{Hz} & \SI{12}{V} continu, $I_{\max} = \SI{1}{A}$ \\
\hline
B & \SI{220}{V} $\sim$ (alternatif), \SI{50}{Hz} & \SI{5}{V} continu, $I_{\max} = \SI{2}{A}$ \\
\hline
C & \SI{220}{V} $\sim$ (alternatif), \SI{50}{Hz} & \SI{5}{V} $\sim$ (alternatif), $I_{\max} = \SI{1}{A}$ \\
\hline
\end{tabularx}
\end{center}

Sur la table se trouvent aussi trois appareils à alimenter :

\begin{itemize}
\item un \textbf{téléphone portable} dont le chargeur d'origine exige : \SI{5}{V} continu, courant jusqu'à \SI{2}{A} ;
\item un \textbf{poste radio} fonctionnant sous \SI{9}{V} continu ;
\item une \textbf{lampe à LED} de camping fonctionnant sous \SI{12}{V} continu, courant \SI{0,8}{A}.
\end{itemize}

\textbf{Document 2 : les oscillogrammes.}

Arnaud, qui dispose au club scientifique de son collège d'un oscilloscope, branche successivement la sortie de deux des adaptateurs et obtient les deux oscillogrammes suivants (sensibilité verticale : \SI{5}{V}/division ; balayage : \SI{5}{ms}/division).

\begin{popfigure}
\begin{center}
\begin{tikzpicture}[scale=0.9]
% Oscillogramme 1 : sinusoïde
\draw[popline] (0,0) grid[step=0.5] (5,3);
\draw[popdark,->] (0,0) -- (5.3,0) node[below left]{\footnotesize $t$};
\draw[popdark,->] (0,0) -- (0,3.3) node[above left]{\footnotesize $u$};
\draw[popblue,thick,domain=0:5,samples=120] plot (\x,{1.5+1.2*sin(deg(2*3.14159*\x/2))});
\draw[popblue,dashed] (0,1.5) -- (5,1.5);
\node[popdark] at (2.5,-0.6) {\footnotesize Oscillogramme 1 : sortie de l'adaptateur C};
% Oscillogramme 2 : continu
\begin{scope}[xshift=7cm]
\draw[popline] (0,0) grid[step=0.5] (5,3);
\draw[popdark,->] (0,0) -- (5.3,0) node[below left]{\footnotesize $t$};
\draw[popdark,->] (0,0) -- (0,3.3) node[above left]{\footnotesize $u$};
\draw[poporange,thick] (0,1) -- (5,1);
\node[popdark] at (2.5,-0.6) {\footnotesize Oscillogramme 2 : sortie de l'adaptateur B};
\node[poporange] at (4.3,1.3) {\footnotesize \SI{5}{V}};
\end{scope}
\end{tikzpicture}
\end{center}
\legende{Oscillogrammes relevés à la sortie de deux adaptateurs du kiosque. Le premier montre une tension sinusoïdale qui prend des valeurs positives et négatives (tension alternative) ; le second montre une droite horizontale (tension continue).}
\end{popfigure}

Monsieur Etoundi avoue à Arnaud : « Je branche au hasard, et la semaine dernière la radio d'un client a grillé ! » Arnaud décide de mettre de l'ordre dans tout cela.

\courssub{Consignes}

\begin{enumerate}
\item Pour chacun des trois adaptateurs A, B et C, indique les fonctions qu'il contient (transformation seule, ou transformation + redressement). Justifie ta réponse à l'aide de la plaque signalétique et, pour B et C, à l'aide des oscillogrammes.
\item Associe chaque appareil (téléphone, radio, lampe LED) à l'adaptateur qui convient. Pour chaque appareil, explique pourquoi les deux autres adaptateurs sont rejetés (tension inadaptée, nature du courant inadaptée, intensité insuffisante).
\item Rédige pour Monsieur Etoundi une affichette comportant \textbf{trois consignes de sécurité} claires, rédigées en phrases courtes, à coller dans le kiosque.
\item Schématise la chaîne de conversion de l'adaptateur B, du secteur \SI{220}{V} jusqu'à la sortie \SI{5}{V} continu, en nommant chaque étage (transformateur, redresseur, filtre) et en précisant la nature de la tension à chaque étape.
\end{enumerate}

\courssub{Ressources à mobiliser}

\begin{itemize}
\item Un \cle{adaptateur secteur} réalise la conversion : réseau \SI{220}{V} alternatif $\rightarrow$ basse tension adaptée à l'appareil.
\item La fonction \cle{transformation} (transformateur) abaisse la tension alternative : \SI{220}{V} $\sim$ $\rightarrow$ par exemple \SI{5}{V} $\sim$ ou \SI{12}{V} $\sim$. Le transformateur ne fonctionne qu'en courant alternatif.
\item La fonction \cle{redressement} (pont de diodes) convertit la tension alternative en tension unidirectionnelle (qui ne change plus de signe) ; le \cle{filtrage} (condensateur) la rend presque continue.
\item Sur une plaque signalétique : le symbole $\sim$ signifie \emph{alternatif} ; la mention \emph{continu} (souvent notée DC ou symbolisée par un trait plein sur des pointillés) signifie \emph{courant continu}.
\item Sur un oscillogramme : une courbe sinusoïdale qui change de signe traduit une tension \emph{alternative} ; une droite horizontale traduit une tension \emph{continue}.
\item Compatibilité appareil--adaptateur : \textbf{même nature de courant}, \textbf{même tension} (la tension de sortie de l'adaptateur doit être égale à la tension de fonctionnement de l'appareil), et $I_{\max}$ de l'adaptateur $\geq$ courant demandé par l'appareil.
\item Puissance électrique : $P = U \times I$, avec $P$ en watts (W), $U$ en volts (V) et $I$ en ampères (A).
\end{itemize}

\courssub{Démarche et corrigé commenté}

\begin{exoresolu}[Diagnostic du poste de recharge]
Énoncé : à partir des documents 1 et 2, (1) identifier les fonctions de chaque adaptateur, (2) associer chaque appareil à l'adaptateur convenable en justifiant les rejets, (3) rédiger trois consignes de sécurité, (4) décrire la chaîne de conversion de l'adaptateur B.

\tcblower

\textbf{Étape 1 : identification des fonctions de chaque adaptateur.}

On lit d'abord la nature du courant de sortie sur chaque plaque signalétique :

\begin{itemize}
\item \textbf{Adaptateur A} : sortie \SI{12}{V} continu. Or le secteur d'entrée est alternatif. Pour obtenir du continu à partir de l'alternatif, il faut un redressement. A contient donc \textbf{transformation + redressement} (avec filtrage).
\item \textbf{Adaptateur B} : sortie \SI{5}{V} continu. Même raisonnement : B contient \textbf{transformation + redressement}. L'oscillogramme 2 confirme : la tension de sortie est une droite horizontale à \SI{5}{V} (soit 1 division $\times$ \SI{5}{V}/division), donc une tension continue, redressée puis filtrée.
\item \textbf{Adaptateur C} : sortie \SI{5}{V} \emph{alternatif}. La tension a été abaissée mais reste alternative : C ne contient que la fonction \textbf{transformation} (transformateur seul). L'oscillogramme 1 confirme : la courbe est sinusoïdale, elle prend des valeurs positives et négatives, c'est bien une tension alternative. Son amplitude crête est d'environ $1{,}2$ division, soit $1{,}2 \times 5 = \SI{6}{V}$ au maximum, ce qui correspond à une tension efficace voisine de \SI{5}{V} : cohérent avec la plaque.
\end{itemize}

\textbf{Étape 2 : association appareil--adaptateur.}

On applique les critères dans l'ordre : même nature de courant, même tension, puis on vérifie l'intensité.

\begin{center}
\begin{tabularx}{\linewidth}{|l|X|X|X|}
\hline
\rowcolor{popblueL} \textbf{Appareil} & \textbf{A (12 V continu, 1 A)} & \textbf{B (5 V continu, 2 A)} & \textbf{C (5 V $\sim$, 1 A)} \\
\hline
Téléphone (5 V continu, 2 A) & Rejet : tension trop élevée (12 V) & \textbf{Convient} : 5 V continu, 2 A & Rejet : courant alternatif \\
\hline
Radio (9 V continu) & Rejet : $12 \neq 9$ V & Rejet : $5 \neq 9$ V & Rejet : tension et nature inadaptées \\
\hline
Lampe LED (12 V continu, 0,8 A) & \textbf{Convient} : 12 V continu, $0{,}8 < 1$ A & Rejet : tension trop faible & Rejet : tension et nature inadaptées \\
\hline
\end{tabularx}
\end{center}

Conclusions :
\begin{itemize}
\item Le \textbf{téléphone} est alimenté par l'adaptateur \textbf{B} : 5 V continu, et $I_{\max} = \SI{2}{A}$ correspond exactement au besoin. A est rejeté car \SI{12}{V} grillerait le circuit de charge ; C est rejeté car un téléphone exige du courant \emph{continu}, pas alternatif.
\item La \textbf{lampe LED} est alimentée par l'adaptateur \textbf{A} : 12 V continu, et le courant demandé (\SI{0,8}{A}) est inférieur à $I_{\max} = \SI{1}{A}$, donc l'adaptateur n'est pas surchargé. Puissance fournie : $P = U \times I = 12 \times 0{,}8 = \SI{9,6}{W}$.
\item La \textbf{radio} (\SI{9}{V} continu) ne peut être alimentée par \textbf{aucun} des trois adaptateurs : A donne \SI{12}{V} (trop, risque de destruction), B donne \SI{5}{V} (trop peu, la radio ne fonctionnera pas), C donne de l'alternatif. C'est très probablement l'adaptateur A branché sur la radio qui a causé la panne de la semaine dernière ! Monsieur Etoundi doit acheter un adaptateur \SI{9}{V} continu.
\end{itemize}

\textbf{Étape 3 : affichette de sécurité pour le kiosque.}

\begin{center}
\fbox{\parbox{0.85\linewidth}{\textbf{SÉCURITÉ — KIOSQUE DE RECHARGE}
\begin{enumerate}
\item Vérifie toujours la plaque signalétique : la tension et la nature du courant (continu ou alternatif $\sim$) de l'adaptateur doivent correspondre exactement à celles de l'appareil.
\item Ne branche jamais un adaptateur mouillé, endommagé ou chauffant anormalement ; débranche-le en tirant sur la fiche, jamais sur le câble.
\item Ne surcharge pas la multiprise : branche un nombre raisonnable d'adaptateurs et débranche ceux qui ne servent pas (un adaptateur branché à vide chauffe et consomme).
\end{enumerate}}}
\end{center}

\textbf{Étape 4 : chaîne de conversion de l'adaptateur B.}

La tension parcourt trois étapes successives : \SI{220}{V} alternatif $\xrightarrow{\text{transformateur}}$ \SI{5}{V} alternatif $\xrightarrow{\text{redresseur}}$ tension redressée (unidirectionnelle mais ondulée) $\xrightarrow{\text{filtre}}$ \SI{5}{V} continu, conforme à l'oscillogramme 2. Le schéma complet est donné dans la figure ci-dessous.
\end{exoresolu}

\begin{popfigure}
\begin{center}
\begin{tikzpicture}[scale=0.95, every node/.style={font=\small}]
\draw[fill=popblueL,draw=popblue,rounded corners] (0,0) rectangle (2.4,1.2);
\node[align=center] at (1.2,0.6) {Secteur\\ \SI{220}{V} $\sim$};
\draw[->,thick,popdark] (2.4,0.6) -- (3.1,0.6);
\draw[fill=poporangeL,draw=poporange,rounded corners] (3.1,0) rectangle (5.5,1.2);
\node[align=center] at (4.3,0.6) {Transformateur\\ (abaisseur)};
\draw[->,thick,popdark] (5.5,0.6) -- (6.2,0.6) node[midway,above]{\footnotesize 5 V $\sim$};
\draw[fill=poppurpleL,draw=poppurple,rounded corners] (6.2,0) rectangle (8.4,1.2);
\node[align=center] at (7.3,0.6) {Redresseur\\ (pont de diodes)};
\draw[->,thick,popdark] (8.4,0.6) -- (9.1,0.6) node[midway,above]{\footnotesize redressée};
\draw[fill=popgreenL,draw=popgreen,rounded corners] (9.1,0) rectangle (11.1,1.2);
\node[align=center] at (10.1,0.6) {Filtre\\ (condensateur)};
\draw[->,thick,popdark] (11.1,0.6) -- (12.1,0.6);
\draw[fill=popgoldL,draw=popgold,rounded corners] (12.1,0) rectangle (14.3,1.2);
\node[align=center] at (13.2,0.6) {Sortie\\ \SI{5}{V} continu};
\end{tikzpicture}
\end{center}
\legende{Chaîne de conversion de l'adaptateur B : le transformateur abaisse la tension alternative de \SI{220}{V} à environ \SI{5}{V}, le redresseur la rend unidirectionnelle, le filtre la lisse pour obtenir une tension continue stable.}
\end{popfigure}

\begin{attention}[Pièges fréquents au BEPC]
\begin{itemize}
\item Ne confonds pas \textbf{tension} et \textbf{intensité} : la tension de sortie de l'adaptateur doit être \emph{égale} à celle de l'appareil ; en revanche, l'intensité maximale de l'adaptateur doit être \emph{supérieure ou égale} au courant demandé. Un adaptateur « 5 V, 2 A » peut alimenter un appareil qui ne consomme que \SI{0,5}{A} : c'est l'appareil qui fixe le courant qu'il tire.
\item Un transformateur seul ne fournit \textbf{jamais} de courant continu : si la sortie est continue, il y a forcément un étage de redressement (et de filtrage) après le transformateur.
\item Une tension alternative de valeur efficace \SI{5}{V} n'est pas équivalente à une tension continue de \SI{5}{V} pour un appareil prévu pour le continu : la nature du courant compte autant que la valeur.
\end{itemize}
\end{attention}

\courssub{Bilan de l'activité}

Cette activité t'a permis de mobiliser l'ensemble des savoir-faire de la leçon dans une situation authentique de la vie quotidienne camerounaise. Tu as mis en évidence les points essentiels suivants :

\begin{itemize}
\item La \cle{plaque signalétique} est la carte d'identité de l'adaptateur : elle renseigne la tension d'entrée, la tension de sortie, la nature du courant ($\sim$ ou continu) et l'intensité maximale. Sa lecture suffit souvent à deviner les fonctions internes : une sortie continue impose la présence d'un \cle{redressement} après le \cle{transformateur}.
\item L'\cle{oscillogramme} est la preuve expérimentale : une sinusoïde révèle une tension alternative, une droite horizontale une tension continue redressée et filtrée.
\item Le choix d'un adaptateur obéit à une règle stricte : \textbf{même tension, même nature de courant, intensité maximale suffisante}. Une tension trop élevée détruit l'appareil (cas de la radio grillée) ; une tension trop faible empêche son fonctionnement.
\item La sécurité électrique n'est pas optionnelle : vérification des plaques, état des câbles, non-surcharge des multiprises sont des réflexes qui protègent les personnes et le matériel.
\end{itemize}

\begin{aretenir}[L'essentiel de l'activité]
Un adaptateur secteur convertit le \SI{220}{V} alternatif du réseau en une basse tension adaptée à l'appareil : le \textbf{transformateur} abaisse la tension, le \textbf{redresseur} (pont de diodes) la rend unidirectionnelle, le \textbf{filtre} (condensateur) la rend continue. Avant tout branchement : lire la plaque signalétique et vérifier la compatibilité tension--courant--intensité.
\end{aretenir}

\newpage

\coursec{Méthodes et exercices résolus}

\courssub{Méthode 1 : Identifier la fonction d'un adaptateur dans une situation concrète}

\begin{definition}[Adaptateur secteur]
Un \cle{adaptateur secteur} est un appareil qui, branché sur la prise de courant ($\SI{220}{V}$ alternatif chez ENEO), fournit à un appareil électrique une \textbf{tension plus faible} et de \textbf{nature adaptée} (le plus souvent continue). Il remplit deux fonctions principales : la \cle{transformation} (abaissement de la tension) et le \cle{redressement} (conversion de l'alternatif en continu).
\end{definition}

\begin{methode}[Identifier et définir la fonction d'un adaptateur]
Pour répondre à une question du type « Quel est le rôle de l'adaptateur dans cette situation ? » :
\begin{enumerate}
\item \textbf{Repérer les deux tensions en présence} : la tension du secteur ENEO ($U_1 = \SI{220}{V}$ alternatif) et la tension de fonctionnement de l'appareil (par exemple $\SI{9}{V}$ continu, inscrite sur l'étiquette).
\item \textbf{Comparer} : si la tension de l'appareil est plus faible que celle du secteur, l'adaptateur joue le rôle d'\cle{abaisseur de tension}.
\item \textbf{Identifier la nature des courants} : le secteur fournit un courant \cle{alternatif} ; la plupart des appareils électroniques (radio, téléphone, ordinateur) fonctionnent en courant \cle{continu}. L'adaptateur assure donc aussi la conversion alternatif $\rightarrow$ continu.
\item \textbf{Rédiger la conclusion} en citant les deux fonctions : \emph{transformation} (abaissement de la tension par le transformateur) et \emph{redressement} (conversion de l'alternatif en continu).
\end{enumerate}
\textbf{Réflexe} : un adaptateur = un \cle{transformateur} + un \cle{redresseur} (+ un filtre). Toujours citer ces deux étages dans la réponse.
\end{methode}

\begin{exoresolu}[Le chargeur de téléphone de la maison]
Énoncé. À Douala, Amina branche le chargeur de son téléphone portable sur une prise ENEO qui délivre une tension alternative de $\SI{220}{V}$. Sur l'étiquette du téléphone, on lit : « Entrée : $5\ \text{V}$ continu — $1\ \text{A}$ ».
\begin{enumerate}
\item Pourquoi ne peut-on pas brancher le téléphone directement sur la prise ?
\item Quelles sont les deux fonctions assurées par le chargeur ?
\end{enumerate}
\tcblower
\textbf{Solution détaillée.}
\begin{enumerate}
\item Le téléphone fonctionne sous une tension \textbf{continue} de $\SI{5}{V}$, alors que la prise délivre une tension \textbf{alternative} de $\SI{220}{V}$. Or $220 \gg 5$ : brancher le téléphone directement sur la prise le soumettrait à une tension 44 fois trop grande, ce qui le détruirait immédiatement (et provoquerait un départ de feu). De plus, le courant alternatif du secteur ne convient pas aux circuits électroniques du téléphone, qui exigent un courant continu. Le branchement direct est donc \textbf{impossible et dangereux}.
\item Le chargeur est un \cle{adaptateur secteur}. Il assure deux fonctions :
\begin{itemize}
\item la \textbf{transformation} : un transformateur abaisse la tension alternative de $\SI{220}{V}$ à une faible tension alternative (voisine de $\SI{5}{V}$) ;
\item le \textbf{redressement} : un pont de diodes convertit cette tension alternative en tension continue, qu'un condensateur de filtrage rend presque constante.
\end{itemize}
\end{enumerate}
\textbf{Conclusion} : le chargeur transforme la tension alternative de $\SI{220}{V}$ du secteur en une tension continue de $\SI{5}{V}$ adaptée au téléphone, grâce à un transformateur (abaissement) puis à un redresseur (conversion alternatif $\rightarrow$ continu).
\end{exoresolu}

\courssub{Méthode 2 : Exploiter la formule du transformateur}

\begin{propriete}[Relation du transformateur parfait]
Pour un transformateur parfait alimenté en tension \textbf{alternative}, les tensions efficaces primaire $U_1$ et secondaire $U_2$ sont proportionnelles aux nombres de spires $N_1$ et $N_2$ des bobinages :
\[ \frac{U_2}{U_1} = \frac{N_2}{N_1} = m \]
où $m$ est le \cle{rapport de transformation} (grandeur sans unité). Si $m < 1$, le transformateur est un \cle{abaisseur} ; si $m > 1$, c'est un \cle{élévateur}. Cette relation n'est \textbf{pas exigible en démonstration} au BEPC, mais son exploitation numérique, si.
\end{propriete}

\begin{methode}[Calculer avec le rapport de transformation]
Pour un transformateur de tensions primaire $U_1$ (secteur) et secondaire $U_2$ (sortie), avec $N_1$ spires au primaire et $N_2$ spires au secondaire :
\begin{enumerate}
\item \textbf{Écrire la relation fondamentale} (transformateur parfait) :
\[ \frac{U_2}{U_1} = \frac{N_2}{N_1} = m \]
\item \textbf{Identifier le type} : si $U_2 < U_1$ (c'est-à-dire $N_2 < N_1$, donc $m < 1$), c'est un \cle{abaisseur} ; si $U_2 > U_1$, c'est un \cle{élévateur}. Dans un adaptateur secteur, c'est toujours un abaisseur.
\item \textbf{Isoler l'inconnue} : $U_2 = U_1 \times \dfrac{N_2}{N_1}$ ou $N_2 = N_1 \times \dfrac{U_2}{U_1}$.
\item \textbf{Vérifier la cohérence} : pour un adaptateur, le résultat doit être une tension faible (quelques volts à quelques dizaines de volts) et un nombre de spires entier.
\end{enumerate}
\textbf{Conditions d'application} : la formule n'est valable que pour un transformateur parfait, alimenté en tension \textbf{alternative} sinusoïdale. Elle est inutilisable en courant continu.
\end{methode}

\begin{exoresolu}[Le transformateur d'un adaptateur de radio]
Énoncé. Le transformateur d'un adaptateur secteur possède $N_1 = \num{1100}$ spires au primaire et $N_2 = 60$ spires au secondaire. Il est branché sur le secteur ENEO de tension efficace $U_1 = \SI{220}{V}$.
\begin{enumerate}
\item Calculer le rapport de transformation $m$ et préciser la nature de ce transformateur.
\item Calculer la tension efficace $U_2$ disponible au secondaire.
\item On veut obtenir $U'_2 = \SI{12}{V}$ au secondaire sans modifier le primaire. Combien de spires $N'_2$ faut-il bobiner ?
\end{enumerate}
\tcblower
\textbf{Solution détaillée.}
\begin{enumerate}
\item Le rapport de transformation est :
\[ m = \frac{N_2}{N_1} = \frac{60}{1100} \approx \num{0,055} \]
Comme $m < 1$ (on a $N_2 < N_1$), la tension secondaire sera plus faible que la tension primaire : c'est un transformateur \textbf{abaisseur}, ce qui est normal pour un adaptateur secteur.
\item D'après la relation du transformateur parfait :
\[ \frac{U_2}{U_1} = \frac{N_2}{N_1} \quad \Longrightarrow \quad U_2 = U_1 \times \frac{N_2}{N_1} = 220 \times \frac{60}{1100} = \frac{\num{13200}}{1100} = \SI{12}{V} \]
Le secondaire délivre une tension alternative efficace de $\SI{12}{V}$. Cette tension sera ensuite redressée pour alimenter la radio en continu.
\item On isole $N'_2$ dans la même relation :
\[ N'_2 = N_1 \times \frac{U'_2}{U_1} = 1100 \times \frac{12}{220} = 1100 \times \num{0,0545} = 60 \ \text{spires} \]
Il faut bobiner $60$ spires au secondaire : c'est cohérent avec la question 2, puisque l'adaptateur fournissait déjà $\SI{12}{V}$.
\end{enumerate}
\textbf{Conclusion} : $m \approx \num{0,055}$ (abaisseur), $U_2 = \SI{12}{V}$, et il faut $N'_2 = 60$ spires pour obtenir $\SI{12}{V}$.
\end{exoresolu}

\courssub{Méthode 3 : Expliquer le redressement (alternatif $\rightarrow$ continu)}

\begin{methode}[Décrire la fonction de redressement]
Pour expliquer comment l'adaptateur convertit le courant alternatif en courant continu :
\begin{enumerate}
\item \textbf{Rappeler les définitions} : un courant \cle{alternatif} change de sens périodiquement ($\num{50}$ fois par seconde au Cameroun, fréquence $f = \SI{50}{Hz}$) ; un courant \cle{continu} circule toujours dans le même sens.
\item \textbf{Citer le composant clé} : la \cle{diode}, qui ne laisse passer le courant que dans un seul sens. Un \cle{pont de diodes} (4 diodes) redresse les deux alternances : c'est le \emph{redressement}.
\item \textbf{Décrire le résultat} : après redressement, la tension garde le même signe mais varie encore (tension « redressée » ondulée).
\item \textbf{Ajouter le filtrage} : un \cle{condensateur} placé en sortie lisse la tension et la rend presque constante, proche d'une tension continue parfaite.
\item \textbf{Schématiser la chaîne} : Secteur $\SI{220}{V}$ alternatif $\rightarrow$ Transformateur $\rightarrow$ faible tension alternative $\rightarrow$ Pont de diodes $\rightarrow$ tension redressée $\rightarrow$ Condensateur $\rightarrow$ tension continue.
\end{enumerate}
\end{methode}

\begin{exoresolu}[Observer les tensions à l'oscilloscope]
Énoncé. Au laboratoire du collège, on visualise à l'oscilloscope la tension à différents points d'un adaptateur ouvert. On obtient trois oscillogrammes : (a) une sinusoïde de valeur efficace $\SI{220}{V}$ ; (b) une sinusoïde de valeur efficace $\SI{9}{V}$ ; (c) une tension de signe constant, presque stable autour de $\SI{9}{V}$.
\begin{enumerate}
\item Attribuer chaque oscillogramme à un point du circuit : entrée (prise), sortie du transformateur, sortie de l'adaptateur.
\item Quel organe permet de passer de (b) à (c) ? Expliquer son rôle en deux lignes.
\end{enumerate}
\tcblower
\textbf{Solution détaillée.}
\begin{enumerate}
\item \begin{itemize}
\item Oscillogramme (a) : sinusoïde de $\SI{220}{V}$ efficace $\rightarrow$ c'est la tension du \textbf{secteur ENEO}, à l'entrée de l'adaptateur (la prise).
\item Oscillogramme (b) : sinusoïde de $\SI{9}{V}$ efficace $\rightarrow$ c'est la tension à la \textbf{sortie du transformateur} : la tension a été abaissée mais elle est encore alternative.
\item Oscillogramme (c) : tension de signe constant, presque stable $\rightarrow$ c'est la tension à la \textbf{sortie de l'adaptateur}, après redressement et filtrage : elle est devenue continue.
\end{itemize}
\item L'organe qui fait passer de (b) à (c) est l'ensemble \textbf{pont de diodes + condensateur}. Le pont de diodes ne laisse passer le courant que dans un seul sens : la tension ne change plus de signe (redressement). Le condensateur se charge quand la tension monte et se décharge quand elle descend, ce qui lisse les variations : la tension devient presque constante (filtrage).
\end{enumerate}
\textbf{Conclusion} : le transformateur abaisse la tension (a $\rightarrow$ b), puis le pont de diodes et le condensateur la convertissent en tension continue (b $\rightarrow$ c).
\end{exoresolu}

\courssub{Méthode 4 : Lire et exploiter la plaque signalétique d'un adaptateur}

\begin{methode}[Interpréter les indications d'une étiquette d'adaptateur]
Sur toute plaque signalétique d'adaptateur, on lit deux lignes essentielles :
\begin{enumerate}
\item \textbf{INPUT (entrée)} : tension et nature du courant d'alimentation, ex. « $220\ \text{V}$ $\sim$ $50\ \text{Hz}$ ». Le symbole $\sim$ signifie \cle{alternatif}.
\item \textbf{OUTPUT (sortie)} : tension et courant maximal délivrés, ex. « $12\ \text{V}$ ⎓ $2\ \text{A}$ ». Le symbole ⎓ signifie \cle{continu}.
\item \textbf{Vérifier la compatibilité} avec l'appareil : la tension de sortie de l'adaptateur doit être \textbf{exactement égale} à la tension de l'appareil ; le courant maximal de l'adaptateur doit être \textbf{supérieur ou égal} au courant demandé par l'appareil.
\item \textbf{Calculer la puissance} si besoin : $P = U \times I$ (en watts, avec $U$ en volts et $I$ en ampères).
\end{enumerate}
\end{methode}

\begin{exoresolu}[Choisir le bon adaptateur pour un poste radio]
Énoncé. Un poste radio porte l'indication « $9\ \text{V}$ continu — $1{,}5\ \text{A}$ ». Au marché de Yaoundé, on propose trois adaptateurs :
\begin{itemize}
\item Adaptateur A : OUTPUT $12\ \text{V}$ ⎓ $2\ \text{A}$ ;
\item Adaptateur B : OUTPUT $9\ \text{V}$ ⎓ $1\ \text{A}$ ;
\item Adaptateur C : OUTPUT $9\ \text{V}$ ⎓ $2\ \text{A}$.
\end{itemize}
\begin{enumerate}
\item Quel adaptateur faut-il choisir ? Justifier en écartant les deux autres.
\item Calculer la puissance électrique consommée par le poste radio en fonctionnement normal.
\end{enumerate}
\tcblower
\textbf{Solution détaillée.}
\begin{enumerate}
\item \textbf{Critère de tension} : la tension de sortie doit être exactement $\SI{9}{V}$ continu.
\begin{itemize}
\item L'adaptateur A délivre $\SI{12}{V}$ : tension trop élevée, il risque de \textbf{griller} le poste radio. À écarter.
\item Les adaptateurs B et C délivrent bien $\SI{9}{V}$ continu : critère de tension satisfait.
\end{itemize}
\textbf{Critère de courant} : le poste demande $\SI{1,5}{A}$ ; l'adaptateur doit pouvoir fournir \emph{au moins} ce courant.
\begin{itemize}
\item L'adaptateur B ne fournit que $\SI{1}{A} < \SI{1,5}{A}$ : il sera surchargé, il chauffera et pourra être détruit. À écarter.
\item L'adaptateur C fournit jusqu'à $\SI{2}{A} > \SI{1,5}{A}$ : il convient. Le poste ne « prendra » que les $\SI{1,5}{A}$ dont il a besoin.
\end{itemize}
\textbf{Choix : l'adaptateur C.}
\item La puissance consommée par le poste radio est :
\[ P = U \times I = 9 \times 1{,}5 = \SI{13,5}{W} \]
\end{enumerate}
\textbf{Conclusion} : il faut choisir l'adaptateur C ($9\ \text{V}$ ⎓ $2\ \text{A}$) : bonne tension, courant suffisant. Le poste consomme alors une puissance de $\SI{13,5}{W}$.
\end{exoresolu}

\courssub{Méthode 5 : Appliquer les consignes de sécurité liées aux adaptateurs}

\begin{methode}[Raisonner sur la sécurité et le bon usage]
Face à une situation de vie courante (adaptateur qui chauffe, rallonge surchargée, adaptateur mouillé...) :
\begin{enumerate}
\item \textbf{Repérer le danger} : échauffement (risque d'incendie), contact avec l'eau (risque d'électrisation), tension inadaptée (destruction de l'appareil), câble dénudé (court-circuit).
\item \textbf{Citer la règle de sécurité correspondante} :
\begin{itemize}
\item ne jamais utiliser un adaptateur mouillé ou les mains mouillées ;
\item ne pas couvrir un adaptateur en fonctionnement (il doit évacuer la chaleur) ;
\item débrancher l'adaptateur après usage : il consomme et chauffe même à vide ;
\item ne jamais ouvrir un adaptateur branché : le primaire est relié au $\SI{220}{V}$, tension mortelle ;
\item utiliser un adaptateur de tension de sortie strictement conforme à l'appareil.
\end{itemize}
\item \textbf{Justifier physiquement} : le corps humain est conducteur ; une tension de $\SI{220}{V}$ peut provoquer une électrisation mortelle, alors que la sortie basse tension ($\leq \SI{24}{V}$) est sans danger.
\end{enumerate}
\end{methode}

\begin{exoresolu}[L'adaptateur sous l'oreiller]
Énoncé. Pendant la nuit, Junior laisse le chargeur de sa lampe rechargeable branché sur la prise, posé sous un oreiller, alors que la lampe est déjà débranchée du chargeur. Le matin, l'adaptateur est très chaud.
\begin{enumerate}
\item Expliquer pourquoi l'adaptateur a chauffé alors qu'il n'alimentait plus la lampe.
\item Citer deux risques de cette pratique et deux règles de bon usage.
\end{enumerate}
\tcblower
\textbf{Solution détaillée.}
\begin{enumerate}
\item Même à vide (sans appareil branché en sortie), le primaire du transformateur reste alimenté par le secteur $\SI{220}{V}$. Il est traversé par un petit courant à vide qui, par \textbf{effet Joule}, dégage de la chaleur dans les bobinages et le circuit magnétique. De plus, l'oreiller empêche l'adaptateur d'évacuer cette chaleur : sa température monte donc fortement.
\item \textbf{Risques} :
\begin{itemize}
\item la chaleur accumulée peut enflammer l'oreiller : \textbf{risque d'incendie} ;
\item l'échauffement détruit l'isolation interne : \textbf{court-circuit} possible et destruction de l'adaptateur (et gaspillage d'énergie ENEO).
\end{itemize}
\textbf{Règles de bon usage} :
\begin{itemize}
\item \textbf{débrancher} l'adaptateur de la prise dès que la charge est terminée ;
\item laisser l'adaptateur \textbf{à l'air libre}, sur une surface dure, jamais sous un tissu ou un oreiller.
\end{itemize}
\end{enumerate}
\textbf{Conclusion} : un adaptateur branché consomme et chauffe même à vide ; il faut le débrancher après usage et ne jamais le couvrir pendant son fonctionnement.
\end{exoresolu}

\courssub{Méthode 6 : Tracer et interpréter la chaîne de conversion d'un adaptateur}

\begin{methode}[Construire la chaîne de conversion]
Pour schématiser ou compléter la chaîne de conversion d'un adaptateur secteur :
\begin{enumerate}
\item \textbf{Ordonner les quatre étages} dans le sens du courant : Prise secteur ($\SI{220}{V}$ alternatif) $\rightarrow$ Transformateur (abaisseur) $\rightarrow$ Pont de diodes (redresseur) $\rightarrow$ Condensateur (filtre) $\rightarrow$ Appareil (tension continue).
\item \textbf{Indiquer sous chaque flèche la nature et la valeur de la tension} : alternative $\SI{220}{V}$, puis alternative faible (ex. $\SI{12}{V}$), puis redressée ondulée, puis continue.
\item \textbf{Associer chaque étage à sa fonction} : transformation (transformateur), redressement (diodes), filtrage (condensateur).
\item \textbf{Vérifier} que la tension finale correspond à l'indication OUTPUT de la plaque signalétique.
\end{enumerate}
\end{methode}

\begin{popfigure}
\begin{center}
\begin{tikzpicture}[font=\small, >=Stealth, node distance=0.55cm]
\node[draw=popblue, fill=popblueL, rounded corners, minimum width=2.1cm, minimum height=1cm, align=center] (secteur) {Secteur ENEO\\ $220\ \text{V}$ $\sim$};
\node[draw=poporange, fill=poporangeL, rounded corners, minimum width=2.3cm, minimum height=1cm, align=center, right=of secteur] (transfo) {Transformateur\\ (abaisseur)};
\node[draw=poppurple, fill=poppurpleL, rounded corners, minimum width=2.3cm, minimum height=1cm, align=center, right=of transfo] (pont) {Pont de diodes\\ (redresseur)};
\node[draw=popgreen, fill=popgreenL, rounded corners, minimum width=2.3cm, minimum height=1cm, align=center, right=of pont] (condo) {Condensateur\\ (filtre)};
\node[draw=popdark, fill=popcream, rounded corners, minimum width=2.1cm, minimum height=1cm, align=center, right=of condo] (app) {Appareil\\ $12\ \text{V}$ ⎓};
\draw[->, thick, popdark] (secteur) -- (transfo) node[midway, above, font=\scriptsize] {$220\ \text{V}$ $\sim$};
\draw[->, thick, popdark] (transfo) -- (pont) node[midway, above, font=\scriptsize] {$12\ \text{V}$ $\sim$};
\draw[->, thick, popdark] (pont) -- (condo) node[midway, above, font=\scriptsize] {redressée};
\draw[->, thick, popdark] (condo) -- (app) node[midway, above, font=\scriptsize] {$12\ \text{V}$ ⎓};
\end{tikzpicture}
\end{center}
\legende{Chaîne de conversion d'un adaptateur secteur : le transformateur abaisse la tension (qui reste alternative), le pont de diodes la redresse, le condensateur la filtre pour obtenir une tension continue.}
\end{popfigure}

\begin{exoresolu}[Compléter la chaîne de conversion]
Énoncé. Voici la chaîne de conversion incomplète d'un adaptateur de sortie $6\ \text{V}$ continu :
\begin{center}
Secteur $220\ \text{V}$ $\sim$ $\longrightarrow$ \fbox{\phantom{XXXX}(1)\phantom{XXXX}} $\longrightarrow$ $6\ \text{V}$ $\sim$ $\longrightarrow$ \fbox{\phantom{XXXX}(2)\phantom{XXXX}} $\longrightarrow$ tension redressée $\longrightarrow$ \fbox{\phantom{XXXX}(3)\phantom{XXXX}} $\longrightarrow$ $6\ \text{V}$ ⎓
\end{center}
\begin{enumerate}
\item Identifier les étages (1), (2) et (3) et donner la fonction de chacun.
\item Le primaire du transformateur comporte $N_1 = \num{2200}$ spires. Calculer le nombre de spires $N_2$ du secondaire.
\end{enumerate}
\tcblower
\textbf{Solution détaillée.}
\begin{enumerate}
\item \begin{itemize}
\item \textbf{(1) Transformateur} : il assure la \textbf{transformation}, c'est-à-dire l'abaissement de la tension alternative de $\SI{220}{V}$ à $\SI{6}{V}$ (la tension reste alternative, d'où le symbole $\sim$).
\item \textbf{(2) Pont de diodes (redresseur)} : il assure le \textbf{redressement} : il force le courant à circuler dans un seul sens ; la tension ne change plus de signe mais reste ondulée.
\item \textbf{(3) Condensateur (filtre)} : il assure le \textbf{filtrage} : il lisse la tension redressée et la rend pratiquement constante, égale à $\SI{6}{V}$ continu (symbole ⎓).
\end{itemize}
\item D'après la relation du transformateur parfait :
\[ \frac{U_2}{U_1} = \frac{N_2}{N_1} \quad \Longrightarrow \quad N_2 = N_1 \times \frac{U_2}{U_1} = 2200 \times \frac{6}{220} = \frac{\num{13200}}{220} = 60 \ \text{spires} \]
Vérification : $N_2 = 60 < N_1 = \num{2200}$, donc $m = \dfrac{60}{2200} \approx \num{0,027} < 1$ : c'est bien un abaisseur, cohérent avec le rôle de l'adaptateur.
\end{enumerate}
\textbf{Conclusion} : (1) transformateur, (2) pont de diodes, (3) condensateur ; le secondaire comporte $N_2 = 60$ spires.
\end{exoresolu}

\begin{aretenir}[L'essentiel de la leçon]
\begin{itemize}
\item Un \cle{adaptateur secteur} convertit la tension alternative du secteur ($\SI{220}{V}$) en une tension faible, le plus souvent continue, adaptée à l'appareil.
\item Il remplit deux fonctions : la \cle{transformation} (transformateur abaisseur : $\dfrac{U_2}{U_1} = \dfrac{N_2}{N_1} = m$, avec $m < 1$) et le \cle{redressement} (pont de diodes), complété par le \cle{filtrage} (condensateur).
\item Le transformateur ne fonctionne qu'en courant \textbf{alternatif}.
\item Sécurité : ne jamais couvrir un adaptateur en fonctionnement, le débrancher après usage, ne jamais l'ouvrir branché, et choisir une tension de sortie strictement conforme à l'appareil.
\end{itemize}
\end{aretenir}

\begin{attention}[Les 5 erreurs qui coûtent des points au BEPC]
\begin{enumerate}
\item \textbf{Confondre transformation et redressement.} Le transformateur abaisse la tension mais elle reste \emph{alternative} ; c'est le pont de diodes qui la rend continue. Écrire « le transformateur convertit l'alternatif en continu » est une erreur grave.
\item \textbf{Appliquer la formule du transformateur en courant continu.} La relation $\dfrac{U_2}{U_1} = \dfrac{N_2}{N_1}$ n'est valable qu'en régime \emph{alternatif} : un transformateur branché sur une pile ne fonctionne pas.
\item \textbf{Inverser le rapport de transformation.} $m = \dfrac{N_2}{N_1} = \dfrac{U_2}{U_1}$ : le secondaire est au numérateur. Pour un adaptateur (abaisseur), on doit trouver $m < 1$ ; si $m > 1$, c'est le signe d'une inversion.
\item \textbf{Oublier les unités et la nature du courant dans la conclusion.} Une tension seule ne suffit pas : il faut préciser « $\SI{9}{V}$ \emph{continu} » ou « $\SI{220}{V}$ \emph{alternatif} ». Les symboles $\sim$ (alternatif) et ⎓ (continu) doivent être correctement lus sur la plaque signalétique.
\item \textbf{Choisir un adaptateur au seul critère de la tension.} Il faut aussi vérifier que le courant maximal de l'adaptateur est \emph{supérieur ou égal} à celui demandé par l'appareil, et justifier le choix en écartant explicitement les autres propositions.
\end{enumerate}
\end{attention}

\newpage

\coursec{Exercices}

\coursec{Leçon 15 : L'adaptateur secteur}

\courssub{Pourquoi un adaptateur ? Fonction générale}

\begin{savaistu}[Situation déclenchante : mon téléphone a faim !]
Tu rentres de l'école, ton téléphone affiche \SI{5}{\%} de batterie. Tu prends ton chargeur, tu le branches sur la prise murale ENEO et... ça charge. Mais attends : la prise ENEO délivre du \SI{220}{V} \emph{alternatif} (symbole $\sim$), alors que ta batterie a besoin d'environ \SI{5}{V} \emph{continu} (symbole $\text{---}$). Brancher le téléphone directement sur le secteur le détruirait instantanément (court-circuit, surchauffe, incendie). L'\textbf{adaptateur secteur} (ou chargeur) est l'interface indispensable qui rend cette cohabitation possible.
\end{savaistu}

\begin{definition}[Adaptateur secteur]
Un \textbf{adaptateur secteur} est un dispositif électronique qui convertit la tension alternative du réseau de distribution (ENEO au Cameroun : \SI{220}{V} $\sim$ \SI{50}{Hz}) en une tension continue de faible valeur adaptée aux besoins d'un appareil (téléphone, radio, routeur, lampe LED, etc.).
\end{definition}

\begin{aretenir}[Les deux fonctions fondamentales]
Un adaptateur secteur assure \textbf{deux fonctions successives} :
\begin{enumerate}
\item \textbf{La transformation (abaissement) :} passer de \SI{220}{V} $\sim$ à une tension alternative faible (ex: \SI{9}{V} $\sim$, \SI{12}{V} $\sim$) grâce à un \textbf{transformateur}.
\item \textbf{Le redressement :} transformer cette tension alternative faible en tension continue (ex: \SI{9}{V} $\text{---}$) grâce à un \textbf{redresseur} (pont de diodes).
\end{enumerate}
Parfois, un troisième étage (régulation/lissage) stabilise la tension, mais le programme de 3ème retient les deux fonctions principales.
\end{aretenir}

\courssub{Première fonction : la transformation (abaissement de tension)}

\begin{experience}[Mesure de la tension à l'entrée et à la sortie du transformateur]
\textbf{Matériel :} Un adaptateur secteur démonté (hors tension !) ou un transformateur de laboratoire, un voltmètre (multimètre) calibre alternatif \SI{10}{V} et \SI{300}{V}, secteur \SI{220}{V}.
\textbf{Protocole :}
\begin{enumerate}
\item Branche le primaire du transformateur sur le secteur (prudence !).
\item Mesure la tension aux bornes du primaire ($U_1$).
\item Mesure la tension aux bornes du secondaire ($U_2$), sans rien y brancher (vide).
\item Relevé : $U_1 = \SI{220}{V}$ $\sim$ ; $U_2 = \SI{9}{V}$ $\sim$ (valeur efficace).
\end{enumerate}
\textbf{Interprétation :} La tension a été abaissée. Le transformateur est un \textbf{transformateur abaisseur}.
\textbf{Conclusion :} Le transformateur ne fonctionne qu'en courant \emph{alternatif}. Il ne change pas la fréquence (\SI{50}{Hz}), il modifie l'amplitude de la tension.
\end{experience}

\begin{propriete}[Relation de transformation (exigible au BEPC)]
Pour un transformateur parfait (sans pertes) :
\[
\frac{U_2}{U_1} = \frac{N_2}{N_1}
\]
où :
\begin{itemize}
\item $U_1$ : tension primaire (entrée) en volts (V) ;
\item $U_2$ : tension secondaire (sortie) en volts (V) ;
\item $N_1$ : nombre de spires du primaire ;
\item $N_2$ : nombre de spires du secondaire.
\end{itemize}
Si $N_2 < N_1$ alors $U_2 < U_1$ : c'est un \textbf{transformateur abaisseur} (cas de l'adaptateur).
Si $N_2 > N_1$ alors $U_2 > U_1$ : c'est un \textbf{transformateur élévateur} (cas des lignes HT).
\end{propriete}

\begin{attention}[Pièges fréquents]
\begin{itemize}
\item \textbf{Unités :} $U_1$ et $U_2$ doivent être dans la \emph{même unité} (tous les deux en V). $N_1$ et $N_2$ sont des nombres purs (sans unité).
\item \textbf{Courant :} La relation ne s'applique qu'aux tensions \emph{alternatives} (valeurs efficaces ou maximales, mais cohérentes).
\item \textbf{Sens :} Ne pas inverser $U_1/U_2$ et $N_1/N_2$. Astuce : « \emph{Secondaire sur Primaire = Secondaire sur Primaire} ».
\end{itemize}
\end{attention}

\courssub{Deuxième fonction : le redressement (alternatif $\rightarrow$ continu)}

\begin{experience}[Visualisation du redressement à l'oscilloscope]
\textbf{Matériel :} Adaptateur secteur (ou montage transformateur + pont de diodes), oscilloscope.
\textbf{Protocole :}
\begin{enumerate}
\item Place la sonde à la sortie du transformateur (entrée du redresseur). Observe.
\item Place la sonde à la sortie de l'adaptateur (sortie du redresseur). Observe.
\end{enumerate}
\textbf{Observations :}
\begin{itemize}
\item \textbf{Entrée redresseur :} Courbe sinusoïdale centrée sur 0 V. La tension change de signe : c'est du \textbf{courant alternatif} ($\sim$).
\item \textbf{Sortie redresseur :} Courbe toujours positive (au-dessus de 0 V). Elle « vibre » mais ne descend pas en dessous de 0 : c'est du \textbf{courant continu pulsé} ($\text{---}$).
\end{itemize}
\textbf{Interprétation :} Le redresseur (composé de 4 diodes en \emph{pont de Graetz}) « redresse » les alternances négatives en positives. Le courant ne circule plus dans les deux sens, mais dans un seul sens.
\end{experience}

\begin{definition}[Redressement]
Le \textbf{redressement} est l'opération qui consiste à transformer une tension alternative en une tension continue (toujours de même signe). Il est réalisé par un \textbf{redresseur} (pont de diodes).
\end{definition}

\begin{aretenir}[Symboles sur la plaque signalétique]
\begin{center}
\begin{tabularx}{\linewidth}{|l|X|X|}
\hline
\rowcolor{popblueL} Symbole & Signification & Exemple \\
\hline
$\sim$ & Courant / Tension \textbf{Alternatif} & INPUT : \SI{220}{V} $\sim$ \\
\hline
$\text{---}$ (trait plein sur trait pointillé) & Courant / Tension \textbf{Continu} & OUTPUT : \SI{5}{V} $\text{---}$ \\
\hline
\end{tabularx}
\end{center}
\emph{Note :} Parfois, le symbole $\Box$ (carré) est utilisé dans certains manuels pour « continu pulsé non lissé », mais le symbole officiel est $\text{---}$.
\end{aretenir}

\courssub{Schéma de principe et chaîne de conversion}

\begin{methode}[Dessiner la chaîne de conversion d'un adaptateur]
\begin{enumerate}
\item Dessine 4 blocs reliés par des flèches : \textbf{Prise secteur} $\rightarrow$ \textbf{Transformateur} $\rightarrow$ \textbf{Redresseur} $\rightarrow$ \textbf{Appareil}.
\item Sous chaque flèche, indique : \textbf{valeur de la tension} + \textbf{nature} ($\sim$ ou $\text{---}$).
\item Exemple pour un chargeur \SI{5}{V} :
\begin{align*}
\text{Prise} &\xrightarrow{\SI{220}{V}\ \sim} \text{Transformateur} \xrightarrow{\SI{9}{V}\ \sim} \text{Redresseur} \xrightarrow{\SI{9}{V}\ \text{---}} \text{Téléphone (\SI{5}{V}\ \text{---})}
\end{align*}
(Note : la régulation finale \SI{9}{V}$\rightarrow$\SI{5}{V} n'est pas au programme mais existe dans les chargeurs modernes « switching »).
\end{enumerate}
\end{methode}

\begin{popfigure}
\begin{tikzpicture}[node distance=2.2cm, auto, block/.style={draw, rectangle, rounded corners, minimum height=1.6cm, minimum width=2.8cm, fill=popblueL, thick, align=center}, arrow/.style={->, thick, popdark, shorten >=2pt, shorten <=2pt}]
\node[block, align=center] (prise){Prise secteur\\ ENEO};
\node[block, right of=prise, align=center] (transfo){Transformateur\\ (abaissement)};
\node[block, right of=transfo, align=center] (redr){Redresseur\\ (pont de diodes)};
\node[block, right of=redr, align=center] (app){Appareil\\ (Téléphone, Radio...)};
\draw[arrow] (prise) -- node[above, align=center] {\SI{220}{V} $\sim$\\ \SI{50}{Hz}} (transfo);
\draw[arrow] (transfo) -- node[above, align=center] {Faible tension\\ $\sim$ (ex \SI{9}{V})} (redr);
\draw[arrow] (redr) -- node[above, align=center] {Même tension\\ $\text{---}$ (ex \SI{9}{V})} (app);
\end{tikzpicture}
\legende{Chaîne de conversion complète d'un adaptateur secteur linéaire classique.}
\end{popfigure}

\courssub{Sécurité et bon usage des adaptateurs}

\begin{attention}[Règles de sécurité vitales]
\begin{enumerate}
\item \textbf{Tension de sortie adaptée :} La tension de sortie de l'adaptateur (OUTPUT) doit être \textbf{exactement égale} à la tension nominale de l'appareil. Une tension plus élevée grille l'appareil ; une tension plus faible ne le fait pas fonctionner (ou mal).
\item \textbf{Intensité suffisante :} L'intensité maximale de l'adaptateur (ex: \SI{1}{A}, \SI{2}{A}) doit être \textbf{supérieure ou égale} à l'intensité demandée par l'appareil. L'adaptateur \emph{ne force pas} le courant ; il \emph{peut fournir jusqu'à} cette valeur. L'appareil « tire » ce dont il a besoin.
\item \textbf{État du boîtier et fils :} Aucun fil dénudé, boîtier fissuré, surchauffe anormale, bruit suspect (bourdonnement fort). Danger d'électrocution ou d'incendie.
\item \textbf{Eau et humidité :} Ne jamais manipuler un adaptateur branché avec les mains mouillées, ni l'utiliser dehors sous la pluie, ni le poser sur un sol mouillé. L'eau conduit le courant \SI{220}{V} $\rightarrow$ électrocution mortelle.
\item \textbf{Surcharge de la prise/multiprise :} Ne pas brancher trop d'appareils puissants (fer, frigo, bouilloire) + adaptateurs sur une même multiprise. Risque d'échauffement des fils, de déclenchement du disjoncteur, ou d'incendie.
\item \textbf{Débrancher :} Débranche l'adaptateur de la prise quand tu ne charges pas (économie d'énergie, usure, risque orage).
\end{enumerate}
\end{attention}

\begin{savaistu}[Contexte Cameroun : Adaptateurs et « courant ENEO »]
\begin{itemize}
\item Au Cameroun, le réseau ENEO est nominalement \SI{220}{V} $\sim$ \SI{50}{Hz}, mais les variations sont fréquentes (\SI{180}{V} à \SI{260}{V}). Les adaptateurs de qualité (marques connues) supportent souvent une plage d'entrée large (ex: \SI{100}{V}-\SI{240}{V} $\sim$), ce qui les protège.
\item Méfie-toi des adaptateurs « contrefaits » vendus à bas prix au marché (Limbé, Marché Central, etc.) : transformateur sous-dimensionné (surchauffe), isolation médiocre, pas de redressement correct (bourdonnement), tension de sortie instable. Ils détruisent les batteries des téléphones et risquent l'incendie.
\item Avec les délestages, beaucoup utilisent des groupes électrogènes. La tension d'un groupe peut être instable (pic au démarrage). Un adaptateur de qualité (ou un onduleur/régulateur en amont) protège tes appareils.
\end{itemize}
\end{savaistu}

\courssub{Je vérifie mes connaissances}

\begin{exercice}[QCM : les fonctions de l'adaptateur]
Pour chaque question, une seule réponse est exacte. Recopie la lettre correspondante.

\textbf{1.} Un adaptateur secteur sert à :
\begin{itemize}
\item a) augmenter la tension du réseau ENEO ;
\item b) adapter la tension du secteur aux besoins d'un appareil ;
\item c) stocker l'énergie électrique ;
\item d) mesurer la tension du secteur.
\end{itemize}

\textbf{2.} La tension délivrée par une prise ENEO au Cameroun est :
\begin{itemize}
\item a) continue et de valeur \SI{220}{V} ;
\item b) alternative et de valeur \SI{220}{V} ;
\item c) continue et de valeur \SI{12}{V} ;
\item d) alternative et de valeur \SI{6}{V}.
\end{itemize}

\textbf{3.} La fonction qui abaisse la tension de \SI{220}{V} à \SI{9}{V} est :
\begin{itemize}
\item a) le redressement ;
\item b) la transformation ;
\item c) la régulation ;
\item d) la filtration.
\end{itemize}

\textbf{4.} Le redressement convertit une tension :
\begin{itemize}
\item a) continue en tension alternative ;
\item b) élevée en tension faible ;
\item c) alternative en tension continue ;
\item d) faible en tension élevée.
\end{itemize}
\end{exercice}

\begin{exercice}[Vrai ou faux ? Justifie]
Réponds par vrai ou faux, puis justifie ta réponse en une ou deux phrases.

\textbf{1.} Un téléphone portable peut être branché directement sur le secteur \SI{220}{V}.

\textbf{2.} Le transformateur d'un adaptateur fonctionne avec une tension alternative à son entrée.

\textbf{3.} À la sortie d'un adaptateur, la tension est alternative.

\textbf{4.} Le symbole $\sim$ sur la plaque d'un adaptateur désigne un courant continu.
\end{exercice}

\begin{exercice}[Définitions]
Définis en une phrase chacun des termes suivants, puis donne un exemple concret de la vie courante :
\begin{enumerate}
\item adaptateur secteur ;
\item transformation (abaissement de tension) ;
\item redressement.
\end{enumerate}
\end{exercice}

\begin{exercice}[Calcul direct sur un transformateur]
Le transformateur d'un adaptateur possède $N_1 = 880$ spires au primaire et $N_2 = 36$ spires au secondaire. Il est alimenté sous une tension primaire $U_1 = \SI{220}{V}$.

On rappelle la relation du transformateur : $\dfrac{U_2}{U_1} = \dfrac{N_2}{N_1}$.

\begin{enumerate}
\item Calcule la tension secondaire $U_2$.
\item Ce transformateur est-il abaisseur ou élévateur de tension ? Justifie.
\end{enumerate}
\end{exercice}

\courssub{Je m'entraîne}

\begin{exercice}[Lecture d'une plaque signalétique]
Voici la plaque signalétique de l'adaptateur d'un poste radio :

\begin{center}
\begin{tabularx}{\linewidth}{|l|X|}
\hline
\rowcolor{popblueL} INPUT (entrée) & \SI{220}{V} $\sim$ \SI{50}{Hz} \\
\hline
OUTPUT (sortie) & \SI{6}{V} $\text{---}$ \SI{800}{mA} \\
\hline
\end{tabularx}
\end{center}

\begin{enumerate}
\item Recopie et complète le tableau suivant :
\begin{center}
\begin{tabularx}{\linewidth}{|X|X|X|}
\hline
\rowcolor{popblueL} & Valeur de la tension & Nature du courant \\
\hline
Entrée de l'adaptateur & \ldots & \ldots \\
\hline
Sortie de l'adaptateur & \ldots & \ldots \\
\hline
\end{tabularx}
\end{center}
\item Nomme les deux fonctions réalisées par cet adaptateur et précise le rôle de chacune.
\item Calcule l'intensité maximale en ampères que peut débiter cet adaptateur.
\end{enumerate}
\end{exercice}

\begin{exercice}[Attribution d'oscillogrammes]
On a relevé trois oscillogrammes aux différentes étapes de la chaîne de conversion d'un adaptateur :

\begin{itemize}
\item Oscillogramme A : tension sinusoïdale de valeur maximale environ \SI{311}{V} (valeur efficace \SI{220}{V}) ;
\item Oscillogramme B : tension sinusoïdale de valeur efficace \SI{9}{V} ;
\item Oscillogramme C : tension de valeur \SI{9}{V}, ne changeant pas de signe (toujours positive).
\end{itemize}

\begin{enumerate}
\item Attribue chaque oscillogramme à l'une des étapes suivantes : tension du secteur, sortie du transformateur, sortie de l'adaptateur.
\item Justifie chaque attribution en une phrase.
\item Quelle étape correspond au redressement ? Explique ce qui a changé entre l'entrée et la sortie de cette étape.
\end{enumerate}
\end{exercice}

\begin{exercice}[Schéma de la chaîne de conversion]
Schématise la chaîne complète allant de la prise secteur jusqu'au téléphone portable en charge. Ton schéma comportera des blocs reliés par des flèches. Pour chaque bloc, indique :
\begin{enumerate}
\item le nom de l'élément (prise secteur, transformateur, dispositif de redressement, téléphone) ;
\item la nature de la tension (alternative ou continue) à l'entrée et à la sortie de chaque bloc ;
\item la valeur de la tension à chaque étape si l'adaptateur délivre \SI{5}{V} à partir du secteur \SI{220}{V}.
\end{enumerate}
\end{exercice}

\begin{exercice}[Sécurité : analyse de situations]
Pour chacune des quatre situations suivantes, identifie le danger encouru et propose la conduite correcte à adopter.

\begin{enumerate}
\item Mballa utilise un adaptateur dont le boîtier en plastique est fissuré, laissant apparaître les fils internes.
\item Le chargeur \SI{5}{V} d'Aïcha est perdu ; son voisin lui prête un adaptateur de sortie \SI{12}{V} pour recharger son téléphone.
\item Pendant la saison des pluies, Kamdem recharge sa lampe dehors, sous la pluie, avec l'adaptateur posé sur le sol mouillé.
\item Chez les Ngo, une seule multiprise alimente à la fois un téléviseur, un réfrigérateur, un fer à repasser et trois adaptateurs de téléphones.
\end{enumerate}
\end{exercice}

\courssub{Je me prépare au Bac}

\begin{exercice}[L'adaptateur du groupe électrogène]
Dans un quartier de Douala où les délestages sont fréquents, M. Etoundi utilise un adaptateur pour recharger la batterie de sa torche. La plaque de l'adaptateur indique : INPUT \SI{220}{V} $\sim$ \SI{50}{Hz} ; OUTPUT \SI{12}{V} $\text{---}$ \SI{1}{A}. Le transformateur interne possède $N_1 = 1100$ spires au primaire.

On rappelle : $\dfrac{U_2}{U_1} = \dfrac{N_2}{N_1}$.

\begin{enumerate}
\item Précise la nature (alternative ou continue) et la valeur de la tension :
\begin{enumerate}
\item à l'entrée de l'adaptateur ;
\item à la sortie de l'adaptateur.
\end{enumerate}
\item Nomme les deux fonctions assurées par l'adaptateur et décris le rôle de chacune.
\item Calcule le nombre de spires $N_2$ du secondaire du transformateur.
\item Calcule la puissance électrique maximale $P = U \times I$ que peut fournir cet adaptateur à la torche.
\item Un jour, M. Etoundi branche par erreur la torche directement sur le secteur \SI{220}{V}. Explique ce qui risque de se produire et pourquoi l'adaptateur est indispensable.
\end{enumerate}
\end{exercice}

\begin{exercice}[Choix d'un adaptateur au marché de Limbé]
Au marché de Limbé, une vendeuse propose trois adaptateurs pour un haut-parleur portatif dont la notice indique : \SI{5}{V} $\text{---}$ \SI{1,5}{A}.

\begin{center}
\begin{tabularx}{\linewidth}{|l|X|X|}
\hline
\rowcolor{popblueL} Adaptateur & Tension de sortie & Intensité maximale \\
\hline
Adaptateur 1 & \SI{5}{V} $\text{---}$ & \SI{1}{A} \\
\hline
Adaptateur 2 & \SI{5}{V} $\text{---}$ & \SI{2}{A} \\
\hline
Adaptateur 3 & \SI{9}{V} $\text{---}$ & \SI{1}{A} \\
\hline
\end{tabularx}
\end{center}

\begin{enumerate}
\item Rappelle la signification des deux grandeurs indiquées sur la sortie d'un adaptateur.
\item Explique pourquoi l'adaptateur 3 est dangereux pour le haut-parleur.
\item Explique pourquoi l'adaptateur 1 ne convient pas, bien que sa tension soit correcte.
\item Quel adaptateur faut-il choisir ? Justifie ta réponse en comparant les trois grandeurs.
\item La vendeuse affirme : « Un adaptateur qui peut donner \SI{2}{A} va forcer \SI{2}{A} dans l'appareil et l'abîmer. » Cette affirmation est-elle correcte ? Explique ce que l'intensité maximale signifie réellement.
\end{enumerate}
\end{exercice}

\begin{exercice}[Étude complète d'une chaîne de conversion]
Un élève de 3\textsuperscript{e} du collège de Bafoussam démonte (hors tension !) un vieil adaptateur de radio-cassette et identifie ses étages. Il relève les tensions suivantes avec un voltmètre :

\begin{itemize}
\item entre la prise et le premier étage : tension alternative \SI{220}{V} ;
\item entre le premier et le deuxième étage : tension alternative \SI{9}{V} ;
\item à la sortie du deuxième étage : tension continue \SI{9}{V}.
\end{itemize}

\begin{enumerate}
\item Nomme les deux étages de cet adaptateur et associe chacun à la fonction correspondante (transformation, redressement).
\item Dessine le schéma de principe de la chaîne complète secteur $\rightarrow$ radio, en précisant sous chaque flèche la valeur et la nature de la tension.
\item Le transformateur possède $N_2 = 45$ spires au secondaire. Calcule le nombre de spires $N_1$ du primaire (on donne $\dfrac{U_2}{U_1} = \dfrac{N_2}{N_1}$ ; arrondis à l'unité).
\item L'élève veut vérifier la nature des tensions avec un oscilloscope. Décris l'allure de la courbe (oscillogramme) qu'il observerait :
\begin{enumerate}
\item à l'entrée de l'adaptateur ;
\item à la sortie de l'adaptateur.
\end{enumerate}
\item Cite deux règles de sécurité que l'élève doit respecter lors de ses mesures sur le secteur \SI{220}{V}, et explique pourquoi chacune est nécessaire.
\end{enumerate}
\end{exercice}

\newpage

\coursec{Corrigés des exercices}

\begin{corrige}[Exercice 1 — QCM : les fonctions de l'adaptateur]
\textbf{1. Réponse b.} L'adaptateur secteur \emph{adapte} la tension du réseau (\SI{220}{V} alternatif) aux besoins de l'appareil (par exemple \SI{5}{V} continu). Il n'augmente pas la tension (a), ne stocke pas d'énergie (c), et ne mesure rien (d).

\textbf{2. Réponse b.} Au Cameroun, le réseau ENEO délivre une tension \textbf{alternative} de valeur efficace \SI{220}{V} et de fréquence \SI{50}{Hz}. Le symbole $\sim$ sur les plaques le confirme.

\textbf{3. Réponse b.} C'est la \textbf{transformation}, réalisée par le transformateur, qui abaisse la tension de \SI{220}{V} à une faible tension alternative. Le redressement (a) ne change pas la valeur de la tension, il change sa nature (alternatif $\rightarrow$ continu).

\textbf{4. Réponse c.} Le redressement convertit une tension \textbf{alternative} en tension \textbf{continue} (toujours de même signe), grâce au pont de diodes.
\end{corrige}

\begin{corrige}[Exercice 2 — Vrai ou faux ? Justifie]
\textbf{1. FAUX.} Un téléphone portable fonctionne sous une faible tension \emph{continue} (environ \SI{5}{V}). Branché directement sur le secteur \SI{220}{V} alternatif, il serait détruit instantanément (surtension, surchauffe, risque d'incendie). L'adaptateur est indispensable.

\textbf{2. VRAI.} Le transformateur ne fonctionne \emph{qu'en courant alternatif} : c'est la variation du courant au primaire qui crée, par induction, une tension au secondaire. Son entrée est donc toujours alimentée en tension alternative (\SI{220}{V} $\sim$).

\textbf{3. FAUX.} À la sortie d'un adaptateur, la tension est \textbf{continue} (symbole $\text{---}$) : c'est précisément le rôle du redresseur, deuxième étage de l'adaptateur, de convertir l'alternatif en continu.

\textbf{4. FAUX.} Le symbole $\sim$ désigne un courant (ou une tension) \textbf{alternatif}. Le courant continu est symbolisé par $\text{---}$ (trait plein sur trait pointillé).
\end{corrige}

\begin{corrige}[Exercice 3 — Définitions]
\textbf{1. Adaptateur secteur :} dispositif électronique qui convertit la tension alternative du réseau (\SI{220}{V} $\sim$) en une tension continue de faible valeur adaptée à un appareil. \emph{Exemple :} le chargeur d'un téléphone portable.

\textbf{2. Transformation (abaissement de tension) :} fonction qui consiste à modifier (ici diminuer) la valeur d'une tension alternative à l'aide d'un transformateur, sans changer sa nature ni sa fréquence. \emph{Exemple :} passage de \SI{220}{V} $\sim$ à \SI{9}{V} $\sim$ dans un chargeur de radio.

\textbf{3. Redressement :} fonction qui consiste à convertir une tension alternative en tension continue (toujours de même signe) à l'aide d'un redresseur (pont de diodes). \emph{Exemple :} passage de \SI{9}{V} $\sim$ à \SI{9}{V} $\text{---}$ à la sortie du chargeur.
\end{corrige}

\begin{corrige}[Exercice 4 — Calcul direct sur un transformateur]
\textbf{1. Calcul de $U_2$.}

On utilise la relation du transformateur : $\dfrac{U_2}{U_1} = \dfrac{N_2}{N_1}$, d'où :
\[
U_2 = U_1 \times \frac{N_2}{N_1} = 220 \times \frac{36}{880} = 220 \times 0{,}0409 = \SI{9}{V}.
\]
La tension secondaire vaut $U_2 = \SI{9}{V}$ (alternative).

\textbf{2. Nature du transformateur.}

On compare les nombres de spires : $N_2 = 36 < N_1 = 880$. Comme $N_2 < N_1$, on a $U_2 < U_1$ (\SI{9}{V} $<$ \SI{220}{V}) : le transformateur est un \textbf{abaisseur de tension}, ce qui est normal pour un adaptateur secteur.

\textbf{Point de vigilance :} vérifier que $U_1$ et $U_2$ sont dans la même unité (ici les deux en volts) et ne pas inverser le rapport des spires.
\end{corrige}

\begin{corrige}[Exercice 5 — Lecture d'une plaque signalétique]
\textbf{1. Tableau complété :}
\begin{center}
\begin{tabularx}{\linewidth}{|X|X|X|}
\hline
\rowcolor{popblueL} & Valeur de la tension & Nature du courant \\
\hline
Entrée de l'adaptateur & \SI{220}{V} & alternatif ($\sim$) \\
\hline
Sortie de l'adaptateur & \SI{6}{V} & continu ($\text{---}$) \\
\hline
\end{tabularx}
\end{center}

\textbf{2. Les deux fonctions :}
\begin{itemize}
\item \textbf{La transformation} (transformateur) : elle abaisse la tension alternative de \SI{220}{V} $\sim$ à une faible tension alternative (environ \SI{6}{V} $\sim$).
\item \textbf{Le redressement} (pont de diodes) : il convertit cette faible tension alternative en tension continue \SI{6}{V} $\text{---}$.
\end{itemize}

\textbf{3. Intensité maximale :} la plaque indique \SI{800}{mA}. Conversion :
\[
I_{\max} = \SI{800}{mA} = \frac{800}{1000} = \SI{0,8}{A}.
\]
L'adaptateur peut débiter au maximum \SI{0,8}{A}.
\end{corrige}

\begin{corrige}[Exercice 6 — Attribution d'oscillogrammes]
\textbf{1. et 2. Attributions et justifications :}
\begin{itemize}
\item \textbf{Oscillogramme A $\rightarrow$ tension du secteur} (entrée de l'adaptateur) : c'est une tension \emph{alternative} (sinusoïdale) de valeur efficace \SI{220}{V}, caractéristique du réseau ENEO.
\item \textbf{Oscillogramme B $\rightarrow$ sortie du transformateur} : la tension est encore \emph{alternative} (sinusoïdale) mais sa valeur efficace n'est plus que \SI{9}{V} : elle a été abaissée, sans changement de nature.
\item \textbf{Oscillogramme C $\rightarrow$ sortie de l'adaptateur} : la tension ne change \emph{plus de signe} (toujours positive) : c'est une tension \emph{continue}, obtenue après redressement.
\end{itemize}

\textbf{3. Étape du redressement :} c'est le passage de B à C (entre la sortie du transformateur et la sortie de l'adaptateur). À l'entrée du redresseur, la tension alternative change de signe périodiquement (positive puis négative) ; à la sortie, les alternances négatives ont été « redressées » : la tension reste toujours positive. La valeur (\SI{9}{V}) est conservée, seule la \emph{nature} de la tension a changé.
\end{corrige}

\begin{corrige}[Exercice 7 — Schéma de la chaîne de conversion]
La chaîne comporte quatre blocs reliés par des flèches, avec la nature et la valeur de la tension indiquées sous chaque flèche :

\begin{popfigure}
\begin{tikzpicture}[node distance=2.1cm, auto, block/.style={draw, rectangle, rounded corners, minimum height=1.5cm, minimum width=2.6cm, fill=popblueL, thick, align=center}, arrow/.style={->, thick, popdark, shorten >=2pt, shorten <=2pt}]
\node[block, align=center] (prise){Prise secteur\\ ENEO};
\node[block, right of=prise, align=center] (transfo){Transformateur\\ (abaissement)};
\node[block, right of=transfo, align=center] (redr){Redresseur\\ (pont de diodes)};
\node[block, right of=redr, align=center] (app){Téléphone\\ portable};
\draw[arrow] (prise) -- node[above, align=center] {\SI{220}{V} $\sim$\\ alternatif} (transfo);
\draw[arrow] (transfo) -- node[above, align=center] {\SI{5}{V} $\sim$\\ alternatif} (redr);
\draw[arrow] (redr) -- node[above, align=center] {\SI{5}{V} $\text{---}$\\ continu} (app);
\end{tikzpicture}
\legende{Chaîne de conversion attendue : les deux flèches internes portent la valeur et la nature de la tension.}
\end{popfigure}

\textbf{Critères de réussite du schéma :}
\begin{enumerate}
\item Les 4 blocs sont nommés dans l'ordre : prise secteur $\rightarrow$ transformateur $\rightarrow$ redresseur $\rightarrow$ téléphone.
\item La tension est \emph{alternative} à l'entrée et à la sortie du transformateur (\SI{220}{V} $\sim$ puis \SI{5}{V} $\sim$), et \emph{continue} à la sortie du redresseur (\SI{5}{V} $\text{---}$).
\item Les valeurs sont cohérentes : abaissement de \SI{220}{V} à \SI{5}{V} par le transformateur, conservation de la valeur \SI{5}{V} par le redresseur.
\end{enumerate}
\end{corrige}

\begin{corrige}[Exercice 8 — Sécurité : analyse de situations]
\textbf{1. Boîtier fissuré (Mballa).} \emph{Danger :} les fils internes sous tension \SI{220}{V} sont accessibles : risque d'\textbf{électrocution} au toucher et de court-circuit. \emph{Conduite :} cesser immédiatement l'utilisation, débrancher l'adaptateur et le remplacer (ou le faire réparer par un professionnel). Ne jamais le réparer soi-même avec du ruban adhésif.

\textbf{2. Adaptateur de \SI{12}{V} au lieu de \SI{5}{V} (Aïcha).} \emph{Danger :} la tension de sortie (\SI{12}{V}) est \textbf{supérieure} à la tension nominale du téléphone (\SI{5}{V}) : risque de destruction des circuits internes et de la batterie (surchauffe, voire incendie). \emph{Conduite :} refuser l'adaptateur et n'utiliser qu'un adaptateur délivrant \emph{exactement} \SI{5}{V} $\text{---}$.

\textbf{3. Recharge sous la pluie (Kamdem).} \emph{Danger :} l'eau est conductrice ; l'adaptateur posé sur le sol mouillé peut laisser passer le courant \SI{220}{V} : risque d'\textbf{électrocution mortelle} et de court-circuit. \emph{Conduite :} recharger à l'intérieur, dans un endroit sec, l'adaptateur posé sur une surface sèche, et ne jamais manipuler la prise avec les mains mouillées.

\textbf{4. Multiprise surchargée (famille Ngo).} \emph{Danger :} le téléviseur, le réfrigérateur et surtout le fer à repasser sont de gros consommateurs ; l'intensité totale dans la multiprise peut dépasser sa limite : \textbf{échauffement des fils, fonte de l'isolant, incendie} (ou déclenchement du disjoncteur). \emph{Conduite :} répartir les appareils puissants sur des prises murales différentes et réserver la multiprise aux appareils de faible puissance (adaptateurs, lampe).
\end{corrige}

\begin{corrige}[Exercice 9 — L'adaptateur du groupe électrogène]
\textbf{1. Nature et valeur des tensions :}
\begin{itemize}
\item a) À l'entrée : tension \textbf{alternative} de valeur \SI{220}{V} (fréquence \SI{50}{Hz}).
\item b) À la sortie : tension \textbf{continue} de valeur \SI{12}{V}.
\end{itemize}

\textbf{2. Les deux fonctions :}
\begin{itemize}
\item \textbf{Transformation} (transformateur) : abaisse la tension alternative de \SI{220}{V} $\sim$ à environ \SI{12}{V} $\sim$.
\item \textbf{Redressement} (pont de diodes) : convertit la tension alternative \SI{12}{V} $\sim$ en tension continue \SI{12}{V} $\text{---}$.
\end{itemize}

\textbf{3. Nombre de spires $N_2$ :}

De $\dfrac{U_2}{U_1} = \dfrac{N_2}{N_1}$, on tire :
\[
N_2 = N_1 \times \frac{U_2}{U_1} = 1100 \times \frac{12}{220} = 1100 \times 0{,}0545 = 60 \text{ spires}.
\]

\textbf{4. Puissance maximale :}
\[
P = U \times I = 12 \times 1 = \SI{12}{W}.
\]
L'adaptateur peut fournir au maximum \SI{12}{W} à la torche.

\textbf{5. Branchement direct sur le secteur :} la torche est conçue pour \SI{12}{V} continu ; soumise à \SI{220}{V} alternatif, soit une tension environ 18 fois trop élevée, ses circuits seraient traversés par un courant très intense : \textbf{destruction immédiate de l'ampoule ou des LED, surchauffe, risque d'incendie}. L'adaptateur est indispensable car il abaisse la tension à la valeur nominale de la torche \emph{et} la rend continue.

\medskip
\textbf{Barème indicatif (sur 10) :} Q1 : 2 pts (1 pt par tension, nature + valeur) ; Q2 : 2 pts (fonctions nommées et rôles décrits) ; Q3 : 2 pts (formule littérale 1 pt, application numérique 1 pt) ; Q4 : 2 pts (formule 1 pt, résultat avec unité 1 pt) ; Q5 : 2 pts.

\textbf{Points de vigilance :} citer la relation $\dfrac{U_2}{U_1} = \dfrac{N_2}{N_1}$ avant de l'appliquer ; isoler $N_2$ littéralement avant le calcul numérique ; ne pas oublier l'unité (watt) pour la puissance ; dans la question 5, évoquer à la fois la surtension \emph{et} la nature inadaptée (alternatif au lieu de continu).
\end{corrige}

\begin{corrige}[Exercice 10 — Choix d'un adaptateur au marché de Limbé]
\textbf{1. Signification des grandeurs :} la \textbf{tension de sortie} (en V) est la tension continue que l'adaptateur délivre à l'appareil ; l'\textbf{intensité maximale} (en A) est le courant maximal que l'adaptateur \emph{peut fournir} sans être endommagé.

\textbf{2. Adaptateur 3 dangereux :} sa tension de sortie (\SI{9}{V}) est \textbf{supérieure} à la tension nominale du haut-parleur (\SI{5}{V}). Cette surtension va griller les circuits internes de l'appareil (surchauffe, destruction des composants).

\textbf{3. Adaptateur 1 inadapté :} sa tension (\SI{5}{V}) est correcte, mais son intensité maximale (\SI{1}{A}) est \textbf{inférieure} à l'intensité demandée par le haut-parleur (\SI{1,5}{A}). L'adaptateur sera surchargé : il va chauffer, se dégrader, et l'appareil fonctionnera mal ou pas du tout.

\textbf{4. Choix :} il faut choisir l'\textbf{adaptateur 2} : sa tension de sortie est \emph{exactement égale} à la tension nominale (\SI{5}{V} = \SI{5}{V}) et son intensité maximale (\SI{2}{A}) est \emph{supérieure} à l'intensité demandée (\SI{1,5}{A}). C'est le seul qui respecte les deux conditions.

\textbf{5. Affirmation de la vendeuse : FAUX.} L'intensité maximale n'est pas un courant « forcé » dans l'appareil : c'est le courant \emph{maximal disponible}. C'est l'appareil qui « tire » le courant dont il a besoin (ici \SI{1,5}{A}). Un adaptateur de \SI{2}{A} alimentera le haut-parleur avec \SI{1,5}{A} seulement, sans aucun danger.

\medskip
\textbf{Barème indicatif (sur 10) :} Q1 : 2 pts ; Q2 : 2 pts ; Q3 : 2 pts ; Q4 : 2 pts (choix 1 pt + double justification 1 pt) ; Q5 : 2 pts.

\textbf{Points de vigilance :} les deux conditions (égalité des tensions, intensité maximale suffisante) doivent être citées \emph{séparément} ; pour la question 5, la justification attendue est « c'est l'appareil qui impose le courant qu'il consomme ».
\end{corrige}

\begin{corrige}[Exercice 11 — Étude complète d'une chaîne de conversion]
\textbf{1. Les deux étages :}
\begin{itemize}
\item \textbf{Premier étage : le transformateur} $\rightarrow$ fonction \textbf{transformation} : il abaisse la tension alternative de \SI{220}{V} $\sim$ à \SI{9}{V} $\sim$.
\item \textbf{Deuxième étage : le redresseur} (pont de diodes) $\rightarrow$ fonction \textbf{redressement} : il convertit la tension alternative \SI{9}{V} $\sim$ en tension continue \SI{9}{V} $\text{---}$.
\end{itemize}

\textbf{2. Schéma de principe :}

\begin{popfigure}
\begin{tikzpicture}[node distance=2.1cm, auto, block/.style={draw, rectangle, rounded corners, minimum height=1.5cm, minimum width=2.6cm, fill=popblueL, thick, align=center}, arrow/.style={->, thick, popdark, shorten >=2pt, shorten <=2pt}]
\node[block] (prise){Prise secteur};
\node[block, right of=prise] (transfo){Transformateur};
\node[block, right of=transfo] (redr){Redresseur};
\node[block, right of=redr] (app){Radio};
\draw[arrow] (prise) -- node[above, align=center] {\SI{220}{V} $\sim$} (transfo);
\draw[arrow] (transfo) -- node[above, align=center] {\SI{9}{V} $\sim$} (redr);
\draw[arrow] (redr) -- node[above, align=center] {\SI{9}{V} $\text{---}$} (app);
\end{tikzpicture}
\legende{Chaîne secteur $\rightarrow$ radio, avec valeur et nature de la tension sous chaque flèche.}
\end{popfigure}

\textbf{3. Nombre de spires $N_1$ :}

De $\dfrac{U_2}{U_1} = \dfrac{N_2}{N_1}$, on tire :
\[
N_1 = N_2 \times \frac{U_1}{U_2} = 45 \times \frac{220}{9} = 45 \times 24{,}44 = 1100 \text{ spires (arrondi à l'unité)}.
\]

\textbf{4. Oscillogrammes attendus :}
\begin{itemize}
\item a) \textbf{À l'entrée :} une courbe \textbf{sinusoïdale} centrée sur \SI{0}{V}, qui prend alternativement des valeurs positives et négatives (tension alternative \SI{220}{V}).
\item b) \textbf{À la sortie :} une courbe située \textbf{toujours au-dessus de \SI{0}{V}} (tension continue \SI{9}{V}), éventuellement légèrement ondulée mais ne changeant jamais de signe.
\end{itemize}

\textbf{5. Deux règles de sécurité (exemples) :}
\begin{itemize}
\item \textbf{Ne jamais toucher les fils dénudés ni les bornes du circuit sous tension} : le secteur \SI{220}{V} peut provoquer une électrocution mortelle.
\item \textbf{Avoir les mains sèches et travailler sur un plan sec} : l'eau conduit le courant et augmenterait le risque d'électrocution.
\item (Autres réponses acceptées : utiliser un voltmètre sur le bon calibre alternatif, faire vérifier le montage par le professeur avant la mise sous tension, ne jamais démonter l'adaptateur branché.)
\end{itemize}

\medskip
\textbf{Barème indicatif (sur 10) :} Q1 : 2 pts ; Q2 : 2 pts (blocs 1 pt, valeurs et natures 1 pt) ; Q3 : 2 pts (expression littérale 1 pt, résultat arrondi 1 pt) ; Q4 : 2 pts ; Q5 : 2 pts (1 pt par règle justifiée).

\textbf{Points de vigilance :} dans la question 3, bien isoler $N_1$ (et non $N_2$) avant l'application numérique ; dans la question 4, le mot-clé attendu pour l'entrée est « sinusoïdale » et pour la sortie « ne change pas de signe » ; chaque règle de sécurité doit être accompagnée de sa justification (le « pourquoi »), pas seulement énoncée.
\end{corrige}

\newpage

\coursec{Fiche bilan}

\begin{popfigure}
\centering
\begin{tikzpicture}[
  mindmap,
  grow cyclic,
  every node/.style={concept, font=\small},
  root concept/.append style={concept color=popblue, font=\bfseries\small, minimum size=2.6cm},
  level 1 concept/.append style={level distance=3.4cm, sibling angle=72, font=\scriptsize},
]
\node[root concept]{L'adaptateur secteur}
  child[concept color=poporange]{ node[align=center]{Pourquoi ? \\ Tension ENEO \\ \SI{220}{V} trop forte} }
  child[concept color=popgreen]{ node[align=center]{Transformation \\ Abaisser la tension \\ (transformateur)} }
  child[concept color=poppurple]{ node[align=center]{Redressement \\ Alternatif $\rightarrow$ continu \\ (diodes)} }
  child[concept color=popgold]{ node[align=center]{Chaîne de conversion \\ Prise $\rightarrow$ transfo \\ $\rightarrow$ pont $\rightarrow$ filtrage} }
  child[concept color=poppink]{ node[align=center]{Sécurité \\ Vérifier $U$ et polarité \\ Pas de surcharge} };
\end{tikzpicture}
\legende{Carte mentale de la leçon : les cinq idées clés sur l'adaptateur secteur.}
\end{popfigure}

\begin{bilan}
\textbf{Définitions clés}
\begin{itemize}
  \item \cle{Adaptateur secteur} : appareil qui adapte la tension du secteur (\SI{220}{V} alternatif chez ENEO) aux besoins d'un appareil (tension continue faible : \SI{5}{V}, \SI{9}{V}, \SI{12}{V}\ldots).
  \item \cle{Transformateur} : dipôle qui modifie la valeur d'une tension \emph{alternative}. Un transformateur \cle{abaisseur} fournit une tension de sortie plus faible que la tension d'entrée.
  \item \cle{Redressement} : conversion d'une tension \cle{alternative} (qui change de sens) en tension \cle{continue} (qui garde le même sens), réalisée par des \cle{diodes}.
  \item \cle{Filtrage} : le condensateur « lisse » la tension redressée pour la rendre presque constante.
\end{itemize}

\textbf{Formules et repères à connaître par c\oe ur}
\begin{center}
\begin{tabularx}{\linewidth}{|l|X|}
\hline
\rowcolor{popblueL} \textbf{Grandeur} & \textbf{Relation / valeur} \\
\hline
Tension du secteur (Cameroun) & $U = \SI{220}{V}$ alternatif, fréquence $f = \SI{50}{Hz}$ \\
\hline
Rapport de transformation & $\dfrac{U_2}{U_1} = \dfrac{N_2}{N_1}$ \quad ($N$ : nombre de spires) \\
\hline
Abaisseur & $N_2 < N_1 \Rightarrow U_2 < U_1$ \\
\hline
Puissance électrique & $P = U \cdot I$ \\
\hline
\end{tabularx}
\end{center}

\textbf{Méthode express : chaîne de conversion d'un adaptateur}
\begin{enumerate}
  \item La prise fournit \SI{220}{V} alternatif.
  \item Le \cle{transformateur} abaisse la tension (ex. \SI{12}{V} alternatif).
  \item Le \cle{pont de diodes} redresse : la tension devant continue mais ondulée.
  \item Le \cle{condensateur} filtre : tension continue quasi constante pour l'appareil.
\end{enumerate}

\textbf{Sécurité}
\begin{itemize}
  \item Toujours vérifier la tension de sortie et la \cle{polarité} ($+$ / $-$) avant de brancher.
  \item Ne jamais ouvrir un adaptateur branché ; ne pas l'utiliser mouillé ou surchauffé.
\end{itemize}
\end{bilan}

\begin{aretenir}[Checklist avant le BEPC]
\begin{itemize}
  \item $\square$ Je sais dire pourquoi on ne branche pas un téléphone directement sur le secteur \SI{220}{V}.
  \item $\square$ Je connais les deux fonctions de l'adaptateur : transformation et redressement.
  \item $\square$ Je sais que le transformateur ne fonctionne qu'en courant alternatif.
  \item $\square$ Je sais appliquer $\dfrac{U_2}{U_1} = \dfrac{N_2}{N_1}$ dans un calcul simple.
  \item $\square$ Je reconnais un transformateur abaisseur ($N_2 < N_1$).
  \item $\square$ Je distingue tension alternative et tension continue (schémas à l'oscilloscope).
  \item $\square$ Je sais que les diodes réalisent le redressement.
  \item $\square$ Je sais que le condensateur filtre (lisse) la tension redressée.
  \item $\square$ Je sais dessiner la chaîne : prise $\rightarrow$ transformateur $\rightarrow$ pont de diodes $\rightarrow$ condensateur $\rightarrow$ appareil.
  \item $\square$ Je sais lire la plaque signalétique d'un adaptateur (entrée / sortie).
  \item $\square$ Je connais les règles de sécurité : polarité, tension adaptée, pas d'ouverture sous tension.
  \item $\square$ Je sais calculer $P = U \cdot I$ pour un adaptateur.
\end{itemize}
\end{aretenir}

\findecours

\end{document}
